% Couple oxydant-réducteur et classification électrochimique — Chimie 1ère D — Cours signé OnBuch+
% Programme officiel MINESEC, Première C, D, E, TI. Compiler avec : tectonic N06-couple-oxydant-reducteur-classification-electrochimique.tex
\newcommand{\DOCMATIERE}{Chimie}
\newcommand{\DOCNIVEAU}{1ère D}
\newcommand{\DOCMODULE}{Module 2 : Oxydoréduction}
\newcommand{\DOCLECON}{06}
\newcommand{\DOCTITRE}{Couple oxydant-réducteur et classification électrochimique}
% OnBuch+ — préambule « cours signés » ENRICHI (compilé avec tectonic / XeLaTeX)
% Système visuel « Pop » d'OnBuch+ : fond crème (jamais blanc), bordures encre
% épaisses, ombres « dures » décalées, titres Archivo Black en capitales.
% Version enrichie pour les cours de chimie : chimie (mhchem, chemfig),
% graphes (pgfplots), boîtes pédagogiques supplémentaires (définition,
% propriété, méthode, attention, le savais-tu, exercice résolu, exercice,
% TP, bilan), page de garde refaite et signature OnBuch+ ancrée en bas.
%
% Macros attendues dans main.tex AVANT \input{preamble} :
%   \DOCMATIERE  \DOCNIVEAU  \DOCMODULE  \DOCLECON (numéro)  \DOCTITRE
\documentclass[11pt]{article}

\usepackage{fontspec}
\usepackage[french]{babel}
\usepackage[a4paper,margin=2cm,top=3.4cm,bottom=2.8cm,headheight=1.4cm,footskip=1.3cm]{geometry}
\usepackage{amsmath,amssymb,amsfonts}
\usepackage[table]{xcolor} % [table] : fournit aussi \rowcolors (alternance de lignes)
\usepackage{pagecolor}
\usepackage{tikz}
\usetikzlibrary{arrows.meta,positioning,calc,shapes.geometric,shapes.misc,shapes.arrows,decorations.pathmorphing,decorations.markings,patterns,shadows,mindmap,trees,fit,backgrounds,matrix,intersections}
\usepackage{pgfplots}
\pgfplotsset{compat=1.18}
\usepackage[version=4]{mhchem}
\usepackage{chemfig}
\usepackage{enumitem}
\usepackage{fancyhdr}
% [most] charge toutes les bibliothèques tcolorbox (listings, minted, raster,
% documentation...) et gonfle inutilement la mémoire TeX sur un document long
% (~300 boîtes) : on ne charge que ce qui est réellement utilisé.
\usepackage[skins,breakable]{tcolorbox}
\usepackage{graphicx}
\usepackage{colortbl}
\usepackage{booktabs}
\usepackage{tabularx}
\usepackage{array}
\usepackage{multicol}
\usepackage{siunitx}
% parse-numbers=false : accepte aussi \SI{2,5 \times 10^{-3}}{...}
\sisetup{locale=FR,output-decimal-marker={,},group-minimum-digits=4,parse-numbers=false}
\DeclareSIUnit\ohm{\ensuremath{\symup{\Omega}}} % U+2126 absent de la police
\usepackage{qrcode}
% \degree : commande fréquemment utilisée par les modèles pour un angle ou une
% température en dehors de \SI{}{} (non fournie par les paquets chargés).
\providecommand{\degree}{^{\circ}}
\usepackage{lastpage}
\usepackage{hyperref}
\hypersetup{hidelinks}

% ============================================================
% Palette « Pop » OnBuch+ — mêmes hex que la classe P (plus_ui.dart)
% ============================================================
\definecolor{poporange}{HTML}{F59321}
\definecolor{popdark}{HTML}{E07A0C}
\definecolor{popcream}{HTML}{FFF6E8}
\definecolor{popink}{HTML}{1C1714}
\definecolor{poppurple}{HTML}{7A5AE0}
\definecolor{popgreen}{HTML}{2E9E6A}
\definecolor{poppink}{HTML}{E0457B}
\definecolor{popblue}{HTML}{2F7DE1}
\definecolor{popgold}{HTML}{C98A12}
\definecolor{popmuted}{HTML}{8A7F76}
\definecolor{popline}{HTML}{EADFD2}
\definecolor{popgray}{HTML}{8A7F76}
\colorlet{popgrey}{popgray}
% Alias tolérants (le modèle nomme parfois les couleurs différemment)
\colorlet{popwhite}{white}
\colorlet{popblack}{popink}
\colorlet{popred}{poppink}
\colorlet{popyellow}{popgold}
% Teintes claires pour fonds de tableaux / zones
\colorlet{popblueL}{popblue!10!white}
\colorlet{poporangeL}{poporange!14!white}
\colorlet{popgreenL}{popgreen!12!white}
\colorlet{poppurpleL}{poppurple!12!white}
\colorlet{poppinkL}{poppink!12!white}
\colorlet{popgoldL}{popgold!14!white}

\pagecolor{popcream}

% ============================================================
% Typographie
% ============================================================
\defaultfontfeatures{Ligatures=TeX}
\setmainfont{PlusJakartaSans-Medium.ttf}[Path=../../../fonts/,BoldFont=PlusJakartaSans-Bold.ttf]
% Maths en sans-serif (Fira Math) avec chiffres et lettres droites de Plus
% Jakarta, pour que les formules se fondent dans le texte.
\usepackage[math-style=ISO,bold-style=ISO]{unicode-math}
\setmathfont{FiraMath-Regular.otf}[Path=../../../fonts/]
\setmathfont{PlusJakartaSans-Medium.ttf}[Path=../../../fonts/,range={up/num,up/{Latin,latin},\mathup}]
\usepackage{icomma} % virgule décimale sans espace en mode math
% Caractères absents de la police de texte (sinon carré vide) : rendus par
% leur équivalent mathématique, en texte comme en formule.
\usepackage{newunicodechar}
\newunicodechar{α}{\ensuremath{\alpha}}
\newunicodechar{Α}{\ensuremath{\alpha}}
\newunicodechar{β}{\ensuremath{\beta}}
\newunicodechar{δ}{\ensuremath{\delta}}
\newunicodechar{✓}{\popcheck{popgreen}}
\newfontfamily\popdisplay{ArchivoBlack-Regular.ttf}[Path=../../../fonts/]
\newfontfamily\poplogo{SpaceGrotesk-Bold.ttf}[Path=../../../fonts/]
\newfontfamily\popsemi{PlusJakartaSans-SemiBold.ttf}[Path=../../../fonts/]
\newfontfamily\popxbold{PlusJakartaSans-ExtraBold.ttf}[Path=../../../fonts/]
\newfontfamily\popxboldls{PlusJakartaSans-ExtraBold.ttf}[Path=../../../fonts/,LetterSpace=3.0]

\setlist[enumerate]{leftmargin=1.7em,itemsep=5pt,topsep=4pt}
\setlist[itemize]{leftmargin=1.7em,itemsep=4pt,topsep=4pt,label={\color{poporange}\textbullet}}
\setlength{\parindent}{0pt}
\setlength{\parskip}{5pt}
\renewcommand{\arraystretch}{1.25}

% chemfig : liaisons un peu plus épaisses, encre Pop
\setchemfig{atom sep=2em,bond style={line width=0.9pt,popink},cram width=3pt}

% pgfplots : style Pop par défaut
\pgfplotsset{
  every axis/.append style={axis line style={popink,line width=1pt},tick style={popink},
    grid style={popline,line width=0.5pt},label style={font=\small},tick label style={font=\footnotesize}},
  popcurve/.style={poporange,line width=1.8pt,no markers,smooth},
}

% ============================================================
% Pill / squiggle
% ============================================================
\newcommand{\poppill}[3]{%
  \tikz[baseline=-0.55ex]\node[draw=popink,line width=1.2pt,rounded corners=2.6pt,%
    fill=#2,inner xsep=6.5pt,inner ysep=3pt]{%
    \popxboldls\fontsize{7}{7}\selectfont\color{#3}\MakeUppercase{#1}};%
}
\newcommand{\popsquiggle}[1]{%
  \tikz[baseline=-0.4ex]{
    \draw[#1,line width=2.6pt,line cap=round]
      plot [smooth] coordinates {(0,0.09) (0.34,0.01) (0.68,0.21) (1.02,0.01) (1.36,0.10)};
  }%
}
% Mot-clé surligné (marqueur orange sous le texte)
\newcommand{\cle}[1]{{\popxbold\color{popdark}#1}}

% ============================================================
% En-tête / pied
% ============================================================
\pagestyle{fancy}
\fancyhf{}
\renewcommand{\headrulewidth}{0pt}
\renewcommand{\footrulewidth}{0pt}
\lhead{\raisebox{-0.15ex}{\poplogo\fontsize{13}{13}\selectfont\color{popink}OnBuch\color{poporange}+}}
\rhead{\poppill{\DOCMATIERE\ \textbullet\ \DOCNIVEAU\ \textbullet\ Leçon \DOCLECON}{white}{popink}}
\lfoot{\popxbold\fontsize{7.6}{7.6}\selectfont\color{popmuted}\MakeUppercase{Cours signé OnBuch+}\ \textbullet\ {\popsemi\DOCTITRE}}
\rfoot{\tikz[baseline=-0.6ex]\node[rounded corners=8.5pt,fill=popink,minimum height=17pt,minimum width=17pt,inner xsep=5pt,inner sep=0]{\popxbold\fontsize{8}{8}\selectfont\color{popcream}\thepage\,/\,\pageref*{LastPage}};}

% ============================================================
% Titres
% ============================================================
\newcommand{\chaptitre}[2]{%
  \begin{tcolorbox}[enhanced,colback=poporange,colframe=popink,boxrule=2.6pt,arc=11pt,
      drop shadow=popink,left=15pt,right=15pt,top=12pt,bottom=11pt,before skip=3pt,after skip=18pt]
    \poppill{#2}{popink}{popcream}\par\vspace{9pt}
    {\raggedright\hyphenpenalty=10000\exhyphenpenalty=10000\popdisplay\fontsize{22}{24}\selectfont\color{popcream}\MakeUppercase{#1}\par}
  \end{tcolorbox}
}
% Section numérotée : pastille orange avec le numéro + titre Archivo
\newcounter{popsec}
\newcommand{\coursec}[1]{%
  \stepcounter{popsec}\setcounter{popsub}{0}%
  \par\vspace{16pt}\noindent
  \begin{minipage}{\linewidth}
  \tikz[baseline=-0.7ex]\node[draw=popink,line width=1.6pt,fill=poporange,rounded corners=4pt,minimum size=22pt,inner sep=0]{\popdisplay\fontsize{12}{12}\selectfont\color{popcream}\thepopsec};%
  \hspace{8pt}{\popdisplay\fontsize{15}{17}\selectfont\color{popink}\MakeUppercase{#1}}\par
  \vspace{4pt}\noindent{\color{poporange}\rule{36pt}{3.6pt}}
  \end{minipage}\par\nopagebreak\vspace{8pt}
}
\newcounter{popsub}
\newcommand{\courssub}[1]{%
  \stepcounter{popsub}%
  \par\vspace{9pt}\noindent{\popxbold\fontsize{11.5}{13}\selectfont\color{popink}{\color{poporange}\thepopsec.\thepopsub}\hspace{6pt}#1}\par\nopagebreak\vspace{4pt}
}

% ============================================================
% Boîtes pédagogiques — même recette : fond blanc, bordure épaisse colorée,
% ombre dure encre, badge plein. Titre optionnel [#1].
% ============================================================
\tcbset{popbase/.style={breakable,enhanced,colback=white,boxrule=2.3pt,arc=9pt,
  shadow={3pt}{-3pt}{0pt}{popink},left=14pt,right=14pt,top=11pt,bottom=11pt,before skip=11pt,after skip=15pt,
  fontupper=\popsemi\fontsize{10.6}{14.5}\selectfont\color{popink},
  fontlower=\popsemi\fontsize{10.6}{14.5}\selectfont\color{popink}}}
% Coche dessinée (le glyphe ✓ n'existe pas dans les polices du document)
\newcommand{\popcheck}[1]{\tikz[baseline=-0.5ex]\draw[#1,line width=1.6pt,line cap=round,line join=round](-0.13,0.0)--(-0.04,-0.09)--(0.13,0.1);}
\newcommand{\popboxhead}[3]{\poppill{#1}{#2}{white}\ifx\relax#3\relax\else\hspace{6pt}{\popxbold\fontsize{10.6}{12}\selectfont\color{popink}#3}\fi\par\vspace{7pt}}

\newtcolorbox{aretenir}[1][]{popbase,colframe=popgreen,before upper={\popboxhead{À retenir}{popgreen}{#1}}}
\newtcolorbox{exemplebox}[1][]{popbase,colframe=poporange,before upper={\popboxhead{Exemple}{poporange}{#1}}}
\newtcolorbox{definition}[1][]{popbase,colframe=popblue,before upper={\popboxhead{Définition}{popblue}{#1}}}
\newtcolorbox{propriete}[1][]{popbase,colframe=poppurple,before upper={\popboxhead{Propriété}{poppurple}{#1}}}
\newtcolorbox{methode}[1][]{popbase,colframe=popgold,before upper={\popboxhead{Méthode}{popgold}{#1}}}
\newtcolorbox{attention}[1][]{popbase,colframe=poppink,before upper={\popboxhead{Attention}{poppink}{#1}}}
\newtcolorbox{experience}[1][]{popbase,colframe=popblue,colback=popblueL,before upper={\popboxhead{Expérience}{popblue}{#1}}}
% « Le savais-tu ? » : carte violette pleine (contexte, culture, Cameroun)
\newtcolorbox{savaistu}[1][]{popbase,colback=poppurple,colframe=popink,
  fontupper=\popsemi\fontsize{10.4}{14.2}\selectfont\color{white},
  before upper={\poppill{Le savais-tu\,?}{white}{poppurple}\ifx\relax#1\relax\else\hspace{6pt}{\popxbold\color{white}#1}\fi\par\vspace{7pt}}}
% Exercice résolu : énoncé en haut, \tcblower puis la solution sur fond crème
\newcounter{popexo}\newcounter{popexr}
\newtcolorbox{exoresolu}[1][]{popbase,colframe=poporange,colbacklower=popcream,
  segmentation style={popink,line width=1.2pt,dashed},
  before upper={\stepcounter{popexr}\popboxhead{Exercice résolu \thepopexr}{poporange}{#1}},
  before lower={\poppill{Solution}{popink}{popcream}\par\vspace{6pt}}}
% Exercice (non résolu en place) — la correction est dans la partie Corrigés
\newtcolorbox{exercice}[1][]{popbase,colframe=popink,
  before upper={\stepcounter{popexo}\popboxhead{Exercice \thepopexo}{popink}{#1}}}
% Corrigé
\newtcolorbox{corrige}[1][]{popbase,colframe=popgreen,colback=popgreenL,
  before upper={\popboxhead{Corrigé}{popgreen}{#1}}}
% Objectifs / prérequis / bilan
\newtcolorbox{objectifs}{popbase,colframe=popink,colback=poporangeL,
  before upper={\popboxhead{Objectifs de la leçon}{popink}{}}}
\newtcolorbox{prerequis}{popbase,colframe=popink,colback=popblueL,
  before upper={\popboxhead{Prérequis}{popblue}{}}}
\newtcolorbox{bilan}{popbase,colframe=popink,colback=white,boxrule=2.8pt,
  before upper={\popboxhead{Fiche bilan}{poporange}{L'essentiel à connaître pour l'examen}}}
% Figure encadrée avec légende
\newtcolorbox{popfigure}[1][]{enhanced,breakable=false,colback=white,colframe=popline,boxrule=1.4pt,arc=8pt,
  left=10pt,right=10pt,top=10pt,bottom=8pt,before skip=10pt,after skip=12pt,halign=center,
  lower separated=false,halign lower=center,
  fontlower=\popsemi\fontsize{9}{11.5}\selectfont\color{popmuted},
  #1}
\newcommand{\legende}[1]{\par\vspace{6pt}{\popsemi\fontsize{9}{11.5}\selectfont\color{popmuted}#1}}

% ============================================================
% Signature OnBuch+
% ============================================================
\newcommand{\poplogoinline}[1]{{\poplogo\fontsize{#1}{#1}\selectfont\color{popink}OnBuch\color{poporange}+}}
\newcommand{\signatureonbuch}{%
  \begin{tcolorbox}[enhanced,colback=popink,colframe=popink,boxrule=2.2pt,arc=10pt,
      drop shadow=poporange,left=14pt,right=14pt,top=10pt,bottom=10pt,before skip=0pt,after skip=0pt]
    \tikz[baseline=-0.7ex]\node[circle,fill=popgreen,draw=popcream,line width=1.3pt,minimum size=22pt,inner sep=0]{\popcheck{white}};%
    \hspace{9pt}\begin{minipage}[c]{0.62\linewidth}
      {\popxbold\fontsize{12}{13}\selectfont\color{popcream}Signé\ \textbullet\ {\poplogo OnBuch\color{poporange}+}}\par\vspace{2pt}
      {\raggedright\popsemi\fontsize{8}{10.5}\selectfont\color{popline}\MakeUppercase{Équipe pédagogique OnBuch+}\\\MakeUppercase{Conforme au programme officiel MINESEC}\par}
    \end{minipage}\hfill
    {\poplogo\fontsize{22}{22}\selectfont\color{popcream}OnBuch\color{poporange}+}
  \end{tcolorbox}%
}

% ============================================================
% Page de garde
% #1 titre · #2 matière · #3 niveau · #4 module · #5 n° leçon · #6 durée ·
% #7 description / objectifs (texte ou itemize)
% ============================================================
\newcommand{\pagegarde}[7]{%
  \thispagestyle{empty}
  \newgeometry{a4paper,margin=1.7cm,top=1.6cm,bottom=1.4cm}%
  \begin{tikzpicture}[remember picture,overlay]
    \fill[poporange!18] ([xshift=-3.2cm,yshift=-3.6cm]current page.north east) circle (4.6cm);
    \fill[poppurple!14] ([xshift=2.4cm,yshift=7.6cm]current page.south west) circle (3.8cm);
  \end{tikzpicture}%
  {\poplogo\fontsize{30}{30}\selectfont\color{popink}OnBuch\color{poporange}+}\hfill
  \raisebox{0.6ex}{\poppill{Cours signé \textbullet\ Premium}{popink}{popcream}}\par
  \vspace{1pt}\popsquiggle{poppurple}\par
  \vspace{8pt}
  \begin{tcolorbox}[enhanced,colback=poporange,colframe=popink,boxrule=2.8pt,arc=13pt,
      drop shadow=popink,left=18pt,right=18pt,top=12pt,bottom=13pt,after skip=10pt]
    \poppill{#2 \textbullet\ #3}{popink}{popcream}\hspace{5pt}\poppill{#4}{white}{popink}\par\vspace{8pt}
    {\popdisplay\fontsize{14}{14}\selectfont\color{popink}LEÇON #5}\par\vspace{4pt}
    {\raggedright\hyphenpenalty=10000\exhyphenpenalty=10000\popdisplay\fontsize{25}{27}\selectfont\color{popcream}\MakeUppercase{#1}\par}
    \vspace{8pt}
    {\popxbold\fontsize{9.5}{9.5}\selectfont\color{popink}\MakeUppercase{Durée conseillée : #6}}
  \end{tcolorbox}
  \begin{tcolorbox}[enhanced,colback=white,colframe=popink,boxrule=2.2pt,arc=10pt,
      drop shadow=popink,left=16pt,right=16pt,top=10pt,bottom=10pt,after skip=8pt]
    \poppill{Dans cette leçon}{poppurple}{white}\par\vspace{5pt}
    \popsemi\fontsize{9.3}{12.6}\selectfont\color{popink}#7
  \end{tcolorbox}
  \noindent\begin{minipage}[t]{0.487\linewidth}
    \begin{tcolorbox}[enhanced,colback=popgreen,colframe=popink,boxrule=2.2pt,arc=10pt,
        drop shadow=popink,left=11pt,right=11pt,top=10pt,bottom=10pt,height=3.1cm,valign=center]
      \tcbox[colback=white,colframe=popink,boxrule=1.3pt,arc=5pt,boxsep=0pt,nobeforeafter,
        left=3pt,right=3pt,top=3pt,bottom=3pt,valign=center]{\qrcode[height=1.9cm]{https://play.google.com/store/apps/details?id=com.luvvix.onbuch&hl=fr}}%
      \hspace{8pt}\begin{minipage}[c]{0.5\linewidth}
        {\popdisplay\fontsize{11}{12.5}\selectfont\color{white}\MakeUppercase{Télécharge OnBuch}}\par\vspace{4pt}
        {\raggedright\popsemi\fontsize{8.4}{11}\selectfont\color{white}Gratuit \textbullet\ Google Play\\Tuteur IA, annales, concours, résultats\par}
      \end{minipage}
    \end{tcolorbox}
  \end{minipage}\hfill
  \begin{minipage}[t]{0.487\linewidth}
    \begin{tcolorbox}[enhanced,colback=poppurple,colframe=popink,boxrule=2.2pt,arc=10pt,
        drop shadow=popink,left=11pt,right=11pt,top=10pt,bottom=10pt,height=3.1cm,valign=center]
      \tikz[baseline=-0.7ex]\node[circle,fill=white,draw=popink,line width=1.3pt,minimum size=1.6cm,inner sep=0]{\popdisplay\fontsize{24}{24}\selectfont\color{poppurple}+};%
      \hspace{8pt}\begin{minipage}[c]{0.6\linewidth}
        {\popdisplay\fontsize{11}{12.5}\selectfont\color{white}\MakeUppercase{Passe à OnBuch+}}\par\vspace{4pt}
        {\raggedright\popsemi\fontsize{8.4}{11}\selectfont\color{white}Tous les cours signés, Tuteur IA illimité, fiches PDF — onglet OnBuch+ de l'app\par}
      \end{minipage}
    \end{tcolorbox}
  \end{minipage}
  \vspace{8pt}
  \popcontenu
  \vfill
  \signatureonbuch
  \restoregeometry
  \newpage
}

% Rangée « Ce cours contient » (page de garde)
\newcommand{\poptile}[3]{% #1 couleur #2 gros texte #3 légende
  \begin{tcolorbox}[enhanced,colback=#1,colframe=popink,boxrule=2pt,arc=9pt,shadow={2.5pt}{-2.5pt}{0pt}{popink},
      left=6pt,right=6pt,top=9pt,bottom=9pt,halign=center,valign=center,height=1.75cm,width=\linewidth,nobeforeafter]
    {\popdisplay\fontsize{12.5}{13}\selectfont\color{white}\MakeUppercase{#2}}\par\vspace{3pt}
    {\centering\popsemi\fontsize{7.8}{9.5}\selectfont\color{white}#3\par}
  \end{tcolorbox}}
\newcommand{\popcontenu}{%
  \par\noindent{\popxboldls\fontsize{7.5}{7.5}\selectfont\color{popmuted}CE COURS SIGNÉ CONTIENT}\par\vspace{7pt}
  \noindent\begin{minipage}[t]{0.235\linewidth}\poptile{popblue}{Cours}{complet, illustré, pas à pas}\end{minipage}\hfill
  \begin{minipage}[t]{0.235\linewidth}\poptile{poporange}{Méthodes}{et exercices résolus}\end{minipage}\hfill
  \begin{minipage}[t]{0.235\linewidth}\poptile{poppink}{Exercices}{type examen + corrigés}\end{minipage}\hfill
  \begin{minipage}[t]{0.235\linewidth}\poptile{popgreen}{Fiche bilan}{l'essentiel à retenir}\end{minipage}\par}

% Sommaire
\newtcolorbox{sommaire}{enhanced,colback=white,colframe=popink,boxrule=2.2pt,arc=10pt,
  drop shadow=popink,left=18pt,right=18pt,top=13pt,bottom=13pt,before skip=6pt,after skip=20pt,
  before upper={\poppill{Sommaire}{poppurple}{white}\par\vspace{9pt}\popsemi\fontsize{10.6}{14.5}\selectfont\color{popink}}}

% Signature de fin de cours — poussée tout en bas de la dernière page
\newcommand{\findecours}{%
  \par\vfill\nopagebreak
  \begin{center}\popsquiggle{poporange}\end{center}\vspace{4pt}
  \signatureonbuch
}
% Glyphes absents de Fira Math : redéfinis à partir de glyphes disponibles
\AtBeginDocument{%
  \def\setminus{\mathbin{\backslash}}\let\smallsetminus\setminus
  \def\vdots{\mathord{\vbox{\baselineskip3.2pt\lineskiplimit0pt\kern4pt\hbox{.}\hbox{.}\hbox{.}}}}%
  \def\ddots{\mathinner{\mkern1mu\raise6.5pt\hbox{.}\mkern2mu\raise3.6pt\hbox{.}\mkern2mu\raise0.7pt\hbox{.}\mkern1mu}}%
  \def\triangle{\mathord{\scalebox{0.9}{$\symup{\Delta}$}}}%
  \def\nabla{\mathord{\raisebox{\depth}{\scalebox{1}[-1]{$\symup{\Delta}$}}}}%
  \def\checkmark{\popcheck{popdark}}%
  \def\star{\mathbin{\ast}}\def\bigstar{\mathord{\ast}}%
  \def\diamond{\mathbin{\scalebox{0.6}{$\Diamond$}}}%
  \def\bigcup{\mathop{\mathchoice{\raisebox{-0.3em}{\scalebox{1.7}{$\cup$}}}{\scalebox{1.25}{$\cup$}}{\cup}{\cup}}\displaylimits}%
  \def\bigcap{\mathop{\mathchoice{\raisebox{-0.3em}{\scalebox{1.7}{$\cap$}}}{\scalebox{1.25}{$\cap$}}{\cap}{\cap}}\displaylimits}%
  \def\lesseqqgtr{\mathrel{\lessgtr}}\def\bigoplus{\mathop{\oplus}\displaylimits}%
}
\providecommand{\ssi}{si et seulement si} % abréviation fréquente

%% FIGFIX-BEGIN v4 — correctifs globaux des figures (docs/onbuch-generation/tools/figfix)
\usepackage{adjustbox}\usepackage{etoolbox}
% toute figure TikZ/pgfplots plus large que la ligne est réduite à la largeur disponible
\BeforeBeginEnvironment{tikzpicture}{\begin{adjustbox}{max width=\linewidth}}
\AfterEndEnvironment{tikzpicture}{\end{adjustbox}}
% légendes pgfplots lisibles par-dessus les courbes
\pgfplotsset{every axis/.append style={legend style={fill=white,fill opacity=0.92,text opacity=1,draw=black!15,inner sep=3pt}}}
% étiquettes d'arêtes (syntaxe edge["..."]) avec fond blanc
\tikzset{every edge quotes/.append style={fill=white,inner sep=1.2pt}}
%% FIGFIX-END
\begin{document}

\pagegarde{Couple oxydant-réducteur et classification électrochimique}{Chimie}{1ère D}{Module 2 : Oxydoréduction}{06}{4 h}{Cette leçon introduit la notion de couple oxydant-réducteur et montre comment comparer expérimentalement les couples entre eux. À partir de réactions entre métaux et ions métalliques ou ions H3O+, on établit une classification électrochimique qualitative permettant de prévoir le sens des réactions d'oxydoréduction.
\vspace{4pt}
\begin{multicols}{2}\small
\begin{itemize}
\item À la fin de cette leçon, je sais définir un couple oxydant-réducteur et écrire sa demi-équation électronique
\item identifier l'oxydant et le réducteur dans un couple Ox/Red et dans une réaction donnée
\item réaliser et interpréter des expériences mettant en jeu des métaux et des ions métalliques (couples Mn+/M)
\item réaliser et interpréter des expériences mettant en jeu le couple H3O+/H2 avec différents métaux
\item déduire d'une réaction spontanée la comparaison des pouvoirs oxydant et réducteur de deux couples
\item établir une classification électrochimique qualitative des couples étudiés
\end{itemize}\end{multicols}}

\begin{sommaire}
\begin{multicols}{2}
\begin{enumerate}
\item Couple oxydant-réducteur : définition et écriture
\item Réaction entre un métal et un ion métallique : couples Mn+/M
\item Action des ions H3O+ sur les métaux : le couple H3O+/H2
\item Construction de la classification électrochimique qualitative
\item Utilisation de la classification : prévision des réactions
\item Synthèse et applications au quotidien camerounais
\item Travaux pratiques
\item Méthodes et exercices résolus
\item Exercices
\item Corrigés des exercices
\item Fiche bilan
\end{enumerate}
\end{multicols}
\end{sommaire}

\begin{objectifs}
\begin{itemize}
\item À la fin de cette leçon, je sais définir un couple oxydant-réducteur et écrire sa demi-équation électronique
\item identifier l'oxydant et le réducteur dans un couple Ox/Red et dans une réaction donnée
\item réaliser et interpréter des expériences mettant en jeu des métaux et des ions métalliques (couples Mn+/M)
\item réaliser et interpréter des expériences mettant en jeu le couple H3O+/H2 avec différents métaux
\item déduire d'une réaction spontanée la comparaison des pouvoirs oxydant et réducteur de deux couples
\item établir une classification électrochimique qualitative des couples étudiés
\item placer correctement le couple H3O+/H2 dans cette classification
\item prévoir si une réaction d'oxydoréduction entre deux couples donnés est possible ou non
\item expliquer des phénomènes quotidiens (corrosion, protection des métaux) à l'aide de la classification
\end{itemize}
\end{objectifs}

\begin{prerequis}
\begin{itemize}
\item Notion d'oxydation et de réduction (transfert d'électrons) vue dans la leçon précédente
\item Écriture et équilibrage des demi-équations électroniques
\item Équilibrage d'une équation-bilan d'oxydoréduction par combinaison des demi-équations
\item Ions courants en solution aqueuse : Cu2+, Zn2+, Fe2+, Ag+, H3O+
\item Tests d'identification simples des ions et du dihydrogène (détonation à la flamme)
\end{itemize}
\end{prerequis}

\newpage

\coursec{Couple oxydant-réducteur : définition et écriture}

Au marché Mokolo à Yaoundé, un vendeur de clous observe chaque semaine le même phénomène : ses clous en fer, exposés à l'air humide de la saison des pluies, se couvrent rapidement d'une poudre brun-rougeâtre — la rouille. À quelques mètres, sa voisine vend des ustensiles en cuivre qui, eux, conservent longtemps leur bel éclat rosé. Intrigué, notre vendeur plonge un jour un clou en fer dans une solution bleue de sulfate de cuivre : en quelques minutes, le clou se recouvre d'un dépôt rougeâtre de cuivre métallique, et la solution bleue pâlit. Mais l'expérience inverse — un fil de cuivre plongé dans une solution de sulfate de fer(II) — ne produit strictement rien.

Pourquoi certains métaux s'oxydent-ils facilement et d'autres non ? Pourquoi le fer arrache-t-il le cuivre à ses ions, alors que le cuivre est incapable de faire de même avec le fer ? Peut-on \emph{prévoir à l'avance} quelles réactions d'oxydoréduction sont possibles ?

Pour répondre à ces questions, il nous faut d'abord un outil conceptuel fondamental : la notion de \cle{couple oxydant-réducteur}. C'est l'objet de cette première section. Une fois cet outil maîtrisé, nous construirons expérimentalement, dans les sections suivantes, la \cle{classification électrochimique} qui permet de prévoir toutes ces réactions.

\courssub{Rappels : oxydation, réduction, oxydant, réducteur}

Avant de définir les couples, remettons en place le vocabulaire de base de l'oxydoréduction, déjà rencontré dans les leçons précédentes. Toute l'oxydoréduction repose sur un seul phénomène microscopique : le \cle{transfert d'électrons} d'une espèce chimique vers une autre.

\begin{definition}[Oxydation, réduction, oxydant, réducteur]
\begin{itemize}
\item Une \cle{oxydation} est une \textbf{perte d'électrons} par une espèce chimique.
\item Une \cle{réduction} est un \textbf{gain d'électrons} par une espèce chimique.
\item Un \cle{oxydant} est une espèce chimique capable de \textbf{capter} un ou plusieurs électrons (il est donc réduit).
\item Un \cle{réducteur} est une espèce chimique capable de \textbf{céder} un ou plusieurs électrons (il est donc oxydé).
\end{itemize}
\end{definition}

Un moyen mnémotechnique très utilisé : \emph{« le réducteur se réduit à céder ses électrons, et c'est lui qui est oxydé »}. Plus simplement : l'oxydant \textbf{prend}, le réducteur \textbf{donne}.

\begin{attention}[Les deux pièges du vocabulaire]
Deux erreurs reviennent chaque année au Probatoire :
\begin{itemize}
\item Croire qu'un oxydant s'oxyde. C'est l'inverse : l'oxydant \textbf{captant} des électrons est \textbf{réduit}. L'oxydant subit une réduction ; le réducteur subit une oxydation.
\item Croire qu'une espèce peut être oxydée seule. Une oxydation ne se produit \textbf{jamais seule} : si une espèce perd des électrons, une autre doit obligatoirement les gagner. Oxydation et réduction sont toujours simultanées, comme les deux faces d'une même pièce.
\end{itemize}
\end{attention}

\begin{popfigure}
\begin{center}
\begin{tikzpicture}[scale=1.05, >=Stealth]
% Réducteur (gauche)
\shade[ball color=poporange!80] (0,0) circle (1.15);
\node[white, font=\bfseries] at (0,0.25) {Réducteur};
\node[white, font=\small] at (0,-0.35) {Red};
\node[popdark, font=\small\itshape, align=center] at (0,-1.55) {cède des électrons\\ (il est \textbf{oxydé})};
% Oxydant (droite)
\shade[ball color=popblue!85] (7,0) circle (1.15);
\node[white, font=\bfseries] at (7,0.25) {Oxydant};
\node[white, font=\small] at (7,-0.35) {Ox};
\node[popdark, font=\small\itshape, align=center] at (7,-1.55) {capte des électrons\\ (il est \textbf{réduit})};
% Flèche des électrons
\draw[->, line width=2.2pt, popgreen!70!black, decorate, decoration={snake, amplitude=0.6mm, segment length=5mm}] (1.35,0.55) -- (5.65,0.55);
\node[popgreen!50!black, font=\bfseries] at (3.5,1.05) {$n\,e^-$};
\node[popmuted, font=\small\itshape] at (3.5,-0.35) {transfert d'électrons};
\end{tikzpicture}
\end{center}
\legende{Le cœur de toute réaction d'oxydoréduction : un transfert d'électrons du réducteur (qui donne) vers l'oxydant (qui prend).}
\end{popfigure}

\courssub{Notion de couple oxydant-réducteur}

Revenons au clou du marché Mokolo. Lorsque le fer Fe s'oxyde en ions \ce{Fe^2+}, il perd deux électrons. Mais l'opération inverse est-elle possible ? Oui : dans une électrolyse, par exemple, des ions \ce{Fe^2+} peuvent capter deux électrons et redevenir du fer métallique. L'ion \ce{Fe^2+} et le métal Fe sont donc deux formes d'un même élément, séparées par un simple échange d'électrons : on dit qu'ils sont \emph{conjugués}. C'est exactement l'idée de couple.

\begin{definition}[Couple oxydant-réducteur]
Un \cle{couple oxydant-réducteur} (ou couple rédox) est l'ensemble formé d'un \textbf{oxydant Ox} et de son \textbf{réducteur conjugué Red}, qui se transforment l'un en l'autre par transfert d'électrons.

On le note \textbf{Ox/Red} : \emph{l'oxydant s'écrit toujours en premier}, séparé du réducteur par une barre oblique.

À chaque couple est associée une \cle{demi-équation électronique} :
\[ \text{Ox} + n\,e^- \ce{<=>} \text{Red} \]
où $n$ est le nombre d'électrons échangés ($n \in \mathbb{N}^*$).
\end{definition}

La double flèche \ce{<=>} est essentielle : elle traduit le fait que, selon les conditions expérimentales, la transformation peut se faire dans un sens (réduction de Ox en Red) ou dans l'autre (oxydation de Red en Ox). Dire qu'une espèce est « un oxydant » n'a donc de sens que \emph{par rapport à une autre espèce} : le même ion \ce{Cu^2+} est un oxydant face au fer, mais le cuivre métallique Cu est un réducteur face aux ions \ce{Ag+}.

\begin{attention}[Convention d'écriture : l'oxydant d'abord !]
L'erreur la plus fréquente consiste à écrire le réducteur en premier dans la notation du couple. Ainsi, écrire \textbf{Cu/\ce{Cu^2+}} est \textbf{incorrect} : la seule écriture acceptée est \ce{Cu^2+}/Cu. Astuce : on lit le couple « \ce{Cu^2+} sur Cu », et l'oxydant est « au-dessus », comme s'il dominait. De même, on écrit \ce{Fe^3+}/\ce{Fe^2+} et jamais \ce{Fe^2+}/\ce{Fe^3+}.
\end{attention}

\courssub{Demi-équations électroniques : écriture et équilibrage}

Écrire la demi-équation d'un couple, c'est traduire le passage de l'oxydant au réducteur par un gain de $n$ électrons. Comme toute équation en chimie, elle doit respecter deux lois de conservation : la \textbf{conservation des éléments} et la \textbf{conservation de la charge électrique}.

\begin{methode}[Écrire et équilibrer une demi-équation électronique]
\begin{enumerate}
\item \textbf{Identifier l'oxydant et le réducteur} du couple, puis écrire l'oxydant à gauche et le réducteur à droite : $\text{Ox} \ce{<=>} \text{Red}$.
\item \textbf{Équilibrer l'élément caractéristique} (métal, hydrogène, etc.) en ajustant les coefficients stœchiométriques.
\item \textbf{Équilibrer les charges} en ajoutant $n$ électrons du côté de l'oxydant (à gauche), de sorte que la charge totale soit la même des deux côtés.
\item \textbf{Vérifier} : même nombre d'atomes de chaque élément et même charge totale des deux côtés.
\end{enumerate}
\end{methode}

\begin{exemplebox}[Demi-équations de quelques couples usuels]
\begin{itemize}
\item Couple \ce{Cu^2+}/Cu : l'ion cuivre(II) capte deux électrons pour donner le métal cuivre :
\[ \ce{Cu^2+ + 2e- <=> Cu} \]
Vérification des charges : à gauche $+2 + (-2) = 0$ ; à droite $0$. L'équation est équilibrée.
\item Couple \ce{Zn^2+}/Zn :
\[ \ce{Zn^2+ + 2e- <=> Zn} \]
\item Couple \ce{Fe^2+}/Fe :
\[ \ce{Fe^2+ + 2e- <=> Fe} \]
\item Couple \ce{Ag+}/Ag : un seul électron échangé :
\[ \ce{Ag+ + e- <=> Ag} \]
\item Couple \ce{H3O+}/H2 : l'ion oxonium, responsable de l'acidité, peut être réduit en dihydrogène gazeux. On équilibre l'oxygène par l'eau, l'hydrogène par \ce{H3O+}, puis la charge par les électrons :
\[ \ce{2H3O+ + 2e- <=> H2 + 2H2O} \]
Vérification : atomes H (6=2+4), O (2=2), charges ($+2-2=0$ à gauche, $0$ à droite).
\end{itemize}
\end{exemplebox}

Le couple \ce{H3O+}/H2 occupe une place particulière : il servira de \textbf{référence} dans la classification électrochimique que nous construirons dans les sections suivantes. Retiens bien sa demi-équation.

\begin{exoresolu}[Écrire les demi-équations des couples \ce{Ag+}/Ag et \ce{Fe^3+}/\ce{Fe^2+}]
\textbf{Énoncé.} Pour chacun des couples oxydant-réducteur suivants, identifier l'oxydant et le réducteur, puis écrire la demi-équation électronique équilibrée : a) \ce{Ag+}/Ag ; b) \ce{Fe^3+}/\ce{Fe^2+}.
\tcblower
\textbf{Solution détaillée.}

\textbf{a) Couple \ce{Ag+}/Ag.}

Étape 1 : dans la notation du couple, l'oxydant est écrit en premier : l'oxydant est \ce{Ag+}, le réducteur conjugué est Ag.

Étape 2 : l'élément argent est déjà équilibré (un atome de chaque côté) : $\ce{Ag+ <=> Ag}$.

Étape 3 : équilibrons les charges. À gauche, la charge est $+1$ ; à droite, elle est nulle. Il faut donc ajouter \textbf{un électron} à gauche :
\[ \ce{Ag+ + e- <=> Ag} \]
Vérification : charges $+1 - 1 = 0$ à gauche, $0$ à droite. L'argent métallique déposé sur le clou de certaines expériences de chimie vient précisément de cette réduction.

\textbf{b) Couple \ce{Fe^3+}/\ce{Fe^2+}.}

Étape 1 : l'oxydant est \ce{Fe^3+} (écrit en premier), le réducteur conjugué est \ce{Fe^2+}. Remarque importante : ici, \textbf{les deux formes sont des ions en solution} ; il n'y a pas de métal.

Étape 2 : l'élément fer est équilibré : $\ce{Fe^3+ <=> Fe^2+}$.

Étape 3 : charges : $+3$ à gauche, $+2$ à droite. Il faut ajouter \textbf{un électron} à gauche :
\[ \ce{Fe^3+ + e- <=> Fe^2+} \]
Vérification : $+3 - 1 = +2$ à gauche, $+2$ à droite. C'est cette réduction qui explique le virage au vert pâle d'une solution jaune-orangée de chlorure de fer(III) lorsqu'on y plonge un réducteur.
\end{exoresolu}

\courssub{Couples métalliques et couples en solution : deux grandes familles}

Les exemples précédents révèlent une distinction importante qu'il faut savoir faire.

\begin{definition}[Couples métalliques et couples en solution]
\begin{itemize}
\item Un \cle{couple métallique} est de la forme \ce{M^n+}/M : l'oxydant est un \textbf{cation métallique en solution} et le réducteur est le \textbf{métal solide} correspondant. Exemples : \ce{Cu^2+}/Cu, \ce{Zn^2+}/Zn, \ce{Fe^2+}/Fe, \ce{Ag+}/Ag, \ce{Al^3+}/Al.
\item Dans d'autres couples, \textbf{les deux formes sont des espèces dissoutes} en solution : l'oxydant et le réducteur ne diffèrent que par leur degré d'oxydation. Exemple : \ce{Fe^3+}/\ce{Fe^2+}, où l'ion fer(III) jaune-orangé est l'oxydant et l'ion fer(II) vert pâle le réducteur.
\end{itemize}
\end{definition}

Cette distinction aura des conséquences pratiques : dans un couple métallique, la réduction se manifeste par un \textbf{dépôt solide} (le métal qui se forme, comme le dépôt rougeâtre de cuivre sur le clou du marché Mokolo) ou par la \textbf{disparition d'un solide} (le métal qui s'oxyde et « fond » dans la solution). Dans un couple en solution, seule la \textbf{couleur de la solution} change, ce qui est souvent plus délicat à observer.

\begin{center}
\begin{tabularx}{\linewidth}{|l|l|l|X|}
\hline
\rowcolor{popblueL} \textbf{Couple Ox/Red} & \textbf{Oxydant} & \textbf{Réducteur} & \textbf{Demi-équation électronique} \\
\hline
\ce{Cu^2+}/\ce{Cu} & \ce{Cu^2+} (bleu) & \ce{Cu} (rougeâtre) & \ce{Cu^2+ + 2e- <=> Cu} \\
\hline
\ce{Zn^2+}/\ce{Zn} & \ce{Zn^2+} (incolore) & \ce{Zn} (gris) & \ce{Zn^2+ + 2e- <=> Zn} \\
\hline
\ce{Fe^2+}/\ce{Fe} & \ce{Fe^2+} (vert pâle) & \ce{Fe} (gris) & \ce{Fe^2+ + 2e- <=> Fe} \\
\hline
\ce{Ag+}/\ce{Ag} & \ce{Ag+} (incolore) & \ce{Ag} (argenté) & \ce{Ag+ + e- <=> Ag} \\
\hline
\ce{Fe^3+}/\ce{Fe^2+} & \ce{Fe^3+} (orangé) & \ce{Fe^2+} (vert pâle) & \ce{Fe^3+ + e- <=> Fe^2+} \\
\hline
\ce{H3O+}/\ce{H2} & \ce{H3O+} (incolore) & \ce{H2} (gaz incolore) & \ce{2H3O+ + 2e- <=> H2 + 2H2O} \\
\hline
\end{tabularx}
\end{center}

\courssub{Des demi-équations à l'équation-bilan : une réaction, deux couples}

Voici maintenant l'idée qui donne toute sa puissance à la notion de couple. Reprenons l'expérience du marché : le clou en fer plongé dans la solution de sulfate de cuivre. Que se passe-t-il au niveau microscopique ?

\begin{itemize}
\item Le fer Fe, réducteur du couple \ce{Fe^2+}/Fe, \textbf{cède} deux électrons : $\ce{Fe -> Fe^2+ + 2e-}$ (oxydation, demi-équation lue de droite à gauche).
\item Les ions \ce{Cu^2+}, oxydant du couple \ce{Cu^2+}/Cu, \textbf{captent} ces deux électrons : $\ce{Cu^2+ + 2e- -> Cu}$ (réduction).
\end{itemize}

Les électrons ne circulent pas librement dans la solution : ils passent \emph{directement} du fer aux ions cuivre. En additionnant membre à membre les deux demi-équations, les électrons s'éliminent (autant sont cédés que captés) et on obtient l'\textbf{équation-bilan} :

\[ \ce{Fe + Cu^2+ -> Fe^2+ + Cu} \]

Cette équation rend compte de toutes les observations : le clou se recouvre de cuivre rougeâtre, la solution bleue (ions \ce{Cu^2+}) pâlit en verdissement (ions \ce{Fe^2+}), et le clou s'amincit.

\begin{aretenir}[Une réaction d'oxydoréduction = deux couples]
Toute réaction d'oxydoréduction met en jeu \textbf{deux couples oxydant-réducteur} : le réducteur de l'un (Red$_1$) cède des électrons à l'oxydant de l'autre (Ox$_2$) :
\[ \ce{Red1 + Ox2 -> Ox1 + Red2} \]
Pour obtenir l'équation-bilan, on écrit les deux demi-équations, on les multiplie si nécessaire par des coefficients pour que le nombre d'électrons cédés égale le nombre d'électrons captés, puis on les additionne : \textbf{les électrons ne doivent jamais apparaître dans l'équation-bilan finale}.
\end{aretenir}

\begin{exemplebox}[Quand les nombres d'électrons diffèrent]
L'argent se dépose sur un fil de cuivre plongé dans une solution de nitrate d'argent. Les couples sont \ce{Cu^2+}/Cu et \ce{Ag+}/Ag. Le cuivre cède \textbf{deux} électrons, mais chaque ion \ce{Ag+} n'en capte qu'\textbf{un}. Il faut donc multiplier la demi-équation de l'argent par deux :
\begin{align*}
\ce{Cu &-> Cu^2+ + 2e-} \\
\ce{2Ag+ + 2e- &-> 2Ag}
\end{align*}
d'où le bilan : $\ce{Cu + 2Ag+ -> Cu^2+ + 2Ag}$. Le fil de cuivre se couvre de beaux cristaux argentés et la solution bleuit lentement.
\end{exemplebox}

\begin{attention}[Les électrons n'apparaissent jamais dans le bilan]
Une équation-bilan d'oxydoréduction contenant des électrons est \textbf{fausse} : c'est le signe que les deux demi-équations n'ont pas été correctement combinées. Vérifie toujours, avant d'additionner, que le nombre d'électrons cédés par le réducteur est exactement égal au nombre d'électrons captés par l'oxydant.
\end{attention}

\begin{savaistu}[D'où vient le mot « rédox » ?]
Les électriciens et les électrochimistes parlent souvent de « réactions rédox » ou de « potentiel rédox ». Ce mot, qui n'existe pas dans le dictionnaire classique, est simplement la contraction de \textbf{RÉD}uction et \textbf{OX}ydation — les deux moitiés inséparables de toute réaction d'oxydoréduction. Il rappelle à lui seul la leçon essentielle de cette section : on ne peut pas réduire sans oxyder, ni oxyder sans réduire. C'est aussi ce principe qui fait fonctionner les batteries des téléphones portables que l'on recharge chaque soir à Yaoundé comme à Douala : au sein de la batterie, deux couples rédox échangent des électrons... que le chargeur force à circuler dans le circuit !
\end{savaistu}

\begin{aretenir}[L'essentiel de la section]
\begin{itemize}
\item Une \textbf{oxydation} est une perte d'électrons ; une \textbf{réduction} est un gain d'électrons. L'\textbf{oxydant} capte, le \textbf{réducteur} cède.
\item Un \textbf{couple oxydant-réducteur} Ox/Red associe un oxydant et son réducteur conjugué ; l'oxydant s'écrit \textbf{toujours en premier}.
\item La \textbf{demi-équation électronique} s'écrit $\text{Ox} + n\,e^- \ce{<=>} \text{Red}$ ; on l'équilibre d'abord pour les éléments, puis pour les charges avec les électrons.
\item On distingue les \textbf{couples métalliques} \ce{M^n+}/M (métal solide + cation) des couples dont les deux formes sont en solution, comme \ce{Fe^3+}/\ce{Fe^2+}.
\item Toute réaction d'oxydoréduction met en jeu \textbf{deux couples} ; l'équation-bilan s'obtient en combinant les deux demi-équations de façon à éliminer les électrons.
\end{itemize}
\end{aretenir}

Nous disposons désormais de l'outil « couple ». Mais une question demeure : face à deux couples donnés, \emph{quel} réducteur cède ses électrons à \emph{quel} oxydant ? Pourquoi le fer attaque-t-il les ions \ce{Cu^2+} et pas l'inverse ? C'est ce que l'expérience va nous révéler dans la section suivante, en étudiant la réaction entre un métal et un ion métallique.

\coursec{Réaction entre un métal et un ion métallique : couples $\ce{M^{n+}}/\ce{M}$}

Dans la section précédente, tu as découvert la notion de \cle{couple oxydant-réducteur} : deux espèces chimiques, un oxydant et un réducteur, reliées par une demi-équation électronique du type $\ce{Ox + n e- <=> Red}$. Mais une question fondamentale reste en suspens : \emph{comment savoir si une réaction d'oxydoréduction peut réellement se produire entre deux espèces données ?} Un bijoutier de Douala qui plonge un fil de cuivre dans une solution contenant des ions argent voit un magnifique dépôt gris se former. En revanche, un clou en fer plongé dans une solution contenant des ions zinc ne subit aucune transformation. Pourquoi ? C'est toute la problématique de cette section : en étudiant expérimentalement les réactions entre un \cle{métal} et un \cle{ion métallique}, nous allons découvrir que certaines réactions sont \cle{spontanées} et d'autres impossibles, et en tirer une règle de comparaison des couples $\ce{M^{n+}}/\ce{M}$.

\courssub{Expérience 1 : le fer déplace le cuivre de sa solution}

\begin{experience}[Lame de fer dans une solution de sulfate de cuivre(II)]
\textbf{Dispositif et protocole.} On introduit une lame de fer bien décapée (à la laine de fer) dans un tube à essai contenant une solution de sulfate de cuivre(II) $\ce{CuSO4}$, de couleur bleue caractéristique due aux ions $\ce{Cu^2+}$.

\textbf{Observations.} Au bout de quelques minutes :
\begin{itemize}
\item la lame de fer se recouvre d'un \textbf{dépôt rouge-orangé} : c'est du cuivre métallique $\ce{Cu}$ ;
\item la solution, initialement bleue, \textbf{se décolore progressivement} (les ions $\ce{Cu^2+}$ disparaissent) ;
\item un test à la soude sur un prélèvement de la solution donne un précipité vert, caractéristique des ions $\ce{Fe^2+}$ : le fer est passé en solution sous forme d'ions fer(II).
\end{itemize}

\textbf{Interprétation.} Le fer métal a cédé des électrons (il a été oxydé en $\ce{Fe^2+}$) et les ions cuivre(II) ont capté ces électrons (ils ont été réduits en cuivre métal). Il s'agit bien d'une réaction d'\cle{oxydoréduction}.
\end{experience}

Écrivons les deux demi-équations électroniques mises en jeu :

\begin{align*}
\text{Oxydation du fer : } & \ce{Fe -> Fe^2+ + 2 e-} \\
\text{Réduction des ions cuivre : } & \ce{Cu^2+ + 2 e- -> Cu}
\end{align*}

Le nombre d'électrons cédés (2) est égal au nombre d'électrons captés (2) : on additionne directement les deux demi-équations pour obtenir l'équation-bilan :

\[ \ce{Fe + Cu^2+ -> Fe^2+ + Cu} \]

Cette réaction est \textbf{totale} et \textbf{spontanée} à température ambiante, sans aucun apport d'énergie extérieur. Les deux couples mis en jeu sont :

\begin{center}
\ce{Cu^2+}/\ce{Cu} \qquad \text{et} \qquad \ce{Fe^2+}/\ce{Fe}
\end{center}

\begin{attention}[Les ions spectateurs]
Dans l'équation-bilan, les ions sulfate $\ce{SO4^2-}$ n'apparaissent pas : ils ne participent pas à la réaction, ce sont des \cle{ions spectateurs}. On écrit donc l'équation sous forme \emph{ionique}, et non $\ce{Fe + CuSO4 -> FeSO4 + Cu}$ (forme « moléculaire » acceptée mais qui masque le mécanisme réel). Retiens : en oxydoréduction, on raisonne toujours sur les espèces qui échangent réellement des électrons.
\end{attention}

\courssub{Expérience 2 : le cuivre ne déplace pas le fer}

Réalisons maintenant l'expérience « inverse ».

\begin{experience}[Lame de cuivre dans une solution de sulfate de fer(II)]
\textbf{Protocole.} On plonge une lame de cuivre bien décapée dans une solution de sulfate de fer(II) $\ce{FeSO4}$, de couleur vert pâle.

\textbf{Observation.} Même après plusieurs heures, \textbf{aucun changement n'est observable} : pas de dépôt sur le cuivre, pas de modification de la couleur de la solution. Le test à la soude ne révèle aucun ion $\ce{Cu^2+}$.

\textbf{Interprétation.} La réaction $\ce{Cu + Fe^2+ -> Cu^2+ + Fe}$ \textbf{ne se produit pas}. Le cuivre métal est incapable de céder ses électrons aux ions $\ce{Fe^2+}$.
\end{experience}

\begin{popfigure}
\begin{center}
\begin{tikzpicture}[scale=0.9, every node/.style={font=\small}]
% Tube 1 : Fe dans CuSO4
\draw[thick, popdark] (-1.6,0) -- (-1.6,3.4) arc[start angle=180, end angle=360, x radius=0.8, y radius=0.5] -- (0,3.4);
\fill[popblue!25] (-1.55,0.35) arc[start angle=180, end angle=360, x radius=0.775, y radius=0.45] -- (1.55-1.6,2.6) -- (-1.55,2.6) -- cycle;
\draw[popblue, thick] (-1.55,2.6) -- (1.55-1.6,2.6);
\fill[popmuted] (-0.35,0.6) rectangle (-0.05,3.9);
\fill[poporange] (-0.38,0.6) rectangle (-0.02,2.3);
\node[popdark] at (0,4.15) {lame de Fe};
\node[align=left, popdark, anchor=west] at (1.2,2.2) {dépôt rouge de $\ce{Cu}$};
\draw[->, poporange, thick] (1.2,2.05) -- (0.05,1.5);
\node[align=left, popblue, anchor=west] at (1.2,0.9) {solution bleue qui\\ se décolore ($\ce{Cu^2+} \to \ce{Fe^2+}$)};
\node[align=center, popdark, font=\bfseries] at (0,-1.1) {Réaction : $\ce{Fe + Cu^2+ -> Fe^2+ + Cu}$};

% Tube 2 : Cu dans FeSO4
\begin{scope}[xshift=8.5cm]
\draw[thick, popdark] (-1.6,0) -- (-1.6,3.4) arc[start angle=180, end angle=360, x radius=0.8, y radius=0.5] -- (0,3.4);
\fill[popgreen!20] (-1.55,0.35) arc[start angle=180, end angle=360, x radius=0.775, y radius=0.45] -- (1.55-1.6,2.6) -- (-1.55,2.6) -- cycle;
\draw[popgreen!60!black, thick] (-1.55,2.6) -- (1.55-1.6,2.6);
\fill[poporange!80!black] (-0.35,0.6) rectangle (-0.05,3.9);
\node[popdark] at (0,4.15) {lame de Cu};
\node[align=left, popdark, anchor=west] at (1.2,2.2) {aucun dépôt};
\node[align=left, popgreen!50!black, anchor=west] at (1.2,0.9) {solution vert pâle\\ inchangée ($\ce{Fe^2+}$)};
\node[align=center, popdark, font=\bfseries] at (0,-1.1) {Aucune réaction};
\end{scope}
\end{tikzpicture}
\end{center}
\legende{À gauche : une lame de fer plongée dans une solution de sulfate de cuivre(II) se recouvre d'un dépôt rouge de cuivre et la solution se décolore. À droite : une lame de cuivre dans une solution de sulfate de fer(II) ne subit aucune transformation.}
\end{popfigure}

\courssub{Conclusion : une réaction spontanée n'a lieu que dans un sens}

La comparaison des deux expériences est capitale :

\begin{itemize}
\item le fer \textbf{réduit} les ions $\ce{Cu^2+}$ : la réaction $\ce{Fe + Cu^2+ -> Fe^2+ + Cu}$ se produit spontanément ;
\item le cuivre \textbf{ne réduit pas} les ions $\ce{Fe^2+}$ : la réaction inverse est impossible.
\end{itemize}

\begin{aretenir}[Règle de comparaison des couples]
Si le métal $\ce{M1}$ réduit l'ion $\ce{M2^{n+}}$ (c'est-à-dire si la réaction $\ce{M1 + M2^{n+} -> M1^{n+} + M2}$ se produit spontanément), alors :
\begin{itemize}
\item $\ce{M1}$ est un \textbf{réducteur plus fort} que $\ce{M2}$ ;
\item $\ce{M2^{n+}}$ est un \textbf{oxydant plus fort} que $\ce{M1^{n+}}$.
\end{itemize}
Autrement dit : \emph{le métal qui déplace un autre métal de sa solution est le réducteur le plus fort.}
\end{aretenir}

Dans notre exemple : le fer déplace le cuivre, donc \textbf{le fer est un réducteur plus fort que le cuivre} et \textbf{l'ion $\ce{Cu^2+}$ est un oxydant plus fort que l'ion $\ce{Fe^2+}$}. On dit aussi que le couple $\ce{Cu^2+}/\ce{Cu}$ est « plus élevé » que le couple $\ce{Fe^2+}/\ce{Fe}$ dans la classification que nous construirons à la section 4.

\begin{propriete}[Réaction spontanée entre deux couples]
Une réaction d'oxydoréduction spontanée se produit toujours entre le \textbf{réducteur d'un couple} et l'\textbf{oxydant d'un autre couple}, à condition que le réducteur soit \emph{plus fort} que le réducteur conjugué de cet oxydant. Le réducteur fort « impose » ses électrons à l'oxydant fort ; les espèces formées (réducteur faible et oxydant faible) ne peuvent pas réagir en sens inverse.
\end{propriete}

\courssub{Expérience 3 : l'arbre de Diane, le cuivre déplace l'argent}

\begin{experience}[Fil de cuivre dans une solution de nitrate d'argent]
\textbf{Protocole.} On plonge un fil de cuivre (par exemple un fil de cuivre enroulé en spirale) dans une solution incolore de nitrate d'argent $\ce{AgNO3}$.

\textbf{Observations.} Après quelques minutes, le fil de cuivre se recouvre d'un magnifique \textbf{dépôt gris-argenté, cristallisé en aiguilles}, qui évoque les branches d'un arbre : c'est le célèbre « \cle{arbre de Diane} ». Simultanément, la solution, initialement incolore, \textbf{bleuit progressivement} : des ions $\ce{Cu^2+}$ apparaissent.

\textbf{Interprétation.} Le cuivre métal est oxydé en ions $\ce{Cu^2+}$ et les ions argent $\ce{Ag+}$ sont réduits en argent métal $\ce{Ag}$. Les couples mis en jeu sont $\ce{Ag+}/\ce{Ag}$ et $\ce{Cu^2+}/\ce{Cu}$.
\end{experience}

Écrivons les demi-équations :

\begin{align*}
\text{Oxydation du cuivre : } & \ce{Cu -> Cu^2+ + 2 e-} \\
\text{Réduction des ions argent : } & \ce{Ag+ + e- -> Ag}
\end{align*}

Ici, le cuivre cède \textbf{deux} électrons alors que chaque ion argent n'en capte qu'\textbf{un}. Il faut donc multiplier la deuxième demi-équation par 2 avant d'additionner :

\[ \ce{Cu + 2 Ag+ -> Cu^2+ + 2 Ag} \]

Le cuivre déplace l'argent de sa solution : \textbf{le cuivre est un réducteur plus fort que l'argent}, et l'ion $\ce{Ag+}$ est un oxydant plus fort que l'ion $\ce{Cu^2+}$.

\begin{attention}[Erreur fréquente : l'équilibrage des électrons]
Le piège classique consiste à écrire $\ce{Cu + Ag+ -> Cu^2+ + Ag}$. Cette équation est \textbf{fausse} : elle n'est équilibrée ni en charge ($+1$ à gauche, $+2$ à droite) ni en électrons échangés. Vérifie toujours que le nombre d'électrons cédés par le réducteur est \emph{exactement égal} au nombre d'électrons captés par l'oxydant. Ici : 2 électrons cédés par $\ce{Cu}$, donc 2 ions $\ce{Ag+}$ nécessaires.
\end{attention}

\begin{popfigure}
\begin{center}
\begin{tikzpicture}[scale=0.95, every node/.style={font=\small}]
% Bécher
\draw[thick, popdark] (-2,0) -- (-2,4.2) (2,0) -- (2,4.2);
\draw[thick, popdark] (-2,0) -- (2,0);
\draw[thick, popdark] (-2,4.2) -- (-2.25,4.5) (2,4.2) -- (2.25,4.5);
% solution
\fill[popblue!15] (-1.95,0.05) rectangle (1.95,3.4);
\draw[popblue!60, thick] (-1.95,3.4) -- (1.95,3.4);
\node[popblue!70!black, anchor=west] at (2.35,3.4) {solution de $\ce{AgNO3}$};
\node[popblue!70!black, anchor=west] at (2.35,2.9) {qui bleuit ($\ce{Cu^2+}$)};
% fil de cuivre
\draw[very thick, poporange!80!black] (0,4.6) -- (0,1.2);
% cristaux d'argent (aiguilles)
\foreach \a/\l in {20/0.9, 55/0.75, 90/1.0, 125/0.8, 160/0.9, 200/0.85, 235/0.7, 270/0.95, 305/0.8, 340/0.9}{
  \draw[popmuted!60!black, thick] (0,2.2) -- (\a:\l);
}
\foreach \a/\l in {30/0.7, 75/0.85, 110/0.65, 150/0.75, 190/0.7, 250/0.8, 290/0.65, 330/0.75}{
  \draw[popmuted!60!black, thick] (0,1.5) -- (\a:\l);
}
\foreach \a/\l in {45/0.6, 135/0.6, 225/0.6, 315/0.6}{
  \draw[popmuted!60!black, thick] (0,2.8) -- (\a:\l);
}
\node[poporange!80!black, anchor=east] at (-0.15,4.3) {fil de cuivre};
\node[popmuted!40!black, anchor=east, align=right] at (-2.3,1.9) {dépôt gris d'argent\\ cristallisé (« arbre »)};
\draw[->, popmuted!60!black, thick] (-2.3,1.75) -- (-0.7,1.7);
\node[align=center, popdark, font=\bfseries] at (0,-0.8) {$\ce{Cu + 2 Ag+ -> Cu^2+ + 2 Ag}$};
\end{tikzpicture}
\end{center}
\legende{L'arbre de Diane : un fil de cuivre plongé dans une solution de nitrate d'argent se recouvre de cristaux d'argent gris brillant, tandis que la solution bleuit par formation d'ions $\ce{Cu^2+}$.}
\end{popfigure}

\begin{savaistu}[L'arbre de Diane, merveille des alchimistes]
Dès le Moyen Âge, les alchimistes connaissaient cette réaction spectaculaire. Ils l'appelaient « arbre de Diane », du nom de la déesse romaine de la Lune, car l'argent était le métal associé à la Lune dans leur symbolique. Ils y voyaient une « croissance végétale » du métal, preuve selon eux que les métaux pouvaient « pousser » et se transmuter. Aujourd'hui, cette réaction illustre simplement la différence de pouvoir réducteur entre le cuivre et l'argent. C'est d'ailleurs un principe voisin qui est utilisé dans la récupération des métaux précieux : on peut faire précipiter l'argent contenu dans des bains de traitement de photographies usagés en y plongeant du cuivre !
\end{savaistu}

\courssub{Méthode d'interprétation et exemple chiffré}

\begin{methode}[Interpréter une expérience métal / ion métallique]
Face à l'expérience « métal $\ce{M1}$ plongé dans une solution contenant des ions $\ce{M2^{n+}}$ », applique systématiquement les quatre étapes suivantes :
\begin{enumerate}
\item \textbf{Observer} : y a-t-il un dépôt sur le métal ? La couleur de la solution change-t-elle ? Si rien ne se passe, conclure immédiatement : pas de réaction, $\ce{M1}$ est un réducteur plus faible que $\ce{M2}$.
\item \textbf{Identifier les produits} : le dépôt est le métal $\ce{M2}$ ; le métal $\ce{M1}$ passe en solution sous forme d'ions $\ce{M1^{n+}}$ (à confirmer par un test caractéristique, par exemple à la soude).
\item \textbf{Écrire les demi-équations} : oxydation $\ce{M1 -> M1^{n+} + n e-}$ et réduction $\ce{M2^{n+} + n e- -> M2}$.
\item \textbf{Équilibrer puis additionner} : multiplier chaque demi-équation pour égaliser les électrons échangés, puis écrire le bilan et conclure sur la force comparée des réducteurs.
\end{enumerate}
\end{methode}

\begin{exoresolu}[Masse de cuivre déposée par le fer]
\textbf{Énoncé.} On plonge une lame de fer de masse $m(\ce{Fe}) = \SI{1,12}{g}$ dans un volume important d'une solution de sulfate de cuivre(II) en excès. La lame est entièrement consommée. Calculer la masse de cuivre déposée sur la lame avant sa disparition. On donne $M(\ce{Fe}) = \SI{56}{g.mol^{-1}}$ et $M(\ce{Cu}) = \SI{64}{g.mol^{-1}}$.
\tcblower
\textbf{Solution détaillée.}

\textbf{Étape 1 : équation-bilan.} La réaction spontanée est :
\[ \ce{Fe + Cu^2+ -> Fe^2+ + Cu} \]

\textbf{Étape 2 : tableau d'avancement.} Les ions $\ce{Cu^2+}$ étant en excès, le fer est le réactif limitant.

\begin{center}
\begin{tabular}{lcccc}
\toprule
 & $\ce{Fe}$ & $\ce{Cu^2+}$ & $\ce{Fe^2+}$ & $\ce{Cu}$ \\
\midrule
État initial (mol) & $n_0$ & excès & $0$ & $0$ \\
État final (mol) & $0$ & excès & $n_0$ & $n_0$ \\
\bottomrule
\end{tabular}
\end{center}

\textbf{Étape 3 : quantité de fer consommée.}
\[ n_0 = \frac{m(\ce{Fe})}{M(\ce{Fe})} = \frac{1,12}{56} = \SI{0,020}{mol} \]

\textbf{Étape 4 : masse de cuivre déposée.} D'après la stœchiométrie, $n(\ce{Cu}) = n_0 = \SI{0,020}{mol}$, d'où :
\[ m(\ce{Cu}) = n(\ce{Cu}) \times M(\ce{Cu}) = 0,020 \times 64 = \SI{1,28}{g} \]

\textbf{Conclusion.} Il se dépose $\SI{1,28}{g}$ de cuivre. Remarque : la masse de cuivre déposée ($\SI{1,28}{g}$) est supérieure à la masse de fer consommée ($\SI{1,12}{g}$) car le cuivre a une masse molaire plus grande que le fer, à quantité de matière égale.
\end{exoresolu}

\begin{exercice}[À toi de jouer]
On plonge un fil de cuivre de masse $m = \SI{0,64}{g}$ dans une solution de nitrate d'argent en excès. Le fil est entièrement consommé. Calculer la masse d'argent déposée. On donne $M(\ce{Cu}) = \SI{64}{g.mol^{-1}}$ et $M(\ce{Ag}) = \SI{108}{g.mol^{-1}}$.
\end{exercice}

\begin{corrige}[Exercice]
L'équation-bilan est $\ce{Cu + 2 Ag+ -> Cu^2+ + 2 Ag}$. La quantité de cuivre consommée est :
\[ n(\ce{Cu}) = \frac{0,64}{64} = \SI{0,010}{mol} \]
D'après la stœchiométrie, $n(\ce{Ag}) = 2\, n(\ce{Cu}) = \SI{0,020}{mol}$ (attention au coefficient 2 !). Donc :
\[ m(\ce{Ag}) = 0,020 \times 108 = \SI{2,16}{g} \]
Il se dépose $\SI{2,16}{g}$ d'argent. L'erreur à éviter était d'oublier le facteur 2 issu de l'équilibrage des électrons.
\end{corrige}

En résumé, les expériences métal/ion métallique nous ont appris que les réactions d'oxydoréduction sont \textbf{sélectives} : elles ne se produisent que dans un sens, celui où le réducteur le plus fort cède ses électrons à l'oxydant le plus fort. Nous avons déjà pu classer trois couples : $\ce{Fe^2+}/\ce{Fe}$, $\ce{Cu^2+}/\ce{Cu}$ et $\ce{Ag+}/\ce{Ag}$, du réducteur le plus fort (Fe) au réducteur le plus faible (Ag). Mais où placer l'hydrogène, ce non-métal si particulier, dans cette hiérarchie ? C'est l'objet de la section suivante, consacrée à l'action des ions $\ce{H3O+}$ sur les métaux.

\coursec{Action des ions \ce{H3O+} sur les métaux : le couple \ce{H3O+}/\ce{H2}}

Dans la section précédente, nous avons vu qu'un métal peut réduire les ions d'un autre métal : le zinc réduit les ions \ce{Cu^2+}, le cuivre réduit les ions \ce{Ag+}, etc. Ces expériences ont fait apparaître des couples du type \ce{M^n+}/\ce{M}. Mais une question surgit naturellement : que se passe-t-il lorsqu'on verse un \emph{acide} sur un métal ? Tout le monde a vu, au laboratoire du lycée, le zinc « bouillonner » dans l'acide chlorhydrique en dégageant un gaz. Ce gaz, c'est du dihydrogène. Cette observation banale cache une information précieuse : les ions oxonium \ce{H3O+}, présents dans toute solution acide, constituent l'oxydant d'un nouveau couple, le couple \ce{H3O+}/\ce{H2}, qui va nous servir de \cle{référence} pour classer tous les couples métalliques. C'est l'objet de cette section.

\courssub{Expérience : action de l'acide chlorhydrique sur quelques métaux}

\begin{experience}[Action de l'acide chlorhydrique dilué sur le zinc, le fer, l'aluminium, le cuivre et l'argent]
\textbf{Dispositif et protocole.} On dispose de cinq tubes à essais contenant chacun quelques millilitres d'une solution diluée d'acide chlorhydrique $(\ce{H3O+} + \ce{Cl-})$. On introduit dans chaque tube un petit morceau de métal différent : zinc (grenaille), fer (clou décapé), aluminium (récupéré sur une feuille de marmite), cuivre (fil de câble électrique dénudé) et argent (bijou ou fil). On bouche éventuellement le tube et on recueille le gaz dégagé, que l'on approche d'une flamme.

\textbf{Observations.}
\begin{itemize}
\item Avec le \textbf{zinc} : réaction vive, effervescence abondante, le tube chauffe légèrement ; le gaz recueilli produit une \emph{légère détonation} (un « aboiement ») à l'approche de la flamme.
\item Avec le \textbf{fer} : dégagement gazeux net mais plus lent ; même test positif à la flamme.
\item Avec l'\textbf{aluminium} : réaction parfois lente à démarrer (à cause de la couche d'alumine protectrice), puis effervescence ; même test positif.
\item Avec le \textbf{cuivre} et l'\textbf{argent} : \emph{aucune réaction observable}, même après plusieurs minutes.
\end{itemize}

\textbf{Interprétation.} Le gaz qui détone est le \cle{dihydrogène} \ce{H2}. Les métaux zinc, fer et aluminium ont donc \emph{réduit} les ions \ce{H3O+} en dihydrogène, tandis qu'ils ont été eux-mêmes \emph{oxydés} en ions métalliques (\ce{Zn^2+}, \ce{Fe^2+}, \ce{Al^3+}). Le cuivre et l'argent, eux, sont incapables de réduire \ce{H3O+}.
\end{experience}

\begin{popfigure}
\begin{center}
\begin{tikzpicture}[scale=0.9, every node/.style={font=\small}]
% Tube à essais (réaction)
\draw[thick, popdark] (-0.9,0) -- (-0.9,-3.4) arc[start angle=180, end angle=360, radius=0.9] -- (0.9,0);
% Acide (liquide)
\fill[popblueL, opacity=0.7] (-0.86,-1.2) -- (-0.86,-3.35) arc[start angle=180, end angle=360, radius=0.86] -- (0.86,-1.2) -- cycle;
\draw[popblue, thick] (-0.86,-1.2) -- (0.86,-1.2);
% Grenaille de zinc au fond
\foreach \x/\y in {-0.45/-3.0,-0.15/-3.15,0.2/-2.95,0.45/-3.1,0.0/-2.8}
  \fill[popmuted] (\x,\y) circle (0.12);
% Bulles d'H2 qui montent
\foreach \x/\y/\r in {-0.3/-2.2/0.07,0.25/-1.9/0.09,0.0/-1.55/0.06,-0.2/-1.75/0.08,0.35/-2.4/0.06}
  \draw[popblue] (\x,\y) circle (\r);
% Étiquettes à droite du tube réaction
\node[align=left, anchor=west] at (1.3,-3.0) {grenaille de zinc \ce{Zn}};
\draw[->, popmuted] (1.3,-3.0) -- (0.3,-3.05);
\node[align=left, anchor=west] at (1.3,-1.6) {acide chlorhydrique dilué $(\ce{H3O+} + \ce{Cl-})$};
\draw[->, popblue] (1.3,-1.5) -- (0.6,-1.25);
\node[align=left, anchor=west] at (1.3,-0.5) {dégagement de \ce{H2}};
\draw[->, poporange] (1.3,-0.6) -- (0.25,-1.05);
% Flèche gaz vers tube recueil
\draw[->, very thick, poporange] (0,0.15) -- (0,0.9) -- (2.6,0.9);
% Tube de recueil renversé
\draw[thick, popdark] (2.6,1.6) -- (2.6,0.35) arc[start angle=180, end angle=360, radius=0.45] -- (3.5,1.6);
% Eau déplacée dans tube recueil
\fill[popgoldL, opacity=0.6] (2.62,0.5) arc[start angle=180, end angle=360, radius=0.43] -- (3.48,1.1) -- (2.62,1.1) -- cycle;
\node[anchor=west, align=left] at (3.8,1.0) {gaz recueilli : \ce{H2}};
% Flamme
\draw[fill=poporange, poporange] (5.6,0.15) .. controls (5.35,0.55) and (5.6,0.75) .. (5.6,1.05)
   .. controls (5.6,0.75) and (5.85,0.55) .. (5.6,0.15);
\draw[fill=poppink, poppink] (5.6,0.2) .. controls (5.45,0.45) and (5.6,0.55) .. (5.6,0.75)
   .. controls (5.6,0.55) and (5.75,0.45) .. (5.6,0.2);
\node[anchor=west, align=left] at (5.95,0.55) {flamme :\\ détonation (« aboiement »)};
% Flèche approche flamme (mouvement de la flamme vers le tube)
\draw[<-, poporange] (5.55,0.9) -- (3.6,0.85);
\end{tikzpicture}
\end{center}
\legende{Action de l'acide chlorhydrique dilué sur le zinc : dégagement de dihydrogène, recueilli et identifié par le test de la détonation à la flamme.}
\end{popfigure}

\courssub{Identification du gaz : le test de la détonation}

\begin{methode}[Identifier le dihydrogène : le test de l'aboiement]
\begin{enumerate}
\item Recueillir le gaz dégagé dans un petit tube à essais renversé (le dihydrogène est beaucoup moins dense que l'air, il « monte »).
\item Approcher rapidement la flamme d'une bougie ou d'une allumette de l'ouverture du tube.
\item Observer : si le gaz est du dihydrogène, on entend une \textbf{légère détonation}, caractéristique, appelée « aboiement ». Le dihydrogène s'enflamme au contact de l'air selon : \[ \ce{2H2 + O2 -> 2H2O} \] réaction très exothermique, d'où la détonation.
\end{enumerate}
\textbf{Précaution} : travailler avec de petites quantités de gaz et un acide \emph{dilué}. Un mélange \ce{H2}/air en grand volume est explosif.
\end{methode}

\courssub{Interprétation : les couples mis en jeu}

Prenons le cas du zinc. Le métal zinc disparaît progressivement : il est \emph{oxydé} en ions \ce{Zn^2+} qui passent en solution (on peut d'ailleurs les caractériser ensuite par la soude, qui donne un précipité blanc d'hydroxyde de zinc). Les ions oxonium, eux, sont \emph{réduits} en dihydrogène. Deux couples sont donc en jeu : \ce{Zn^2+}/\ce{Zn} et \ce{H3O+}/\ce{H2}.

Écrivons les demi-équations :
\begin{align*}
\text{oxydation du zinc : } & \ce{Zn -> Zn^2+ + 2e-}\\
\text{réduction de l'ion oxonium : } & \ce{2H3O+ + 2e- -> H2 + 2H2O}
\end{align*}
En additionnant membre à membre (les deux électrons échangés s'éliminent), on obtient l'équation-bilan :
\[ \ce{Zn + 2H3O+ -> Zn^2+ + H2 ^ + 2H2O} \]

Cette réaction est \textbf{totale}, \textbf{exothermique} (le tube chauffe) et se déroule à température ambiante, sans catalyseur. Remarque que les ions chlorure \ce{Cl-} ne participent pas : ce sont des \cle{ions spectateurs}.

\begin{definition}[Le couple \ce{H3O+}/\ce{H2}]
L'ion oxonium \ce{H3O+} est l'\textbf{oxydant} et le dihydrogène \ce{H2} est le \textbf{réducteur} du couple \ce{H3O+}/\ce{H2}, dont la demi-équation électronique s'écrit :
\[ \ce{2H3O+ + 2e- <=> H2 + 2H2O} \]
Ce couple joue un rôle de \textbf{référence} dans la classification électrochimique : c'est par rapport à lui que l'on situe tous les couples métalliques \ce{M^n+}/\ce{M}.
\end{definition}

\begin{attention}[Erreurs fréquentes sur la demi-équation du couple \ce{H3O+}/\ce{H2}]
\begin{itemize}
\item \textbf{Ne jamais écrire \ce{H+} en solution aqueuse.} Le proton \ce{H+} n'existe pas librement dans l'eau : il se fixe sur une molécule d'eau pour donner l'ion oxonium \ce{H3O+}. Écrire $\ce{2H+} + 2e^- \rightleftharpoons \ce{H2}$ est toléré dans certains ouvrages étrangers, mais au programme camerounais, c'est bien \ce{H3O+} qu'il faut écrire.
\item \textbf{Ne pas oublier les molécules d'eau.} La demi-équation correcte est \ce{2H3O+ + 2e- <=> H2 + 2H2O}. Si tu écris \ce{2H3O+ + 2e- <=> H2}, l'équation n'est équilibrée ni en oxygène ni en hydrogène.
\item \textbf{Vérifier toujours la double conservation} : conservation des éléments (6 H à gauche, $2 + 4 = 6$ à droite ; 2 O de chaque côté) et conservation de la charge ($+2 - 2 = 0$ à gauche, $0$ à droite).
\end{itemize}
\end{attention}

\courssub{Exemples chiffrés et équilibrage des équations}

\begin{exemplebox}[Équilibrer l'action de l'acide sur le fer]
Le fer s'oxyde en ions \ce{Fe^2+}. Cherchons l'équation-bilan de la réaction entre le fer et l'acide chlorhydrique dilué.

\textbf{Étape 1 — demi-équations :}
\[ \ce{Fe -> Fe^2+ + 2e-} \qquad \ce{2H3O+ + 2e- -> H2 + 2H2O} \]
\textbf{Étape 2 — les électrons échangés sont déjà en même nombre (2)} : on additionne directement.
\textbf{Étape 3 — bilan :}
\[ \ce{Fe + 2H3O+ -> Fe^2+ + H2 ^ + 2H2O} \]
Vérification : 1 Fe de chaque côté ; 6 H à gauche et $2+4=6$ à droite ; charge $+2$ de chaque côté. L'équation est correcte.
\end{exemplebox}

\begin{exoresolu}[Action de l'acide chlorhydrique sur l'aluminium]
On plonge une plaque d'aluminium de masse $m = \SI{2,7}{g}$ dans un excès d'acide chlorhydrique dilué. L'aluminium s'oxyde en ions \ce{Al^3+}.\\
1) Écrire l'équation-bilan de la réaction.\\
2) Calculer le volume de dihydrogène dégagé, mesuré dans les conditions où le volume molaire vaut $V_m = \SI{24,0}{L.mol^{-1}}$.\\
On donne : $M(\ce{Al}) = \SI{27,0}{g.mol^{-1}}$.
\tcblower
\textbf{1) Équation-bilan.} Demi-équations :
\[ \ce{Al -> Al^3+ + 3e-} \qquad \ce{2H3O+ + 2e- -> H2 + 2H2O} \]
Pour égaliser les électrons, on multiplie la première par 2 et la seconde par 3 (6 électrons échangés) :
\[ \ce{2Al + 6H3O+ -> 2Al^3+ + 3H2 ^ + 6H2O} \]
\textbf{2) Volume de dihydrogène.} Quantité d'aluminium :
\[ n(\ce{Al}) = \frac{m}{M} = \frac{2,7}{27,0} = \SI{0,10}{mol} \]
D'après les coefficients stœchiométriques, $n(\ce{H2}) = \dfrac{3}{2}\,n(\ce{Al}) = \SI{0,15}{mol}$.
\[ V(\ce{H2}) = n \times V_m = 0,15 \times 24,0 = \SI{3,6}{L} \]
La plaque d'aluminium dégage donc \SI{3,6}{L} de dihydrogène : c'est un volume considérable, qui explique l'effervescence spectaculaire observée.
\end{exoresolu}

\courssub{L'expérience inverse : le cuivre et l'argent ne réagissent pas}

Reprenons nos tubes : avec le cuivre et l'argent, rien ne se passe, même en chauffant légèrement, même en attendant longtemps. Aucune bulle, aucune détonation, aucune dissolution du métal. La réaction
\[ \ce{Cu + 2H3O+ -> Cu^2+ + H2 + 2H2O} \]
\emph{ne se produit pas} : le cuivre est incapable de céder ses électrons aux ions oxonium. Dit autrement, le cuivre est un réducteur \emph{plus faible} que le dihydrogène, tandis que le zinc, le fer et l'aluminium sont des réducteurs \emph{plus forts} que \ce{H2}.

\begin{aretenir}[Ce que l'expérience nous apprend]
\begin{itemize}
\item Les métaux \textbf{zinc, fer, aluminium} (comme le magnésium) \textbf{réduisent} les ions \ce{H3O+} : ils sont attaqués par les acides dilués non oxydants, avec dégagement de dihydrogène. On dit que leurs couples \ce{M^n+}/\ce{M} sont placés \emph{au-dessus} (côté réducteurs forts) du couple \ce{H3O+}/\ce{H2}.
\item Les métaux \textbf{cuivre, argent} (comme le mercure et l'or) \textbf{ne réduisent pas} \ce{H3O+} : ils résistent aux acides dilués non oxydants. Leurs couples sont placés \emph{en dessous} du couple \ce{H3O+}/\ce{H2}.
\item Le couple \ce{H3O+}/\ce{H2} apparaît ainsi comme une \textbf{frontière} dans la classification électrochimique : il sépare les métaux « attaquables par les acides » des métaux « nobles ».
\end{itemize}
\end{aretenir}

\begin{center}
\begin{tabularx}{\linewidth}{|l|X|X|}
\hline
\rowcolor{popblueL} \textbf{Métal} & \textbf{Observation avec \ce{H3O+} dilué} & \textbf{Équation-bilan} \\
\hline
Zinc \ce{Zn} & Effervescence vive, détonation & \ce{Zn + 2H3O+ -> Zn^2+ + H2 + 2H2O} \\
\hline
Fer \ce{Fe} & Dégagement modéré, détonation & \ce{Fe + 2H3O+ -> Fe^2+ + H2 + 2H2O} \\
\hline
Aluminium \ce{Al} & Effervescence après amorçage & \ce{2Al + 6H3O+ -> 2Al^3+ + 3H2 + 6H2O} \\
\hline
Cuivre \ce{Cu} & Aucune réaction & --- \\
\hline
Argent \ce{Ag} & Aucune réaction & --- \\
\hline
\end{tabularx}
\end{center}

\begin{attention}[Précautions expérimentales et vocabulaire]
\begin{itemize}
\item Toujours utiliser des acides \textbf{dilués} : un acide concentré réagit violemment avec les métaux et les projections sont dangereuses. Manipuler avec des lunettes et des gants.
\item Attention : l'acide nitrique \ce{HNO3} est un cas particulier. Même dilué, c'est l'ion nitrate \ce{NO3-} (oxydant puissant) qui attaque le cuivre, et non \ce{H3O+} : le gaz dégagé n'est alors pas \ce{H2} mais des oxydes d'azote. Notre classification ne s'applique qu'aux acides \emph{non oxydants} comme \ce{HCl} ou \ce{H2SO4} dilués.
\item Ne confonds jamais \ce{H3O+} (ion oxonium, oxydant du couple) avec \ce{H+} (notation simplifiée à éviter) ni avec \ce{H2} (le réducteur du couple).
\end{itemize}
\end{attention}

\begin{savaistu}[Le dihydrogène des piles artisanales]
L'action d'un acide sur un métal est parfois exploitée dans des démonstrations artisanales : en plongeant une lame de zinc et une lame de cuivre dans un citron, un jus de bissap acidulé ou un vinaigre (qui contiennent des ions \ce{H3O+}), on fabrique une \emph{pile} rudimentaire capable d'allumer une petite diode électroluminescente. Le zinc, réducteur plus fort que \ce{H2}, s'oxyde et fournit des électrons ; le cuivre, lui, ne sert que d'électrode collectrice. Ce principe — deux métaux de pouvoirs réducteurs différents plongés dans une solution acide — est exactement celui de la pile Volta, inventée en 1800, et tu le retrouveras en classe de Terminale. Le dihydrogène est d'ailleurs aujourd'hui étudié comme \emph{carburant propre} : sa combustion ne produit que de l'eau !
\end{savaistu}

En résumé, l'expérience de cette section nous a offert un \textbf{point de repère universel} : le couple \ce{H3O+}/\ce{H2}. Nous savons désormais placer les couples \ce{Zn^2+}/\ce{Zn}, \ce{Fe^2+}/\ce{Fe} et \ce{Al^3+}/\ce{Al} d'un côté de cette frontière, et les couples \ce{Cu^2+}/\ce{Cu} et \ce{Ag+}/\ce{Ag} de l'autre. Il ne reste plus, dans la section suivante, qu'à assembler toutes ces observations expérimentales en une \emph{classification électrochimique qualitative} cohérente.

\coursec{Construction de la classification électrochimique qualitative}

Dans la section précédente, nous avons vu que les ions \ce{H3O+} des solutions acides attaquent certains métaux (fer, zinc, aluminium) avec dégagement de dihydrogène, alors que le cuivre et l'argent restent inertes. Cette différence de comportement n'est pas un caprice de la nature : elle révèle une \cle{hiérarchie} entre les couples oxydant-réducteur. L'objectif de cette section est ambitieux et magnifique : à partir d'un petit nombre d'expériences simples, nous allons construire une \cle{échelle} sur laquelle tous les couples seront rangés, comme des équipes de football sur un classement. Cette échelle, appelée \cle{classification électrochimique}, nous permettra ensuite de \emph{prévoir} les réactions sans même les réaliser.

\courssub{Le principe : chaque réaction spontanée est un verdict}

\textbf{Intuition d'abord.}\par

Imagine deux lutteurs traditionnels du Nord-Cameroun qui s'affrontent : le vainqueur est déclaré « plus fort » que le vaincu. En oxydoréduction, c'est exactement pareil. Quand on met en présence deux couples, par exemple \ce{Cu^2+}/\ce{Cu} et \ce{Fe^2+}/\ce{Fe}, une réaction spontanée peut se produire :

\[ \ce{Cu^2+ + Fe -> Cu + Fe^2+} \]

Le fait que cette réaction se produise \emph{dans ce sens-là} (et pas dans l'autre) est un verdict expérimental : l'ion \ce{Cu^2+} a « arraché » des électrons au fer. On dit que \ce{Cu^2+} est un oxydant \textbf{plus fort} que \ce{Fe^2+}, et que le fer est un réducteur \textbf{plus fort} que le cuivre. Chaque expérience réussie impose donc une inégalité entre deux couples.

\begin{definition}[Classification électrochimique qualitative]
La \cle{classification électrochimique qualitative} est le classement des couples oxydant-réducteur selon le \textbf{pouvoir oxydant croissant de leurs oxydants} (ou, ce qui est équivalent, selon le pouvoir réducteur croissant de leurs réducteurs). On la construit expérimentalement : chaque réaction d'oxydoréduction spontanée observée impose une inégalité entre les deux couples mis en jeu.
\end{definition}

\begin{methode}[Construire la classification pas à pas]
\begin{enumerate}
\item \textbf{Réaliser l'expérience} : mettre en contact le réducteur d'un couple avec l'oxydant d'un autre couple (par exemple un clou en fer dans une solution de sulfate de cuivre).
\item \textbf{Observer} : s'il y a réaction (dépôt métallique, dégagement gazeux, changement de couleur), écrire l'équation bilan.
\item \textbf{Tirer le verdict} : l'oxydant qui a réagi est plus fort que l'oxydant de l'autre couple ; le réducteur qui a réagi est plus fort que l'autre réducteur.
\item \textbf{Cumuler les inégalités} : enchaîner les résultats de plusieurs expériences pour obtenir un classement complet, comme on enchaîne les résultats d'un tournoi.
\end{enumerate}
\end{methode}

\courssub{Les résultats expérimentaux cumulés}

Reprenons, sous forme de bilan, toutes les expériences réalisées dans les sections précédentes. Chacune d'elles est une « rencontre » dont nous connaissons le vainqueur.

\begin{center}
\begin{tabularx}{\linewidth}{|l|X|X|}
\hline
\rowcolor{popblueL} \textbf{Expérience} & \textbf{Observation} & \textbf{Verdict} \\
\hline
Fer dans une solution de \ce{Cu^2+} & Dépôt rouge de cuivre sur le clou : \ce{Cu^2+ + Fe -> Cu + Fe^2+} & \ce{Cu^2+} oxydant plus fort que \ce{Fe^2+} ; Fe réducteur plus fort que Cu \\
\hline
Cuivre dans une solution de \ce{Ag+} & Dépôt gris d'argent, solution qui bleuit : \ce{2Ag+ + Cu -> 2Ag + Cu^2+} & \ce{Ag+} oxydant plus fort que \ce{Cu^2+} ; Cu réducteur plus fort que Ag \\
\hline
Zinc dans une solution de \ce{Fe^2+} & Dépôt de fer sur le zinc : \ce{Fe^2+ + Zn -> Fe + Zn^2+} & \ce{Fe^2+} oxydant plus fort que \ce{Zn^2+} ; Zn réducteur plus fort que Fe \\
\hline
Zinc dans un acide dilué & Dégagement de dihydrogène : \ce{2H3O+ + Zn -> H2 + Zn^2+ + 2H2O} & \ce{H3O+} oxydant plus fort que \ce{Zn^2+} ; Zn réducteur plus fort que \ce{H2} \\
\hline
Fer dans un acide dilué & Dégagement de dihydrogène : \ce{2H3O+ + Fe -> H2 + Fe^2+ + 2H2O} & \ce{H3O+} oxydant plus fort que \ce{Fe^2+} ; Fe réducteur plus fort que \ce{H2} \\
\hline
Cuivre dans un acide dilué & \textbf{Aucune réaction} & \ce{Cu^2+} oxydant plus fort que \ce{H3O+} ; \ce{H2} réducteur plus fort que Cu \\
\hline
\end{tabularx}
\end{center}

\begin{attention}[Une absence de réaction est aussi un verdict]
Lorsque le cuivre plongé dans l'acide chlorhydrique ne réagit \emph{pas}, ce n'est pas une expérience ratée : c'est une information précieuse ! Elle prouve que la réaction inverse, \ce{Cu^2+ + H2 -> Cu + 2H3O+}, serait spontanée. Donc \ce{Cu^2+} est un oxydant plus fort que \ce{H3O+}. En oxydoréduction, « rien ne se passe » signifie « la réaction spontanée irait dans l'autre sens ».
\end{attention}

\courssub{Le classement obtenu}

Enchaînons les inégalités comme les maillons d'une corde. Pour les \textbf{réducteurs}, du plus faible au plus fort :

\[ \ce{Ag} < \ce{Cu} < \ce{H2} < \ce{Fe} < \ce{Zn} \]

D'autres expériences (que nous admettons) montrent que l'aluminium et le magnésium réduisent les ions \ce{Zn^2+} : ce sont des réducteurs \textbf{encore plus forts que le zinc}, ils se placent donc \textbf{en dessous du zinc} dans le classement (pouvoir réducteur croissant vers le bas).

Pour les \textbf{oxydants}, l'ordre est exactement inversé (c'est logique : un couple dont le réducteur est très fort possède un oxydant très faible, et réciproquement) :

\[ \ce{Zn^2+} < \ce{Fe^2+} < \ce{H3O+} < \ce{Cu^2+} < \ce{Ag+} \]

\begin{aretenir}[La règle d'or de la classification]
Dans un couple oxydant-réducteur, \textbf{plus l'oxydant est fort, plus le réducteur conjugué est faible}, et inversement. C'est pourquoi les deux échelles de la classification sont toujours tracées en sens opposés : la force des oxydants croît vers le haut, celle des réducteurs croît vers le bas.
\end{aretenir}

\begin{popfigure}
\begin{center}
\begin{tikzpicture}[scale=1]
% Colonne oxydants
\node[font=\bfseries\large, popblue] at (-3.2,6.6) {Oxydants};
\draw[-{Stealth[length=3mm]}, line width=1.5pt, popblue] (-5.6,0.4) -- (-5.6,6.2);
\node[rotate=90, popblue, font=\small] at (-6.15,3.3) {pouvoir oxydant croissant};
% Colonne réducteurs
\node[font=\bfseries\large, poporange] at (3.2,6.6) {Réducteurs};
\draw[-{Stealth[length=3mm]}, line width=1.5pt, poporange] (5.6,6.2) -- (5.6,0.4);
\node[rotate=-90, poporange, font=\small] at (6.15,3.3) {pouvoir réducteur croissant};
% Couples
\foreach \y/\ox/\red in {5.6/{\ce{Ag+}}/{\ce{Ag}}, 4.6/{\ce{Cu^2+}}/{\ce{Cu}}, 3.4/{\ce{H3O+}}/{\ce{H2}}, 2.2/{\ce{Fe^2+}}/{\ce{Fe}}, 1.2/{\ce{Zn^2+}}/{\ce{Zn}}, 0.2/{\ce{Al^3+}}/{\ce{Al}}}{
  \node[popblue, font=\large] at (-3.2,\y) {\ox};
  \node[font=\large] at (0,\y) {/};
  \node[poporange, font=\large] at (3.2,\y) {\red};
  \draw[popline, dashed] (-2.2,\y) -- (2.2,\y);
}
% Mise en évidence du couple H3O+/H2
\draw[rounded corners=4pt, line width=1.6pt, popgreen] (-4.6,2.85) rectangle (4.6,3.95);
\node[popgreen, font=\small\bfseries, anchor=west] at (4.75,3.4) {couple de référence};
% Métaux nobles
\draw[rounded corners=4pt, line width=1.4pt, popgold] (-4.6,4.15) rectangle (4.6,6.05);
\node[popgold, font=\small\bfseries, anchor=east] at (-4.75,5.1) {métaux nobles};
\end{tikzpicture}
\end{center}
\legende{Classification électrochimique qualitative : les oxydants sont classés par force croissante vers le haut, les réducteurs par force croissante vers le bas. Le couple \ce{H3O+}/\ce{H2} (encadré vert) sépare les métaux attaqués par les acides des métaux nobles (encadré doré).}
\end{popfigure}

\courssub{La place capitale du couple \ce{H3O+}/\ce{H2}}

Le couple \ce{H3O+}/\ce{H2} joue le rôle de \cle{frontière} dans la classification, un peu comme la ligne de milieu de terrain sur un terrain de football :

\begin{itemize}
\item \textbf{Au-dessus de \ce{H2}} (du côté des réducteurs forts) : Zn, Fe, Al, Mg... Ces métaux sont des réducteurs plus forts que \ce{H2}, donc ils \textbf{réduisent les ions \ce{H3O+}} : ils sont attaqués par les acides non oxydants (HCl dilué, \ce{H2SO4} dilué) avec dégagement de dihydrogène.
\item \textbf{En dessous de \ce{H2}} : Cu, Ag, Au, Pt... Ce sont les \cle{métaux nobles}. Leur pouvoir réducteur est plus faible que celui de \ce{H2} : ils \textbf{ne réduisent pas \ce{H3O+}} et restent inattaqués par les acides non oxydants.
\end{itemize}

\begin{popfigure}
\begin{center}
\begin{tikzpicture}[scale=1]
\draw[-{Stealth[length=3mm]}, line width=1.5pt, popdark] (0,0) -- (0,10.4);
\node[rotate=90, font=\small, popdark] at (-0.45,5.2) {pouvoir oxydant de l'oxydant croissant};
% Graduations et couples
\foreach \y/\c/\col in {9.2/{\ce{Ag+}/\ce{Ag}}/popgold, 7.4/{\ce{Cu^2+}/\ce{Cu}}/popgold, 5.2/{\ce{H3O+}/\ce{H2}}/popgreen, 3.0/{\ce{Fe^2+}/\ce{Fe}}/popblue, 1.0/{\ce{Zn^2+}/\ce{Zn}}/popblue}{
  \draw[line width=1.2pt, \col] (-0.25,\y) -- (0.25,\y);
  \node[anchor=west, font=\large, \col] at (0.5,\y) {\c};
}
% Annotations
\draw[decorate, decoration={brace, amplitude=6pt}, popblue, line width=1.2pt] (3.6,0.5) -- (3.6,4.6);
\node[anchor=west, popblue, font=\small, align=left] at (4.0,2.55) {métaux attaqués\\par les acides\\(dégagement de \ce{H2})};
\draw[decorate, decoration={brace, amplitude=6pt}, popgold, line width=1.2pt] (3.6,6.3) -- (3.6,9.7);
\node[anchor=west, popgold, font=\small, align=left] at (4.0,8.0) {métaux nobles :\\non attaqués\\par les acides};
\node[popgreen, font=\small\bfseries, align=center] at (-2.6,5.2) {frontière\\\ce{H3O+}/\ce{H2}};
\draw[-{Stealth}, popgreen] (-1.5,5.2) -- (-0.35,5.2);
\end{tikzpicture}
\end{center}
\legende{Frise verticale de la classification : le couple \ce{H3O+}/\ce{H2} sépare les métaux attaqués par les acides dilués (Fe, Zn) des métaux nobles (Cu, Ag).}
\end{popfigure}

\begin{exemplebox}[Trois expériences suffisent à classer quatre couples]
À partir des trois seules expériences Fe/\ce{Cu^2+}, Cu/\ce{Ag+} et Zn/\ce{H3O+}, complétées par l'absence de réaction entre Cu et \ce{H3O+}, établissons l'ordre des couples.
\begin{itemize}
\item Fe réduit \ce{Cu^2+} $\Rightarrow$ Fe $>$ Cu (réducteurs).
\item Cu réduit \ce{Ag+} $\Rightarrow$ Cu $>$ Ag.
\item Cu ne réduit pas \ce{H3O+} $\Rightarrow$ \ce{H2} $>$ Cu.
\item Zn réduit \ce{H3O+} $\Rightarrow$ Zn $>$ \ce{H2}.
\end{itemize}
En chaînant : Ag $<$ Cu $<$ \ce{H2} et Zn $>$ \ce{H2}, Fe $>$ Cu. L'expérience Fe/acide précise Fe $>$ \ce{H2}, et Zn/\ce{Fe^2+} donne Zn $>$ Fe. D'où le classement complet : Ag $<$ Cu $<$ \ce{H2} $<$ Fe $<$ Zn.
\end{exemplebox}

\begin{exoresolu}[Classer trois couples à partir d'observations expérimentales]
\textbf{Énoncé.} On réalise trois expériences à \SI{25}{\celsius} :
\begin{enumerate}
\item Un fil de cuivre plongé dans une solution de nitrate d'argent se recouvre d'un dépôt gris et la solution bleuit.
\item Une lame de zinc plongée dans une solution de sulfate de cuivre se recouvre d'un dépôt rouge.
\item Un fil de cuivre plongé dans une solution de chlorure de zinc ne subit aucune transformation.
\end{enumerate}
Écrire les équations des réactions et établir le classement des trois couples \ce{Ag+}/\ce{Ag}, \ce{Cu^2+}/\ce{Cu}, \ce{Zn^2+}/\ce{Zn} par pouvoir oxydant croissant.
\tcblower
\textbf{Solution détaillée.}
\begin{enumerate}
\item Le dépôt gris est de l'argent et la teinte bleue révèle les ions \ce{Cu^2+} :
\[ \ce{2Ag+ + Cu -> 2Ag + Cu^2+} \]
\ce{Ag+} est donc un oxydant plus fort que \ce{Cu^2+} (et Cu un réducteur plus fort que Ag).
\item Le dépôt rouge est du cuivre :
\[ \ce{Cu^2+ + Zn -> Cu + Zn^2+} \]
\ce{Cu^2+} est un oxydant plus fort que \ce{Zn^2+} (et Zn un réducteur plus fort que Cu).
\item Aucune réaction entre Cu et \ce{Zn^2+} : cela confirme que Cu est un réducteur \emph{plus faible} que Zn, donc que \ce{Zn^2+} est un oxydant plus faible que \ce{Cu^2+}. Cohérent avec l'expérience 2 !
\end{enumerate}
\textbf{Classement par pouvoir oxydant croissant :}
\[ \ce{Zn^2+} < \ce{Cu^2+} < \ce{Ag+} \]
et, en sens inverse, par pouvoir réducteur croissant : $\ce{Ag} < \ce{Cu} < \ce{Zn}$.
\end{exoresolu}

\begin{attention}[Le piège classique : inverser les flèches]
L'erreur la plus fréquente au Probatoire est d'inverser le sens des flèches de force croissante. Retiens bien :
\begin{itemize}
\item côté \textbf{oxydants} : la force croît \textbf{vers le haut} (\ce{Ag+} en haut est le plus fort oxydant) ;
\item côté \textbf{réducteurs} : la force croît \textbf{vers le bas} (Zn, Al, Mg en bas sont les plus forts réducteurs).
\end{itemize}
Astuce de mémorisation : « l'oxydant monte, le réducteur descend ». Et n'oublie jamais qu'un couple placé haut côté oxydant est forcément placé haut côté réducteur \emph{faible} : les deux membres d'un couple sont toujours sur la \textbf{même ligne horizontale} du tableau.
\end{attention}

\begin{savaistu}[L'or et le platine, champions de l'inertie]
Tout en bas de l'échelle des réducteurs se trouvent l'or (Au) et le platine (Pt). Leur pouvoir réducteur est si faible qu'ils ne s'oxydent \emph{jamais} à l'air, même chauffés, et résistent à presque tous les acides. Seule l'« eau régale » (mélange concentré d'acide chlorhydrique et d'acide nitrique) parvient à dissoudre l'or ! C'est précisément cette inertie chimique qui explique leur usage en joaillerie : une bague en or traverse les générations sans ternir, contrairement à un bijou en fer qui rouille en quelques saisons des pluies. À Yaoundé comme ailleurs, quand on dit qu'une amitié est « en or », on dit donc, en langage de chimiste, qu'elle est tout en bas de la classification électrochimique : inoxydable !
\end{savaistu}

\begin{aretenir}[Ce qu'il faut retenir de cette section]
\begin{itemize}
\item La classification électrochimique range les couples selon le pouvoir oxydant croissant de leurs oxydants ; chaque réaction spontanée observée impose une inégalité.
\item Classement des réducteurs (croissant) : $\ce{Ag} < \ce{Cu} < \ce{H2} < \ce{Fe} < \ce{Zn}$ (puis Al, Mg).
\item Classement des oxydants (croissant) : $\ce{Zn^2+} < \ce{Fe^2+} < \ce{H3O+} < \ce{Cu^2+} < \ce{Ag+}$.
\item Le couple \ce{H3O+}/\ce{H2} est la frontière : au-dessus de \ce{H2}, les métaux sont attaqués par les acides dilués ; en dessous se trouvent les métaux nobles (Cu, Ag, Au, Pt).
\item Une absence de réaction est un verdict expérimental aussi précieux qu'une réaction observée.
\end{itemize}
\end{aretenir}

\coursec{Utilisation de la classification : prévision des réactions}

Dans la section précédente, nous avons construit pas à pas la \cle{classification électrochimique qualitative} : à partir d'expériences simples (une lame de zinc dans une solution de sulfate de cuivre(II), un clou en fer dans l'acide chlorhydrique), nous avons rangé les couples oxydant-réducteur les uns par rapport aux autres, en plaçant le couple \ce{H3O+}/\ce{H2} comme repère central. Mais une classification n'est pas une fin en soi : c'est un \textbf{outil de prévision}. C'est précisément ce que nous allons apprendre ici : à partir du seul tableau de classification, dire si une réaction entre deux espèces est possible, et écrire son équation-bilan. C'est une compétence \emph{incontournable} au Probatoire.

\courssub{La règle fondamentale}

\textbf{Intuition : qui cède, qui prend ?}\par

Une réaction d'oxydoréduction est un \textbf{transfert d'électrons} : le réducteur \emph{donne} des électrons, l'oxydant les \emph{prend}. Pour que la réaction ait lieu spontanément, il faut un donneur généreux et un receveur avide. Or la classification range précisément les espèces selon cette générosité : plus un réducteur est placé \emph{haut} dans la classification, plus il donne facilement ses électrons ; plus un oxydant est placé \emph{bas}, plus il les capte facilement.

\begin{propriete}[Règle fondamentale]
Une réaction d'oxydoréduction \cle{spontanée} a lieu entre l'\textbf{oxydant le plus fort} et le \textbf{réducteur le plus fort} présents dans le milieu. Elle produit l'oxydant le plus faible et le réducteur le plus faible. Autrement dit : la réaction naturelle va toujours du couple le plus « haut placé » (pour le réducteur) vers le couple le plus « bas placé » (pour l'oxydant).
\end{propriete}

\begin{attention}[Deux oxydants ou deux réducteurs : rien ne se passe]
L'erreur la plus fréquente au Probatoire est de vouloir faire réagir \textbf{deux oxydants entre eux} (par exemple \ce{Cu^2+} avec \ce{Fe^3+}) ou \textbf{deux réducteurs entre eux} (par exemple le zinc métallique avec le fer métallique). C'est \emph{impossible} : un oxydant ne peut que capter des électrons, un réducteur ne peut que les céder. Deux capteurs ou deux donneurs face à face : aucun transfert d'électrons n'est possible, donc aucune réaction. Avant tout calcul, vérifie que ton mélange contient bien \textbf{un oxydant et un réducteur}.
\end{attention}

\courssub{La règle du « gamma »}

Il existe une image très pratique pour appliquer la règle fondamentale sans se tromper : la \cle{règle du gamma} (du nom de la lettre grecque $\gamma$, dont la forme rappelle la diagonale).

\begin{propriete}[Règle du gamma]
On écrit les couples en colonne, l'oxydant le plus fort en bas, le réducteur le plus fort en haut. La réaction spontanée possible se produit entre les espèces situées en \textbf{diagonale} : l'oxydant du couple inférieur réagit avec le réducteur du couple supérieur. Si les deux espèces mises en présence ne sont pas en diagonale (oxydant du couple supérieur face au réducteur du couple inférieur), il n'y a \textbf{pas de réaction}.
\end{propriete}

\begin{popfigure}
\begin{center}
\begin{tikzpicture}[scale=1.0]
% Colonne oxydants / réducteurs
\node[font=\bfseries, popblue] at (-1.6,4.6) {Oxydants};
\node[font=\bfseries, poporange] at (1.6,4.6) {Réducteurs};
% Couples
\node[font=\large] at (-1.6,3.6) {\ce{Cu^2+}};
\node[font=\large] at (1.6,3.6) {\ce{Cu}};
\node[font=\large] at (-1.6,2.6) {\ce{H3O+}};
\node[font=\large] at (1.6,2.6) {\ce{H2}};
\node[font=\large] at (-1.6,1.6) {\ce{Fe^2+}};
\node[font=\large] at (1.6,1.6) {\ce{Fe}};
\node[font=\large] at (-1.6,0.6) {\ce{Zn^2+}};
\node[font=\large] at (1.6,0.6) {\ce{Zn}};
% Flèches de force
\draw[-{Stealth[length=3mm]}, very thick, popblue] (-2.9,0.4) -- (-2.9,3.8) node[midway, left, font=\small, align=center]{oxydant de\\plus en plus\\fort};
\draw[-{Stealth[length=3mm]}, very thick, poporange] (2.9,3.8) -- (2.9,0.4) node[midway, right, font=\small, align=center]{réducteur de\\plus en plus\\fort};
% Flèche gamma : Zn (réducteur fort) -> Cu2+ (oxydant fort)
\draw[-{Stealth[length=4mm]}, line width=2pt, popgreen] (1.35,0.55) -- (-1.35,3.5);
\node[popgreen, font=\small\bfseries, align=center] at (0,2.1) {réaction\\spontanée};
\end{tikzpicture}
\end{center}
\legende{Règle du gamma : le réducteur le plus fort (\ce{Zn}, en haut à droite) réagit spontanément avec l'oxydant le plus fort (\ce{Cu^2+}, en bas à gauche) : la flèche diagonale dessine un $\gamma$. En revanche, \ce{Cu} et \ce{Zn^2+} ne réagissent pas entre eux.}
\end{popfigure}

\begin{methode}[Prévoir une réaction d'oxydoréduction en quatre étapes]
\begin{enumerate}
\item \textbf{Repérer les espèces présentes} dans le milieu (métal, ions, solution acide) et identifier \textbf{les deux couples} mis en jeu.
\item \textbf{Placer les deux couples} dans la classification électrochimique (du réducteur le plus fort en haut à l'oxydant le plus fort en bas).
\item \textbf{Appliquer la règle du gamma} : si l'oxydant et le réducteur présents sont en diagonale (oxydant du couple du bas, réducteur du couple du haut), la réaction est possible ; sinon, il ne se passe rien.
\item \textbf{Écrire l'équation-bilan} en combinant les deux demi-équations électroniques de façon à éliminer les électrons, puis vérifier la conservation des atomes et des charges.
\end{enumerate}
\end{methode}

\courssub{Application 1 : quel récipient pour quelle solution ?}

Voici une question typique de Probatoire, et aussi un vrai problème de laboratoire : peut-on stocker une solution dans un récipient métallique sans que le récipient soit attaqué ?

\begin{exemplebox}[Un récipient en cuivre peut-il contenir une solution de \ce{Fe^2+} ? de \ce{Ag+} ?]
\textbf{Cas de la solution de \ce{Fe^2+}.} Les espèces présentes sont \ce{Fe^2+} (oxydant du couple \ce{Fe^2+}/\ce{Fe}) et \ce{Cu} (réducteur du couple \ce{Cu^2+}/\ce{Cu}). Dans la classification, \ce{Cu} est placé \emph{au-dessous} de \ce{Fe} (le cuivre est un réducteur plus faible que le fer) : les deux espèces ne sont \textbf{pas en diagonale}. Il n'y a \textbf{pas de réaction} : le récipient en cuivre peut contenir la solution de \ce{Fe^2+} sans être attaqué.

\textbf{Cas de la solution de \ce{Ag+}.} Les espèces présentes sont \ce{Ag+} (oxydant du couple \ce{Ag+}/\ce{Ag}, placé tout en bas de la classification, donc oxydant fort) et \ce{Cu} (réducteur). Cette fois, les deux espèces \textbf{sont en diagonale} : la réaction est spontanée. Le cuivre s'oxyde et les ions argent se réduisent :
\begin{align*}
\text{oxydation : } & \ce{Cu -> Cu^2+ + 2e-}\\
\text{réduction : } & \ce{Ag+ + e- -> Ag} \quad (\times 2)
\end{align*}
\[ \ce{Cu + 2Ag+ -> Cu^2+ + 2Ag} \]
Le récipient se dissout lentement et se recouvre d'un dépôt gris d'argent métallique : la solution est perdue et le récipient détruit. Conclusion : \textbf{jamais de solution de \ce{Ag+} dans un récipient en cuivre}.
\end{exemplebox}

\courssub{Application 2 : choix du métal pour stocker une solution acide}

Une solution acide contient des ions \ce{H3O+}, oxydant du couple \ce{H3O+}/\ce{H2}. D'après la règle du gamma, tout métal placé \textbf{au-dessus} de \ce{H2} dans la classification (zinc, fer, aluminium, magnésium\ldots) est un réducteur assez fort pour être attaqué par l'acide, avec dégagement de dihydrogène :
\[ \ce{M + n H3O+ -> M^{n+} + n/2 H2 + n H2O} \]
En revanche, les métaux placés \textbf{au-dessous} de \ce{H2} (cuivre, argent, or) ne sont pas attaqués par les acides non oxydants comme \ce{HCl} dilué.

\begin{attention}[Conséquence pratique]
On ne conserve jamais une solution acide dans un bidon en fer ou en zinc : le récipient serait rongé et le dégagement de dihydrogène (gaz inflammable !) créerait un risque d'explosion. C'est pourquoi les acides du laboratoire sont stockés dans des flacons en verre ou en matière plastique, chimiquement inertes.
\end{attention}

\courssub{Application 3 : corrosion du fer et protection par le zinc}

Le fer est partout autour de nous : tôles de nos maisons, carrosseries, pylônes, charpentes. Or le fer, réducteur assez fort, s'oxyde facilement au contact de l'air humide : c'est la \cle{corrosion}, qui produit la rouille. Comment protéger le fer ? L'idée, géniale, consiste à lui associer un métal \textbf{réducteur plus fort que lui}, qui s'oxydera \emph{à sa place}.

\begin{propriete}[Protection par un métal plus réducteur]
Le zinc est placé au-dessus du fer dans la classification : c'est un réducteur plus fort. Si du zinc est en contact avec du fer exposé à l'humidité, c'est le zinc qui s'oxyde (\ce{Zn -> Zn^2+ + 2e-}) et non le fer. Le zinc « se sacrifie » pour protéger le fer : on parle de \cle{protection cathodique} ou d'\textbf{anode sacrificielle}.
\end{propriete}

C'est le principe de la \cle{galvanisation} : on recouvre les tôles de fer d'une fine couche de zinc (tôles « galvanisées » très utilisées pour les toitures au Cameroun). Même si la couche de zinc est rayée, le fer reste protégé tant qu'il reste du zinc à proximité. Le \textbf{fer blanc} (boîtes de conserve) est, lui, recouvert d'étain : protection moins durable, car l'étain est un réducteur plus faible que le fer — si le revêtement est rayé, le fer rouille alors \emph{plus vite}.

\begin{popfigure}
\begin{center}
\begin{tikzpicture}[scale=0.95]
% ===== Fer galvanisé =====
\begin{scope}
\fill[popgreen!25] (0,0) rectangle (5,0.7);
\fill[popblue!45] (0,0.7) rectangle (5,1.0);
\node[font=\small] at (2.5,0.35) {fer};
\node[font=\small, popblue!80!black] at (2.5,0.85) {couche de zinc};
% gouttes d'eau
\foreach \x in {0.8,2.0,3.2,4.4}{
  \draw[fill=popblue!50, draw=popblue] (\x,1.55) .. controls (\x-0.13,1.3) and (\x-0.13,1.15) .. (\x,1.12) .. controls (\x+0.13,1.15) and (\x+0.13,1.3) .. (\x,1.55);
}
\node[font=\small, popblue] at (2.5,2.0) {air humide};
% rayure
\fill[popgreen!25] (2.3,0.7) rectangle (2.7,1.0);
\draw[-{Stealth[length=2.5mm]}, thick, popgreen!60!black] (3.4,0.35) -- (2.55,0.85);
\node[font=\scriptsize, popgreen!60!black, align=center] at (4.1,0.3) {même rayé,\\le zinc s'oxyde\\à la place du fer};
\node[font=\bfseries, popgreen!60!black] at (2.5,-0.5) {Fer galvanisé : protégé};
\end{scope}
% ===== Fer non protégé =====
\begin{scope}[xshift=7.5cm]
\fill[popgreen!25] (0,0) rectangle (5,1.0);
\node[font=\small] at (2.5,0.5) {fer};
% rouille
\fill[poporange!70, decorate, decoration={random steps, segment length=3pt, amplitude=1.5pt}] (0.6,1.0) rectangle (2.2,1.35);
\fill[poporange!70, decorate, decoration={random steps, segment length=3pt, amplitude=1.5pt}] (3.0,1.0) rectangle (4.5,1.3);
\node[font=\scriptsize, poporange!80!black] at (1.4,1.6) {rouille};
\node[font=\scriptsize, poporange!80!black] at (3.75,1.55) {rouille};
% gouttes d'eau
\foreach \x in {0.8,2.0,3.2,4.4}{
  \draw[fill=popblue!50, draw=popblue] (\x,2.15) .. controls (\x-0.13,1.9) and (\x-0.13,1.75) .. (\x,1.72) .. controls (\x+0.13,1.75) and (\x+0.13,1.9) .. (\x,2.15);
}
\node[font=\small, popblue] at (2.5,2.6) {air humide};
\node[font=\bfseries, poporange!80!black] at (2.5,-0.5) {Fer nu : il rouille};
\end{scope}
\end{tikzpicture}
\end{center}
\legende{À gauche, le fer galvanisé : le zinc, réducteur plus fort que le fer, s'oxyde à sa place, même en cas de rayure. À droite, le fer non protégé s'oxyde directement au contact de l'air humide : il se couvre de rouille.}
\end{popfigure}

\begin{savaistu}[Les anodes sacrificielles]
Les coques des navires, les pipelines enterrés et les chauffe-eau sont protégés par des blocs de zinc ou de magnésium fixés sur la structure métallique : ce sont des \textbf{anodes sacrificielles}. Le magnésium et le zinc, réducteurs plus forts que le fer, s'oxydent à sa place et sont « rongés » à la place de la coque. On les remplace périodiquement. Sans elles, un navire rouillerait en quelques années ! Le même principe protège les pipelines qui transportent le pétrole, par exemple ceux de la SONARA ou du corridor Tchad--Cameroun.
\end{savaistu}

\courssub{Exemple chiffré complet : zinc et acide chlorhydrique}

\begin{exoresolu}[Réaction du zinc avec l'acide chlorhydrique]
On plonge une masse $m = \SI{6,5}{g}$ de zinc en grenaille dans un volume $V = \SI{200}{mL}$ d'acide chlorhydrique de concentration $C = \SI{1,0}{mol.L^{-1}}$. On donne $M(\ce{Zn}) = \SI{65}{g.mol^{-1}}$ et le volume molaire $V_m = \SI{24}{L.mol^{-1}}$ dans les conditions de l'expérience.\\
1. La réaction est-elle possible ? Écrire son équation-bilan.\\
2. Déterminer le réactif limitant.\\
3. Calculer le volume de dihydrogène dégagé.
\tcblower
\textbf{1. Possibilité de la réaction.} Le zinc est placé au-dessus du couple \ce{H3O+}/\ce{H2} dans la classification : c'est un réducteur plus fort que \ce{H2}. Les espèces \ce{Zn} et \ce{H3O+} sont en diagonale (règle du gamma) : la réaction est \textbf{spontanée}. Demi-équations :
\begin{align*}
\text{oxydation : } & \ce{Zn -> Zn^2+ + 2e-}\\
\text{réduction : } & \ce{2H3O+ + 2e- -> H2 + 2H2O}
\end{align*}
\[ \ce{Zn + 2H3O+ -> Zn^2+ + H2 + 2H2O} \]
On observe une effervescence (dégagement de \ce{H2}) et la dissolution progressive du zinc.

\textbf{2. Réactif limitant.} Quantités initiales :
\[ n(\ce{Zn}) = \frac{m}{M} = \frac{6,5}{65} = \SI{0,10}{mol} \qquad n(\ce{H3O+}) = C \times V = 1,0 \times 0,200 = \SI{0,20}{mol} \]
Tableau d'avancement :
\begin{center}
\begin{tabularx}{\linewidth}{|l|X|X|X|X|}
\hline
\rowcolor{popblueL}
État & \ce{Zn} & \ce{2H3O+} & \ce{Zn^2+} & \ce{H2} \\
\hline
Initial (mol) & $0,10$ & $0,20$ & $0$ & $0$ \\
\hline
Final (mol) & $0,10 - x_{max}$ & $0,20 - 2x_{max}$ & $x_{max}$ & $x_{max}$ \\
\hline
\end{tabularx}
\end{center}
Hypothèses : si le zinc est limitant, $x_{max} = 0,10$ mol ; si \ce{H3O+} est limitant, $x_{max} = \dfrac{0,20}{2} = 0,10$ mol. Les deux valeurs sont égales : les réactifs sont en \textbf{proportions stœchiométriques}, $x_{max} = \SI{0,10}{mol}$. Les deux réactifs disparaissent entièrement.

\textbf{3. Volume de dihydrogène.}
\[ n(\ce{H2}) = x_{max} = \SI{0,10}{mol} \quad\Longrightarrow\quad V(\ce{H2}) = n \times V_m = 0,10 \times 24 = \SI{2,4}{L} \]
Le dégagement de dihydrogène est de $\boxed{\SI{2,4}{L}}$.
\end{exoresolu}

\courssub{Limites de la classification qualitative}

La classification électrochimique qualitative est un outil puissant, mais il faut connaître ses limites pour ne pas lui faire dire ce qu'elle ne dit pas.

\begin{attention}[Ce que la classification ne prévoit pas]
\begin{itemize}
\item Elle prévoit le \textbf{sens} de la réaction spontanée, mais \textbf{pas sa vitesse} : certaines réactions prévues sont infiniment lentes à température ordinaire (l'aluminium, pourtant réducteur très fort, résiste à l'air grâce à une couche protectrice d'alumine : c'est la \emph{passivation}).
\item Elle ne dit rien du \textbf{rendement} ni des conditions expérimentales (température, concentration, catalyseurs).
\item Elle est \textbf{qualitative} : elle classe les couples sans mesurer l'écart de force entre eux.
\end{itemize}
\end{attention}

\begin{savaistu}[Complément Probatoire : vers la classification quantitative]
En Terminale, tu rendras cette classification \textbf{quantitative} : à chaque couple on associera un nombre, le \textbf{potentiel standard d'électrode} $E^\circ$ (en volts), mesuré par rapport au couple \ce{H3O+}/\ce{H2} pris comme référence ($E^\circ = \SI{0,00}{V}$). On pourra alors non seulement prévoir le sens des réactions, mais aussi calculer des constantes d'équilibre et comprendre le fonctionnement des piles. En Première, restons qualitatifs : la règle du gamma suffit à résoudre tous les exercices du programme.
\end{savaistu}

\begin{aretenir}[L'essentiel de la section]
\begin{itemize}
\item Une réaction spontanée a lieu entre l'\textbf{oxydant le plus fort} et le \textbf{réducteur le plus fort} présents en solution.
\item \textbf{Règle du gamma} : la réaction se produit entre les espèces placées en diagonale dans la classification ; sinon, il n'y a pas de réaction.
\item Deux oxydants (ou deux réducteurs) ne réagissent jamais entre eux.
\item Applications : choix des récipients (pas de \ce{Ag+} dans le cuivre, pas d'acide dans le fer), protection du fer par le zinc (galvanisation, anodes sacrificielles).
\item La classification prévoit le sens de la réaction, mais ni sa vitesse ni son rendement.
\end{itemize}
\end{aretenir}

\coursec{Synthèse et applications au quotidien camerounais}

Dans la section précédente, nous avons appris à \emph{utiliser} la classification électrochimique pour prévoir si une réaction d'oxydoréduction peut avoir lieu entre deux couples. Il est temps maintenant de rassembler toutes les pièces du puzzle : qu'est-ce qu'un couple oxydant-réducteur ? Comment lire une classification ? Et surtout, à quoi tout cela sert-il concrètement, du bidon de fongicide du cultivateur de Foumbot aux tôles des cases de Kribi ? Cette section de synthèse est aussi ta fiche de révision idéale pour le Probatoire.

\courssub{Bilan des savoirs essentiels}

\courssub{Le couple oxydant-réducteur}

\begin{definition}[Couple oxydant-réducteur]
Un \cle{couple oxydant-réducteur}, noté \cle{Ox/Red}, est constitué de deux espèces chimiques, un \cle{oxydant} Ox et un \cle{réducteur} Red, qui se transforment l'un en l'autre par gain ou perte d'électrons. On l'associe à une \cle{demi-équation électronique} :
\[ \mathrm{Ox} + n\,e^- \rightleftharpoons \mathrm{Red} \]
L'oxydant est l'espèce qui \textbf{capte} des électrons (il se réduit) ; le réducteur est l'espèce qui \textbf{cède} des électrons (il s'oxyde).
\end{definition}

Retiens l'ordre d'écriture : on écrit toujours le couple sous la forme \textbf{Ox/Red}, l'oxydant en premier. Ainsi on écrit \ce{Cu^2+}/\ce{Cu} et jamais \ce{Cu}/\ce{Cu^2+}. Quelques couples rencontrés dans cette leçon :

\begin{center}
\begin{tabularx}{\linewidth}{|l|X|}
\hline
\rowcolor{popblueL} \textbf{Couple Ox/Red} & \textbf{Demi-équation électronique} \\
\hline
\ce{Cu^2+}/\ce{Cu} & \ce{Cu^2+ + 2e- <=> Cu} \\
\hline
\ce{Zn^2+}/\ce{Zn} & \ce{Zn^2+ + 2e- <=> Zn} \\
\hline
\ce{Fe^2+}/\ce{Fe} & \ce{Fe^2+ + 2e- <=> Fe} \\
\hline
\ce{Ag+}/\ce{Ag} & \ce{Ag+ + e- <=> Ag} \\
\hline
\ce{H3O+}/\ce{H2} & \ce{2H3O+ + 2e- <=> H2 + 2H2O} \\
\hline
\end{tabularx}
\end{center}

\courssub{Réaction spontanée entre deux couples}

\begin{propriete}[Règle de la réaction spontanée]
Une réaction d'oxydoréduction \cle{spontanée} a lieu entre l'oxydant d'un couple et le réducteur d'un \textbf{autre} couple : c'est toujours l'oxydant le \textbf{plus fort} qui réagit avec le réducteur le \textbf{plus fort}. Dans l'équation bilan, les électrons n'apparaissent jamais : on multiplie chaque demi-équation par un coefficient tel que le nombre d'électrons cédés égale le nombre d'électrons captés, puis on additionne.
\end{propriete}

Par exemple, entre les couples \ce{Cu^2+}/\ce{Cu} et \ce{Fe^2+}/\ce{Fe} :
\begin{align*}
\text{Réduction : } & \ce{Cu^2+ + 2e- -> Cu} \\
\text{Oxydation : } & \ce{Fe -> Fe^2+ + 2e-} \\
\text{Bilan : } & \ce{Cu^2+ + Fe -> Cu + Fe^2+}
\end{align*}
Les deux électrons cédés par le fer sont captés par l'ion cuivre : la comptabilité électronique est équilibrée.

\courssub{Classification électrochimique qualitative et place de \ce{H3O+}/\ce{H2}}

\begin{aretenir}[Classification qualitative]
La \cle{classification électrochimique qualitative} range les couples Ox/Red par \textbf{pouvoir oxydant croissant} (ou pouvoir réducteur décroissant). Elle se construit expérimentalement : si le métal A déplace le métal B de ses ions (A s'oxyde, les ions B se réduisent), alors le réducteur A est plus fort que B, et le couple B$^{n+}$/B est placé au-dessus du couple A$^{m+}$/A. Le couple \ce{H3O+}/\ce{H2} sert de \textbf{frontière} :
\begin{itemize}
\item les métaux situés \textbf{au-dessous} de \ce{H2} (Fe, Zn, Al, Pb\ldots) sont attaqués par les ions \ce{H3O+} des acides avec dégagement de dihydrogène ;
\item les métaux situés \textbf{au-dessus} de \ce{H2} (Cu, Ag, Au\ldots) ne réagissent pas avec les acides non oxydants.
\end{itemize}
Ordre obtenu expérimentalement (pouvoir oxydant croissant) :
\[ \ce{Zn^2+}/\ce{Zn} \;<\; \ce{Fe^2+}/\ce{Fe} \;<\; \ce{Pb^2+}/\ce{Pb} \;<\; \ce{H3O+}/\ce{H2} \;<\; \ce{Cu^2+}/\ce{Cu} \;<\; \ce{Ag+}/\ce{Ag} \]
\end{aretenir}

La carte mentale ci-dessous résume l'architecture de toute la leçon.

\begin{popfigure}
\begin{center}
\begin{tikzpicture}[
  font=\small,
  central/.style={ellipse, draw=popdark, fill=popblue, text=white, thick, align=center, minimum width=4.2cm, minimum height=1.4cm},
  branche/.style={rectangle, draw=popdark, thick, rounded corners=3pt, align=center, minimum width=3.6cm, minimum height=1.1cm},
  lien/.style={thick, popmuted}
]
\node[central] (c) at (0,0) {\textbf{Couple Ox/Red}\\ $\mathrm{Ox} + n\,e^- \rightleftharpoons \mathrm{Red}$};
\node[branche, fill=poporangeL] (de) at (-5.2,2.6) {\textbf{Demi-équation}\\ gain/perte d'électrons\\ (ex. \ce{Cu^2+ + 2e- <=> Cu})};
\node[branche, fill=popgreenL] (cl) at (5.2,2.6) {\textbf{Classification}\\ pouvoir oxydant croissant\\ frontière \ce{H3O+}/\ce{H2}};
\node[branche, fill=poppinkL] (pr) at (-5.2,-2.6) {\textbf{Prévision}\\ Ox fort $+$ Red fort\\ $\rightarrow$ réaction spontanée};
\node[branche, fill=popgoldL] (ap) at (5.2,-2.6) {\textbf{Applications}\\ corrosion, cimentation,\\ protection des métaux};
\draw[lien] (c) -- (de);
\draw[lien] (c) -- (cl);
\draw[lien] (c) -- (pr);
\draw[lien] (c) -- (ap);
\end{tikzpicture}
\end{center}
\legende{Carte mentale de synthèse : autour du couple oxydant-réducteur s'articulent la demi-équation électronique, la classification, la prévision des réactions et les applications.}
\end{popfigure}

\courssub{Démarche type d'un exercice au Probatoire}

Les exercices de classification électrochimique suivent presque toujours le même schéma : on te décrit des expériences (dépôt, dégagement gazeux, changement de couleur) et on te demande de classer des couples. Voici la méthode infaillible.

\begin{methode}[Résolution d'un exercice de classification en 5 étapes]
\begin{enumerate}
\item \textbf{Analyser les observations.} Traduire chaque observation en langage chimique : un dépôt rouge sur le fer = du cuivre métallique formé ; un dégagement qui détone à la flamme = \ce{H2} ; une solution qui se décolore = disparition d'un ion.
\item \textbf{Identifier les produits et les réactifs.} Déterminer quelle espèce a disparu (réactif) et quelle espèce est apparue (produit) pour chaque élément concerné.
\item \textbf{Écrire les demi-équations.} Une oxydation (perte d'électrons, le métal devient ion) et une réduction (gain d'électrons, l'ion devient métal ou \ce{H3O+} devient \ce{H2}).
\item \textbf{Écrire l'équation bilan.} Multiplier pour équilibrer les électrons, additionner, vérifier la conservation des atomes et des charges.
\item \textbf{Conclure sur la classification.} Si la réaction a lieu, le réducteur qui s'oxyde est plus fort : son couple est placé plus bas. Si rien ne se passe, conclure à l'impossibilité de la réaction spontanée dans ce sens (et donc à l'ordre inverse).
\end{enumerate}
\end{methode}

\begin{attention}[Erreur fréquente : conclure trop vite]
Ne conclus \textbf{jamais} sur la position relative de deux couples à partir d'une seule expérience sans vérifier le \textbf{sens} de la réaction observée. Si tu observes un dépôt de cuivre sur une lame de fer plongée dans \ce{CuSO4}, c'est bien \ce{Fe} qui s'oxyde et \ce{Cu^2+} qui se réduit : le couple \ce{Cu^2+}/\ce{Cu} est donc \emph{au-dessus} de \ce{Fe^2+}/\ce{Fe}. Inverser la conclusion (placer le fer au-dessus parce qu'il « réagit ») est la faute la plus fréquente au Probatoire. Autre piège : l'absence de réaction est aussi une information précieuse — elle signifie que la réaction spontanée aurait lieu dans le \emph{sens inverse}. Vérifie toujours la cohérence avec au moins deux expériences croisées.
\end{attention}

\begin{exoresolu}[Classement de trois couples à partir d'expériences]
On réalise trois expériences à \SI{25}{\celsius} :
\begin{itemize}
\item \textbf{Expérience 1 :} une lame de zinc plongée dans une solution de sulfate de cuivre se recouvre d'un dépôt rouge et la solution bleue se décolore.
\item \textbf{Expérience 2 :} une lame de cuivre plongée dans une solution de nitrate d'argent se recouvre d'un dépôt gris brillant et la solution bleuit.
\item \textbf{Expérience 3 :} une lame de cuivre plongée dans une solution de sulfate de zinc : aucune transformation observable.
\end{itemize}
Écrire les équations des réactions et classer les trois couples \ce{Zn^2+}/\ce{Zn}, \ce{Cu^2+}/\ce{Cu}, \ce{Ag+}/\ce{Ag} par pouvoir oxydant croissant.
\tcblower
\textbf{Étape 1 — Observations.} Exp. 1 : dépôt rouge = \ce{Cu} formé, décoloration = disparition de \ce{Cu^2+}. Exp. 2 : dépôt gris = \ce{Ag}, bleuissement = apparition de \ce{Cu^2+}. Exp. 3 : rien.

\textbf{Étapes 2 et 3 — Demi-équations.}

Expérience 1 : \ce{Cu^2+ + 2e- -> Cu} (réduction) et \ce{Zn -> Zn^2+ + 2e-} (oxydation).

Expérience 2 : \ce{Ag+ + e- -> Ag} (réduction, à multiplier par 2) et \ce{Cu -> Cu^2+ + 2e-} (oxydation).

\textbf{Étape 4 — Bilans.}
\[ \ce{Cu^2+ + Zn -> Cu + Zn^2+} \qquad \ce{2Ag+ + Cu -> 2Ag + Cu^2+} \]

\textbf{Étape 5 — Classification.} Exp. 1 : Zn s'oxyde, donc Zn est un réducteur plus fort que Cu : \ce{Cu^2+}/\ce{Cu} est au-dessus de \ce{Zn^2+}/\ce{Zn}. Exp. 2 : Cu s'oxyde devant \ce{Ag+} : \ce{Ag+}/\ce{Ag} est au-dessus de \ce{Cu^2+}/\ce{Cu}. L'expérience 3 confirme : Cu ne peut pas réduire \ce{Zn^2+}, ce qui est cohérent. D'où, par pouvoir oxydant croissant :
\[ \ce{Zn^2+}/\ce{Zn} \;<\; \ce{Cu^2+}/\ce{Cu} \;<\; \ce{Ag+}/\ce{Ag} \]
\end{exoresolu}

\courssub{Applications au quotidien camerounais}

\courssub{Pourquoi les tôles en fer rouillent-elles plus vite à Kribi et Douala ?}

Tout habitant du littoral le sait : une tôle neuve posée à Kribi ou à Limbe se perce de rouille en quelques années, alors qu'à Bafoussam ou à Ngaoundéré, elle tient bien plus longtemps. La classification électrochimique explique ce phénomène.

Le fer est un réducteur : son couple \ce{Fe^2+}/\ce{Fe} est placé \emph{au-dessous} de \ce{H3O+}/\ce{H2} et du couple \ce{O2}/\ce{H2O}. Le dioxygène de l'air, en présence d'eau, peut donc oxyder le fer spontanément :
\begin{align*}
\text{Oxydation : } & \ce{Fe -> Fe^2+ + 2e-} \\
\text{Réduction : } & \ce{O2 + 2H2O + 4e- -> 4OH-}
\end{align*}
Les ions \ce{Fe^2+} formés s'oxydent ensuite en rouille, un oxyde de fer hydraté. Deux facteurs accélèrent tout sur la côte :
\begin{itemize}
\item \textbf{l'humidité permanente} : l'eau est indispensable, elle transporte les ions et permet la circulation des charges ;
\item \textbf{le sel marin} : les ions \ce{Na+} et \ce{Cl-} du brouillard salin rendent l'eau bien plus conductrice, ce qui accélère le transfert d'électrons, exactement comme un pont salin accélère le fonctionnement d'une pile.
\end{itemize}
D'où les protections usuelles : peinture (barrière contre \ce{O2} et \ce{H2O}) ou galvanisation — le recouvrement de la tôle par du zinc, métal \emph{plus réducteur} que le fer, qui s'oxyde à sa place : c'est la \cle{protection sacrificielle}.

\courssub{Le bidon en fer et le sulfate de cuivre : un mélange interdit}

Le sulfate de cuivre \ce{CuSO4} est un fongicide très utilisé par les maraîchers (bouillie bordelaise pour traiter les tomates et les pommes de terre). Un cultivateur qui stockerait sa solution dans un bidon en fer aurait une mauvaise surprise : le bidon se perce et le produit perd son efficacité. En effet, \ce{Cu^2+}/\ce{Cu} est au-dessus de \ce{Fe^2+}/\ce{Fe}, donc la réaction spontanée est :
\[ \ce{Cu^2+ + Fe -> Cu + Fe^2+} \]
Le fer du bidon se dissout (trous, fuites) tandis que le cuivre se dépose en poudre rouge au fond : le fongicide est détruit. Il faut donc utiliser des récipients en plastique ou en inox.

\begin{popfigure}
\begin{center}
\begin{tikzpicture}[font=\small, scale=0.95]
% Bidon
\draw[thick, popdark, fill=popmuted!25] (0,0) -- (0,3.4) arc (180:0:1.5 and 0.35) -- (3,0) arc (360:180:1.5 and 0.35) -- cycle;
\draw[thick, popdark] (0,3.4) arc (180:360:1.5 and 0.35);
% Solution
\fill[popblue!35] (0.06,0.15) -- (0.06,2.3) arc (180:360:1.44 and 0.32) -- (2.94,0.15) arc (360:180:1.44 and 0.32) -- cycle;
\draw[popblue, thick] (0.06,2.3) arc (180:360:1.44 and 0.32);
\node[popblue!70!black] at (1.5,1.7) {solution de \ce{CuSO4}};
\node[popblue!70!black] at (1.5,1.25) {(ions \ce{Cu^2+} bleus)};
% Trou et fuite
\draw[fill=popdark] (0,0.7) circle (0.09);
\draw[-{Stealth}, poppink, very thick] (-0.05,0.7) -- (-1.1,0.35);
\node[poppink, align=center] at (-1.7,0.9) {\textbf{trou :}\\ le fer\\ s'est dissous};
% Dépôt de cuivre
\fill[poppink!70] (0.7,0.12) .. controls (1.1,0.45) and (1.9,0.45) .. (2.3,0.12) -- cycle;
\node[poppink!80!black, align=center] at (4.6,0.55) {dépôt de cuivre\\ métallique \ce{Cu}};
\draw[-{Stealth}, poppink!80!black] (3.7,0.5) -- (2.1,0.28);
% Paroi
\node[popdark, align=center] at (4.6,3.1) {paroi en fer \ce{Fe}};
\draw[-{Stealth}, popdark] (3.6,3.0) -- (2.98,2.9);
% Réaction
\node[align=center, draw=poporange, fill=poporangeL, rounded corners, thick] at (1.5,-1.15) {\ce{Cu^2+ + Fe -> Cu + Fe^2+}\\[-1pt] \footnotesize oxydant fort $+$ réducteur fort $\rightarrow$ réaction spontanée};
\end{tikzpicture}
\end{center}
\legende{Dégradation d'un bidon en fer contenant une solution de sulfate de cuivre : le fer, réducteur plus fort que le cuivre, s'oxyde et le bidon se perce.}
\end{popfigure}

\begin{exemplebox}[Marmites en cuivre et jus acides : pourquoi le cuivre résiste]
Dans certaines cuisines traditionnelles et chez les confiseurs, on utilise des marmites en cuivre. Or le jus de bissap, le jus de citron ou la sauce tomate contiennent des acides, donc des ions \ce{H3O+}. Pourquoi la marmite n'est-elle pas attaquée, alors qu'une marmite en fer s'userait vite ?

La réponse tient en une ligne de classification : le couple \ce{Cu^2+}/\ce{Cu} est placé \textbf{au-dessus} de \ce{H3O+}/\ce{H2}. Le cuivre est donc un réducteur \emph{trop faible} pour céder ses électrons aux ions \ce{H3O+} : la réaction
\[ \ce{2H3O+ + Cu -> H2 + Cu^2+ + 2H2O} \]
n'est \textbf{pas spontanée} : elle n'a pas lieu. À l'inverse, le couple \ce{Fe^2+}/\ce{Fe} est au-dessous de \ce{H3O+}/\ce{H2} : le fer est attaqué par les acides avec dégagement de dihydrogène :
\[ \ce{2H3O+ + Fe -> H2 ^ + Fe^2+ + 2H2O} \]
C'est exactement le critère expérimental qui a permis de placer le couple \ce{H3O+}/\ce{H2} dans la classification. Remarque pratique : à l'air, le cuivre finit quand même par se recouvrir lentement de vert-de-gris, car le dioxygène est un oxydant plus puissant que \ce{H3O+} — c'est pourquoi on étame (on recouvre d'étain) l'intérieur des marmites en cuivre.
\end{exemplebox}

\courssub{La cimentation : récupérer un métal précieux par déplacement}

La \cle{cimentation} consiste à récupérer un métal noble d'une solution en le faisant déplacer par un métal moins noble et bon marché, presque toujours le fer (souvent de la ferraille recyclée : vieilles tôles, clous, chutes de ferblanterie). Le principe est une application directe de la classification : tout métal placé au-dessous dans la classification réduit les ions d'un métal placé au-dessus.

Exemple chiffré : un artisan fait passer \SI{10}{L} d'une solution contenant \SI{0,05}{mol.L^{-1}} d'ions \ce{Cu^2+} sur de la ferraille. La quantité d'ions cuivre est $n(\ce{Cu^2+}) = C \times V = 0{,}05 \times 10 = \SI{0,5}{mol}$. D'après l'équation $\ce{Cu^2+ + Fe -> Cu + Fe^2+}$, il se dépose $n(\ce{Cu}) = \SI{0,5}{mol}$, soit une masse :
\[ m(\ce{Cu}) = n \times M = 0{,}5 \times 63{,}5 \approx \SI{32}{g} \text{ de cuivre récupéré,} \]
en consommant $m(\ce{Fe}) = 0{,}5 \times 55{,}8 \approx \SI{28}{g}$ de ferraille sans valeur. Avec des ions argent \ce{Ag+} (couple \ce{Ag+}/\ce{Ag}, encore plus haut dans la classification), le même principe permet de récupérer l'argent des bains de photographie ou de radiographie usagés :
\[ \ce{2Ag+ + Fe -> 2Ag + Fe^2+} \]

\begin{savaistu}[La cimentation, une technique vieille de plusieurs siècles]
La cimentation du cuivre par le fer est connue depuis l'Antiquité : les mineurs chinois de la dynastie Song et, plus tard, les exploitants des mines de Rio Tinto en Espagne récupéraient le cuivre des eaux de mine en y plongeant des fers de lance ou de la ferraille. Au XIX\textsuperscript{e} siècle, ce procédé fournissait une part importante du cuivre mondial. Il est encore utilisé aujourd'hui dans l'industrie minière moderne (hydrométallurgie) et dans le traitement des déchets d'ateliers de bijouterie et de galvanoplastie. Au Cameroun, où la ferraille abonde et où l'artisanat de récupération est très développé, ce principe illustre parfaitement comment une loi de chimie — la classification électrochimique — se transforme en technique économique : un déchet de fer contre un métal de valeur.
\end{savaistu}

\courssub{Vers la suite du programme : piles et dosages}

Tout ce que tu viens d'apprendre est le socle de deux grands chapitres à venir.

\begin{itemize}
\item \textbf{Les piles électrochimiques (Terminale).} Quand on sépare physiquement les deux couples d'une réaction spontanée et qu'on force les électrons à passer par un circuit extérieur, on fabrique une \cle{pile} : la réaction $\ce{Cu^2+ + Zn -> Cu + Zn^2+}$ devient la pile Daniell, qui débite un courant électrique. La classification qualitative de Première deviendra quantitative grâce au \cle{potentiel standard d'électrode} $E^\circ$ (en volts), mesuré par rapport au couple \ce{H3O+}/\ce{H2} pris comme référence ($E^\circ = \SI{0,00}{V}$). Tu comprends maintenant pourquoi ce couple joue un rôle si central : c'est la « borne zéro » de toute l'électrochimie.
\item \textbf{Les dosages d'oxydoréduction.} De la même façon qu'on dose un acide par une base, on pourra doser un réducteur par un oxydant de concentration connue (par exemple le fer(II) par le permanganate de potassium, dont la couleur violette sert d'indicateur). L'équivalence correspondra à l'égalité des quantités d'électrons échangés.
\end{itemize}

\begin{aretenir}[L'essentiel de la leçon]
\begin{itemize}
\item Un couple Ox/Red est lié par la demi-équation $\mathrm{Ox} + n\,e^- \rightleftharpoons \mathrm{Red}$.
\item La réaction spontanée oppose l'oxydant le plus fort au réducteur le plus fort.
\item La classification électrochimique se construit par des expériences de déplacement ; le couple \ce{H3O+}/\ce{H2} sépare les métaux attaqués par les acides (Fe, Zn\ldots) de ceux qui y résistent (Cu, Ag\ldots).
\item Applications directes : corrosion et protection des métaux, choix des récipients, cimentation, et bientôt les piles et les dosages.
\end{itemize}
\end{aretenir}

\begin{exercice}[Pour vérifier ta maîtrise]
Un bijoutier artisan de Douala récupère l'argent de ses bains usagés contenant des ions \ce{Ag+} en y plongeant de la ferraille.
\begin{enumerate}
\item Écrire les deux demi-équations et l'équation bilan de la réaction.
\item Justifier, à l'aide de la classification, que cette réaction est spontanée. Pourrait-on utiliser du cuivre à la place du fer ?
\item Calculer la masse d'argent récupérée à partir d'un bain contenant \SI{0,20}{mol} d'ions \ce{Ag+} (on donne $M(\ce{Ag}) = \SI{107,9}{g.mol^{-1}}$).
\end{enumerate}
\end{exercice}

\begin{corrige}[Exercice : récupération de l'argent]
\textbf{1.} Réduction : \ce{Ag+ + e- -> Ag}. Oxydation : \ce{Fe -> Fe^2+ + 2e-}. En multipliant la première par 2 puis en additionnant :
\[ \ce{2Ag+ + Fe -> 2Ag + Fe^2+} \]
\textbf{2.} Le couple \ce{Ag+}/\ce{Ag} est placé au-dessus de \ce{Fe^2+}/\ce{Fe} : \ce{Ag+} est un oxydant plus fort que \ce{Fe^2+} et Fe un réducteur plus fort que Ag, donc la réaction est spontanée. Oui, le cuivre conviendrait aussi, car \ce{Ag+}/\ce{Ag} est également au-dessus de \ce{Cu^2+}/\ce{Cu} : $\ce{2Ag+ + Cu -> 2Ag + Cu^2+}$.
\textbf{3.} $m(\ce{Ag}) = n \times M = 0{,}20 \times 107{,}9 = \SI{21,6}{g}$ d'argent récupéré.
\end{corrige}

\newpage

\coursec{Travaux pratiques}

\courssub{Objectifs du TP}

Ce travail pratique te permet de passer de la théorie à l'expérience. À l'issue de la séance, tu devras être capable de :

\begin{itemize}
\item réaliser des expériences de mise en évidence de quelques \cle{couples oxydant-réducteur} ;
\item réaliser des expériences mettant en jeu les couples $\ce{M^{n+}}/\ce{M}$ et $\ce{H3O+}/\ce{H2}$ ;
\item identifier expérimentalement les produits formés (dépôt métallique, dégagement de dihydrogène) ;
\item écrire les \cle{demi-équations électroniques} et le \cle{bilan} de chaque réaction observée ;
\item établir une \cle{classification électrochimique qualitative} des couples étudiés et y placer le couple $\ce{H3O+}/\ce{H2}$.
\end{itemize}

\courssub{Matériel et produits}

\begin{center}
\rowcolors{2}{popblueL}{white}
\begin{tabularx}{\linewidth}{|l|X|}
\hline
\rowcolor{popblue!40} \textbf{Matériel} & \textbf{Produits chimiques} \\
\hline
6 tubes à essai + portoir & Solution de sulfate de cuivre(II) \ce{CuSO4} à \SI{0,5}{mol.L^{-1}} (bleue) \\
\hline
1 clou en fer décapé (ou limaille de fer) & Solution de sulfate de fer(II) \ce{FeSO4} à \SI{0,5}{mol.L^{-1}} (verdâtre) \\
\hline
Fil de cuivre dénudé (câble électrique) & Solution de nitrate d'argent \ce{AgNO3} à \SI{0,1}{mol.L^{-1}} (incolore) \\
\hline
Grenaille de zinc (ou lame de zinc) & Acide chlorhydrique dilué \ce{HCl} à \SI{1}{mol.L^{-1}} \\
\hline
Pipettes compte-gouttes, spatule & Eau distillée \\
\hline
Allumettes ou bec Bunsen & \\
\hline
Lunettes de protection, gants, blouse & \\
\hline
\end{tabularx}
\end{center}

\begin{attention}[Alternatives avec du matériel local]
Si le nitrate d'argent n'est pas disponible (produit coûteux), remplace l'expérience 3 par un fil de cuivre plongé dans une solution de sulfate de zinc : l'absence de réaction fournit aussi une information de classement. Le zinc peut être récupéré sur le boîtier d'une vieille pile usagée (électrode de zinc des piles Leclanché), soigneusement nettoyé. Le cuivre provient d'un câble électrique dénudé, le fer d'un clou de menuiserie frotté à la toile émeri.
\end{attention}

\courssub{Sécurité}

\begin{attention}[Consignes de sécurité obligatoires]
\begin{itemize}
\item \textbf{Lunettes et blouse} portées pendant toute la séance ; cheveux longs attachés.
\item \textbf{Acide chlorhydrique} : produit \emph{corrosif} (pictogramme : deux gouttes rongeant une surface et une main). Verser toujours l'acide doucement le long de la paroi du tube, jamais le visage au-dessus. En cas de contact avec la peau, rincer abondamment à l'eau pendant plusieurs minutes et prévenir l'enseignant.
\item \textbf{Nitrate d'argent} : produit \emph{tachant} (taches noires tenaces sur la peau) et \emph{dangereux pour l'environnement}. Manipuler avec des gants, ne pas rejeter à l'évier : récupérer les résidus dans un bidon prévu à cet effet.
\item \textbf{Test à la flamme} : le dihydrogène est un gaz \emph{inflammable} (pictogramme flamme). N'approcher la flamme que de l'extrémité du tube à essai, jamais du flacon d'acide. Tenir le tube incliné, ouverture dirigée vers personne.
\item Ne jamais goûter ni sentir directement un produit ; signaler immédiatement toute casse de verrerie.
\end{itemize}
\end{attention}

\courssub{Schéma du dispositif}

\begin{popfigure}
\begin{center}
\begin{tikzpicture}[scale=0.9]
% Tube 1 : clou dans CuSO4
\draw[thick,popdark] (0,0) -- (0,2.6);
\draw[thick,popdark] (1.2,0) -- (1.2,2.6);
\draw[thick,popdark] (0,0) arc (180:360:0.6);
\fill[popblue!35] (0.03,0.12) -- (0.03,1.6) -- (1.17,1.6) -- (1.17,0.12) arc (180:360:0.57);
\draw[thick,popmuted] (0.45,2.9) -- (0.75,0.25);
\node[font=\small,align=center] at (0.6,-0.9) {Exp. 1 : clou Fe\\ dans \ce{CuSO4}};
% Tube 2 : fil cuivre dans FeSO4
\begin{scope}[xshift=3.2cm]
\draw[thick,popdark] (0,0) -- (0,2.6);
\draw[thick,popdark] (1.2,0) -- (1.2,2.6);
\draw[thick,popdark] (0,0) arc (180:360:0.6);
\fill[popgreen!25] (0.03,0.12) -- (0.03,1.6) -- (1.17,1.6) -- (1.17,0.12) arc (180:360:0.57);
\draw[thick,poporange] (0.55,2.9) -- (0.65,0.3);
\node[font=\small,align=center] at (0.6,-0.9) {Exp. 2 : fil Cu\\ dans \ce{FeSO4}};
\end{scope}
% Tube 3 : Cu dans AgNO3
\begin{scope}[xshift=6.4cm]
\draw[thick,popdark] (0,0) -- (0,2.6);
\draw[thick,popdark] (1.2,0) -- (1.2,2.6);
\draw[thick,popdark] (0,0) arc (180:360:0.6);
\fill[popcream] (0.03,0.12) -- (0.03,1.6) -- (1.17,1.6) -- (1.17,0.12) arc (180:360:0.57);
\draw[thick,popmuted] (0.5,2.9) -- (0.6,0.3);
\draw[popmuted,fill=popmuted!50] (0.4,0.4) rectangle (0.8,0.6);
\node[font=\small,align=center] at (0.6,-0.9) {Exp. 3 : fil Cu\\ dans \ce{AgNO3}};
\end{scope}
% Tube 4 : zinc dans HCl + flamme (nouvelle ligne)
\begin{scope}[yshift=-4.5cm]
\draw[thick,popdark] (0,0) -- (0,2.6);
\draw[thick,popdark] (1.2,0) -- (1.2,2.6);
\draw[thick,popdark] (0,0) arc (180:360:0.6);
\fill[popcream] (0.03,0.12) -- (0.03,1.4) -- (1.17,1.4) -- (1.17,0.12) arc (180:360:0.57);
\fill[popmuted] (0.35,0.25) circle (0.09);
\fill[popmuted] (0.7,0.2) circle (0.09);
\fill[popmuted] (0.55,0.45) circle (0.09);
\foreach \x/\y in {0.4/0.9,0.8/1.1,0.6/1.3}{\draw[popblue] (\x,\y) circle (0.04);}
\draw[thick,popgold] (1.9,2.2) -- (1.9,2.7);
\draw[poporange,fill=poporange!60] (1.9,2.95) ellipse (0.09 and 0.25);
\draw[->,thick,popdark] (1.75,2.4) -- (1.3,2.4);
\node[font=\small,align=center] at (0.6,-0.9) {Exp. 4 : Zn + \ce{HCl}\\ test à la flamme};
\end{scope}
\end{tikzpicture}
\end{center}
\legende{Dispositifs expérimentaux : métaux plongés dans des solutions d'ions métalliques ou d'acide chlorhydrique dilué, avec test d'identification du dihydrogène à la flamme.}
\end{popfigure}

\courssub{Protocole}

\begin{enumerate}
\item \textbf{Expérience 1 — Fer et ions cuivre(II).} Verser environ \SI{5}{mL} de solution de sulfate de cuivre(II) dans un tube à essai. Y plonger le clou en fer préalablement décapé. Attendre 5 minutes en agitant doucement de temps en temps. Observer la surface du clou et la couleur de la solution.
\item \textbf{Expérience 2 — Cuivre et ions fer(II).} Verser environ \SI{5}{mL} de solution de sulfate de fer(II) dans un deuxième tube. Y plonger le fil de cuivre dénudé. Attendre 5 minutes. Observer.
\item \textbf{Expérience 3 — Cuivre et ions argent.} Verser environ \SI{3}{mL} de solution de nitrate d'argent dans un troisième tube. Y plonger le fil de cuivre. Observer pendant 5 minutes (dépôt gris-argenté en « aiguilles »).
\item \textbf{Expérience 4 — Action de l'acide sur les métaux.} Dans trois tubes séparés, introduire respectivement quelques grains de zinc, un petit clou en fer et un morceau de fil de cuivre. Ajouter dans chacun environ \SI{3}{mL} d'acide chlorhydrique dilué. Observer l'effervescence éventuelle et comparer sa vigueur.
\item \textbf{Test du gaz.} Pour les tubes où un dégagement gazeux est visible, boucher le tube avec un bouchon (ou le pouce ganté) pendant 30 secondes, puis approcher rapidement une flamme d'allumette de l'ouverture : une \cle{détonation} caractéristique (petit « aboiement ») prouve la présence de dihydrogène \ce{H2}.
\item Noter immédiatement toutes les observations dans le tableau des résultats.
\end{enumerate}

\courssub{Observations et résultats attendus}

\begin{center}
\rowcolors{2}{popblueL}{white}
\begin{tabularx}{\linewidth}{|l|X|X|}
\hline
\rowcolor{popblue!40} \textbf{Expérience} & \textbf{Observations} & \textbf{Conclusion immédiate} \\
\hline
1. Fe dans \ce{CuSO4} & Dépôt rougeâtre de cuivre sur le clou ; la solution bleue pâlit progressivement & Le fer réduit les ions \ce{Cu^2+} ; réaction spontanée \\
\hline
2. Cu dans \ce{FeSO4} & Aucun changement visible après 5 min & Pas de réaction : le cuivre ne réduit pas \ce{Fe^2+} \\
\hline
3. Cu dans \ce{AgNO3} & Dépôt gris-argenté (cristaux d'argent) sur le fil ; la solution bleuit lentement & Le cuivre réduit les ions \ce{Ag+} ; réaction spontanée \\
\hline
4a. Zn + \ce{HCl} & Effervescence vive ; détonation à la flamme & Dégagement de \ce{H2} ; le zinc est attaqué \\
\hline
4b. Fe + \ce{HCl} & Effervescence modérée ; détonation à la flamme & Dégagement de \ce{H2} ; le fer est attaqué \\
\hline
4c. Cu + \ce{HCl} & Aucune effervescence & Le cuivre n'est pas attaqué par \ce{H3O+} \\
\hline
\end{tabularx}
\end{center}

\courssub{Exploitation}

\begin{exoresolu}[Questions guidées d'exploitation]
\textbf{Question 1.} Écrire les demi-équations électroniques et l'équation-bilan de la réaction de l'expérience 1. Quels couples interviennent ?

\textbf{Question 2.} L'expérience 2 montre-t-elle que le cuivre est un réducteur plus fort ou plus faible que le fer ? Justifier.

\textbf{Question 3.} Écrire le bilan de l'expérience 3 et classer les couples $\ce{Ag+}/\ce{Ag}$ et $\ce{Cu^2+}/\ce{Cu}$.

\textbf{Question 4.} Écrire le bilan de l'action de l'acide chlorhydrique sur le zinc. Quel couple contenant \ce{H3O+} intervient ?

\textbf{Question 5.} Déduire de l'ensemble des expériences la classification qualitative des cinq couples étudiés, du réducteur le plus fort au réducteur le plus faible. Où se situe le couple $\ce{H3O+}/\ce{H2}$ ?
\tcblower
\textbf{Réponse 1.} Le fer perd des électrons (oxydation) : $\ce{Fe -> Fe^2+ + 2e-}$. Les ions cuivre(II) les captent (réduction) : $\ce{Cu^2+ + 2e- -> Cu}$. Bilan :
\[ \ce{Fe + Cu^2+ -> Fe^2+ + Cu} \]
Couples mis en jeu : $\ce{Cu^2+}/\ce{Cu}$ et $\ce{Fe^2+}/\ce{Fe}$.

\textbf{Réponse 2.} Aucune réaction n'a lieu entre le cuivre et les ions \ce{Fe^2+} : le cuivre est incapable de céder ses électrons à \ce{Fe^2+}. Le cuivre est donc un réducteur \textbf{plus faible} que le fer. On écrit : $\ce{Cu^2+}/\ce{Cu}$ est au-dessus de $\ce{Fe^2+}/\ce{Fe}$ dans la classification.

\textbf{Réponse 3.} Oxydation du cuivre : $\ce{Cu -> Cu^2+ + 2e-}$ ; réduction de l'argent : $\ce{Ag+ + e- -> Ag}$ (à multiplier par 2). Bilan :
\[ \ce{Cu + 2Ag+ -> Cu^2+ + 2Ag} \]
Le cuivre réduit \ce{Ag+}, donc $\ce{Ag+}/\ce{Ag}$ est classé au-dessus de $\ce{Cu^2+}/\ce{Cu}$.

\textbf{Réponse 4.} Oxydation : $\ce{Zn -> Zn^2+ + 2e-}$ ; réduction : $\ce{2H3O+ + 2e- -> H2 + 2H2O}$. Bilan :
\[ \ce{Zn + 2H3O+ -> Zn^2+ + H2 ^ + 2H2O} \]
Le couple contenant l'ion oxonium est $\ce{H3O+}/\ce{H2}$.

\textbf{Réponse 5.} Le zinc et le fer réduisent \ce{H3O+} (dégagement de \ce{H2}), mais pas le cuivre ; le fer réduit \ce{Cu^2+} ; le cuivre réduit \ce{Ag+}. La classification, du réducteur le plus fort au plus faible, est donc :
\[ \ce{Zn^2+}/\ce{Zn} \;<\; \ce{Fe^2+}/\ce{Fe} \;<\; \ce{H3O+}/\ce{H2} \;<\; \ce{Cu^2+}/\ce{Cu} \;<\; \ce{Ag+}/\ce{Ag} \]
Le couple $\ce{H3O+}/\ce{H2}$ se situe \textbf{entre} les couples $\ce{Fe^2+}/\ce{Fe}$ et $\ce{Cu^2+}/\ce{Cu}$ : c'est la frontière entre les métaux attaqués par les acides (Zn, Fe) et ceux qui ne le sont pas (Cu, Ag).
\end{exoresolu}

\begin{methode}[Règle du gamma pour prévoir une réaction]
\begin{enumerate}
\item Placer les deux couples sur l'axe de classification vertical (réducteur le plus fort en bas, oxydant le plus fort en haut).
\item Tracer le « gamma » ($\gamma$) : le réducteur le plus fort (en bas à droite) réagit avec l'oxydant le plus fort (en haut à gauche), placés en diagonale.
\item Conclure : la réaction spontanée est celle qui relie ces deux espèces ; l'autre sens est impossible.
\end{enumerate}
Exemple : \ce{Fe} (bas) et \ce{Cu^2+} (haut) sont en diagonale $\Rightarrow$ réaction. \ce{Cu} et \ce{Fe^2+} sont sur l'autre diagonale $\Rightarrow$ aucune réaction.
\end{methode}

\begin{aretenir}[À retenir de la séance]
\begin{itemize}
\item Un \cle{couple oxydant-réducteur} s'écrit $\ce{Ox}/\ce{Red}$ (ex. $\ce{Cu^2+}/\ce{Cu}$). L'oxydant capte les électrons, le réducteur en cède.
\item Une réaction d'oxydoréduction spontanée met en jeu un oxydant et un réducteur tels que l'oxydant est \textbf{plus fort} que l'oxydant du couple réducteur (règle du gamma).
\item La \cle{classification électrochimique qualitative} ordonne les couples par pouvoir oxydant/réducteur croissant. Pour les couples métalliques $\ce{M^{n+}}/\ce{M}$ et le couple $\ce{H3O+}/\ce{H2}$ :
\[ \ce{Zn^2+}/\ce{Zn} < \ce{Fe^2+}/\ce{Fe} < \ce{H3O+}/\ce{H2} < \ce{Cu^2+}/\ce{Cu} < \ce{Ag+}/\ce{Ag} \]
\item Les métaux situés \textbf{au-dessous} de $\ce{H3O+}/\ce{H2}$ (Zn, Fe, Mg, Al...) sont attaqués par les acides avec dégagement de \ce{H2}. Ceux \textbf{au-dessus} (Cu, Ag, Au...) ne le sont pas.
\end{itemize}
\end{aretenir}

\begin{savaistu}[Applications au quotidien camerounais]
\begin{itemize}
\item \textbf{Corrosion et protection :} Le fer (toitures, charpentes, fûts à bili-bili) rouille car $\ce{Fe^2+}/\ce{Fe}$ est sous $\ce{O2}/\ce{H2O}$ et $\ce{H3O+}/\ce{H2}$. La galvanisation (zinc sur fer) protège : le zinc ($\ce{Zn^2+}/\ce{Zn}$ plus bas) s'oxyde à la place du fer (anode sacrificielle).
\item \textbf{Cuisine et conservation :} Le vinaigre (acide acétique, source de \ce{H3O+}) attaque les casseroles en aluminium ($\ce{Al^3+}/\ce{Al}$ très bas) mais pas celles en inox (fer passivé) ni en cuivre. C'est pourquoi le \emph{vin de palme} ou le \emph{bissap} acides ne se conservent pas longtemps dans des récipients métalliques non adaptés.
\item \textbf{Bijouterie et monnaie :} L'argent et le cuivre (alliages) résistent aux acides faibles (jus de fruit, sueur) car leurs couples sont au-dessus de $\ce{H3O+}/\ce{H2}$. En revanche, le nitrate d'argent (rayons X, pharmacopée) noircit la peau par réduction en Ag métallique au contact des réducteurs organiques.
\item \textbf{Piles et batteries :} Les piles alcalines (Zn/MnO2) et les accumulateurs plomb-acide (voitures, groupes électrogènes CIMENCAM/SONARA) reposent sur des couples choisis aux extrémités de la classification pour maximiser la tension.
\end{itemize}
\end{savaistu}

\courssub{Conclusion}

Les expériences réalisées montrent qu'un métal ne réagit avec les ions d'un autre métal que s'il est un \cle{réducteur plus fort} que celui-ci. En comparant deux à deux les couples $\ce{M^{n+}}/\ce{M}$, on construit une \cle{classification électrochimique qualitative} :

\[ \ce{Zn^2+}/\ce{Zn} \;<\; \ce{Fe^2+}/\ce{Fe} \;<\; \ce{H3O+}/\ce{H2} \;<\; \ce{Cu^2+}/\ce{Cu} \;<\; \ce{Ag+}/\ce{Ag} \]

Le couple $\ce{H3O+}/\ce{H2}$ occupe une place charnière : les métaux situés \emph{en dessous} de lui (Zn, Fe, Mg, Al, etc.) sont attaqués par les acides à ions \ce{H3O+} avec dégagement de dihydrogène ; ceux situés \emph{au-dessus} (Cu, Ag, Au, Pt) y sont inertes. C'est précisément pourquoi les bijoux en cuivre ou en argent ne sont pas rongés par le vinaigre de la cuisine, alors qu'un clou en fer laissé dans un jus acide (bissap, vin de palme) se dégrade rapidement. La classification électrochimique est donc un outil de \emph{prévision} puissant : elle permet d'annoncer, sans refaire l'expérience, si une réaction d'oxydoréduction entre un métal et un ion est possible ou non.

\newpage

\coursec{Méthodes et exercices résolus}

\coursec{Leçon 06 : Couple oxydant-réducteur et classification électrochimique}

\courssub{1. Couple oxydant-réducteur : définition et écriture}

\begin{definition}[Couple oxydant-réducteur]
Un \cle{couple oxydant-réducteur} (noté Ox/Red) est un ensemble de deux espèces chimiques (l'\cle{oxydant} Ox et le \cle{réducteur} Red) qui ne diffèrent que par un nombre entier $n$ d'électrons. L'oxydant est l'espèce qui \textbf{capte} les électrons (il s'\cle{réduit}), le réducteur est l'espèce qui \textbf{cède} les électrons (il s'\cle{oxyde}).
La transformation s'écrit sous forme de \cle{demi-équation électronique} (ou demi-réaction) dans le sens de la réduction :
\[ \ce{Ox + n e- <=> Red} \]
\end{definition}

\begin{propriete}[Écriture et équilibrage d'une demi-équation]
Pour écrire la demi-équation d'un couple Ox/Red en milieu aqueux :
\begin{enumerate}
\item Écrire l'oxydant à gauche, le réducteur à droite, séparés par $\ce{<=>}$.
\item Équilibrer les atomes autres que O et H.
\item Équilibrer l'oxygène en ajoutant $\ce{H2O}$.
\item Équilibrer l'hydrogène en ajoutant $\ce{H3O+}$ (milieu acide) ou $\ce{OH-}$ (milieu basique).
\item Équilibrer la charge en ajoutant des électrons $\ce{e-}$ du côté le plus positif.
\end{enumerate}
Le couple se note toujours \textbf{Ox/Red} (ex. : \ce{Cu^2+}/\ce{Cu}, \ce{MnO4-}/\ce{Mn^2+}, \ce{H3O+}/\ce{H2}).
\end{propriete}

\begin{propriete}[Couples métalliques $\ce{M^{n+}}/\ce{M}$]
Pour un métal M et son cation $\ce{M^{n+}}$, le couple s'écrit $\ce{M^{n+}}/\ce{M}$.
La demi-équation est : $\ce{M^{n+} + n e- <=> M}$.
\begin{itemize}
\item L'oxydant est le cation métallique $\ce{M^{n+}}$.
\item Le réducteur est le métal $\ce{M}$ (solide, état standard).
\end{itemize}
\end{propriete}

\begin{methode}[Identifier un couple oxydant-réducteur et écrire sa demi-équation]
Pour reconnaître et écrire correctement un couple Ox/Red :
\begin{enumerate}
\item Repérer dans l'énoncé les deux espèces conjuguées : l'une contient l'élément à un degré d'oxydation élevé (l'\cle{oxydant}), l'autre à un degré d'oxydation plus faible (le \cle{réducteur}).
\item Vérifier que les deux espèces ne diffèrent que par des électrons (éventuellement avec $\ce{H3O+}$ et $\ce{H2O}$ pour équilibrer O et H en milieu acide).
\item Écrire la demi-équation dans le sens \textbf{oxydant} $+ n\,e^-$ $\ce{<=>}$ \textbf{réducteur} : l'oxydant est toujours écrit à gauche, les électrons à côté de lui.
\item Équilibrer dans l'ordre : les atomes autres que O et H, puis O avec $\ce{H2O}$, puis H avec $\ce{H3O+}$, enfin la charge avec les électrons.
\item Noter le couple sous la forme Ox/Red : $\ce{Cu^2+}/\ce{Cu}$, $\ce{Zn^2+}/\ce{Zn}$, $\ce{H3O+}/\ce{H2}$.
\end{enumerate}
\textbf{Réflexe :} dans un couple $\ce{M^{n+}}/\ce{M}$, le cation métallique est l'oxydant, le métal est le réducteur.
\end{methode}

\begin{exoresolu}[Reconnaître les couples et écrire les demi-équations]
Lors d'une séance de travaux pratiques au lycée de Nkongsamba, on réalise les expériences suivantes : (a) un clou en fer plongé dans une solution de sulfate de cuivre(II) se recouvre d'un dépôt rougeâtre ; (b) une lame de zinc plongée dans une solution de sulfate de fer(II) se recouvre d'un dépôt gris foncé.
\begin{enumerate}
\item Identifier les couples oxydant-réducteur mis en jeu dans chaque expérience.
\item Écrire la demi-équation électronique de chaque couple.
\item Écrire l'équation-bilan de la réaction de l'expérience (a).
\end{enumerate}
\tcblower
\textbf{1. Identification des couples.}

Expérience (a) : le dépôt rougeâtre est du cuivre métallique $\ce{Cu}$ ; le fer $\ce{Fe}$ s'est transformé en ions $\ce{Fe^2+}$. Couples mis en jeu : $\ce{Cu^2+}/\ce{Cu}$ et $\ce{Fe^2+}/\ce{Fe}$.

Expérience (b) : le dépôt gris est du fer $\ce{Fe}$ ; le zinc $\ce{Zn}$ s'est transformé en ions $\ce{Zn^2+}$. Couples mis en jeu : $\ce{Fe^2+}/\ce{Fe}$ et $\ce{Zn^2+}/\ce{Zn}$.

\textbf{2. Demi-équations électroniques} (oxydant à gauche, électrons à gauche) :
\begin{align*}
\ce{Cu^2+ + 2e- &<=> Cu} \\
\ce{Fe^2+ + 2e- &<=> Fe} \\
\ce{Zn^2+ + 2e- &<=> Zn}
\end{align*}

\textbf{3. Équation-bilan de l'expérience (a).}

Le fer est le réducteur qui cède des électrons (oxydation) : $\ce{Fe -> Fe^2+ + 2e-}$.

L'ion cuivre(II) est l'oxydant qui capte des électrons (réduction) : $\ce{Cu^2+ + 2e- -> Cu}$.

Les électrons échangés sont au même nombre ($2\,e^-$) : on additionne membre à membre et on simplifie les électrons :
\[ \ce{Fe(s) + Cu^2+(aq) -> Fe^2+(aq) + Cu(s)} \]

\textbf{Conclusion :} le fer, réducteur plus fort que le cuivre, réduit les ions $\ce{Cu^2+}$ ; l'équation-bilan est $\ce{Fe + Cu^2+ -> Fe^2+ + Cu}$.
\end{exoresolu}

\begin{popfigure}
\begin{tikzpicture}[scale=0.9, every node/.style={font=\small}]
% Echelle verticale
\draw[popline, thick, ->] (0,0) -- (0,7) node[above] {\textbf{Pouvoir réducteur croissant}};
\draw[popline, thick, ->] (2,0) -- (2,7) node[above] {\textbf{Pouvoir oxydant croissant}};
% Couples
\foreach \y/\ox/\red/\col in {1/\ce{Cu^2+}/\ce{Cu}/popblue, 2/\ce{H3O+}/\ce{H2}/poporange, 3/\ce{Fe^2+}/\ce{Fe}/popgreen, 4/\ce{Zn^2+}/\ce{Zn}/poppurple}{
    \draw[\col, thick] (-0.5,\y) -- (0.5,\y) node[midway, left=2mm] {\ox/\red};
    \draw[\col, thick] (1.5,\y) -- (2.5,\y);
    \node[circle, fill=\col, inner sep=1.5pt] at (0,\y) {};
    \node[circle, fill=\col, inner sep=1.5pt] at (2,\y) {};
}
% Flèche diagonale gamma
\draw[poppink, very thick, ->, >=Stealth] (0.5,1) .. controls (1,1.5) and (1,2.5) .. (1.5,3) node[midway, right, poppink] {\textbf{Règle du $\gamma$}};
\node[align=center, popmuted, font=\footnotesize] at (3.5,3.5) {Réaction spontanée :\\ Oxydant fort + Réducteur fort};
\end{tikzpicture}
\legende{Représentation schématique de la règle du $\gamma$ (diagonale) dans une classification où le pouvoir réducteur croît vers le haut.}
\end{popfigure}

\courssub{2. Réaction entre un métal et un ion métallique : couples $\ce{M^{n+}}/\ce{M}$}

\begin{propriete}[Règle du $\gamma$ (gamma) et déplacement]
Dans une classification électrochimique où les réducteurs sont classés par pouvoir réducteur \textbf{croissant vers le bas} (ou vers la gauche) :
\begin{itemize}
\item L'oxydant le plus fort est en haut (ou à droite).
\item Le réducteur le plus fort est en bas (ou à gauche).
\item La réaction spontanée (naturelle) a lieu entre l'\textbf{oxydant le plus fort} et le \textbf{réducteur le plus fort} : c'est la \cle{règle du $\gamma$} (diagonale descendante).
\item Un métal $\ce{M_1}$ réduit les ions $\ce{M_2^{n+}}$ \textbf{si et seulement si} $\ce{M_1}$ est un réducteur plus fort que $\ce{M_2}$ (c'est-à-dire placé \textbf{au-dessous} de $\ce{M_2}$ dans la classification « réducteur fort en bas »).
\item On dit que $\ce{M_1}$ \textbf{déplace} $\ce{M_2}$ de sa solution.
\end{itemize}
\end{propriete}

\begin{methode}[Prévoir si un métal réagit avec un ion métallique (couples $\ce{M^{n+}}/\ce{M}$)]
Pour savoir si un métal $\ce{M_1}$ peut réduire les ions $\ce{M_2^{n+}}$ :
\begin{enumerate}
\item Écrire les deux couples concernés : $\ce{M_1^{n+}}/\ce{M_1}$ et $\ce{M_2^{n+}}/\ce{M_2}$.
\item Placer les deux couples dans la \cle{classification électrochimique} (de mémoire ou fournie) : le réducteur le plus fort est le métal placé le plus \textbf{bas} (convention « réducteur fort en bas »).
\item Appliquer la \cle{règle du $\gamma$} : la réaction naturelle a lieu entre l'oxydant le plus fort et le réducteur le plus fort, c'est-à-dire en diagonale.
\item Conclure : le métal $\ce{M_1}$ réduit les ions $\ce{M_2^{n+}}$ \textbf{uniquement si} $\ce{M_1}$ est un réducteur plus fort que $\ce{M_2}$ ($\ce{M_1}$ placé au-dessous de $\ce{M_2}$).
\item Écrire l'équation-bilan en égalisant le nombre d'électrons échangés.
\end{enumerate}
\textbf{Réflexe :} un métal déplace de ses solutions tout métal situé \emph{au-dessus} de lui dans la classification (réducteur moins fort).
\end{methode}

\begin{exoresolu}[Le métal zinc peut-il débarrasser une eau de ses ions cuivre(II) ?]
Pour traiter une eau de rejet d'un atelier de gravure contenant des ions $\ce{Cu^2+}$, un technicien de Douala propose d'y plonger de la grenaille de zinc. On donne l'extrait de classification électrochimique (réducteur de plus en plus fort vers le bas) : $\ce{Cu^2+}/\ce{Cu}$ puis $\ce{Zn^2+}/\ce{Zn}$.
\begin{enumerate}
\item Justifier que le zinc réduit les ions $\ce{Cu^2+}$ et écrire l'équation-bilan.
\item On plonge $\SI{6,54}{g}$ de zinc dans $\SI{500}{mL}$ d'une solution de sulfate de cuivre(II) de concentration $C = \SI{0,20}{mol.L^{-1}}$. Déterminer le réactif limitant.
\item Calculer la masse de cuivre déposée.
\end{enumerate}
Données : $M(\ce{Zn}) = \SI{65,4}{g.mol^{-1}}$ ; $M(\ce{Cu}) = \SI{63,5}{g.mol^{-1}}$.
\tcblower
\textbf{1. Prévision de la réaction.}

Dans la classification, le zinc est placé \emph{au-dessous} du cuivre : $\ce{Zn}$ est un réducteur plus fort que $\ce{Cu}$, et $\ce{Cu^2+}$ est un oxydant plus fort que $\ce{Zn^2+}$. Par la règle du $\gamma$, la réaction spontanée a lieu entre $\ce{Cu^2+}$ (oxydant le plus fort) et $\ce{Zn}$ (réducteur le plus fort) :
\begin{align*}
\ce{Cu^2+ + 2e- &-> Cu} \quad \text{(réduction)}\\
\ce{Zn &-> Zn^2+ + 2e-} \quad \text{(oxydation)}
\end{align*}
\[ \ce{Zn(s) + Cu^2+(aq) -> Zn^2+(aq) + Cu(s)} \]
Le procédé du technicien est donc valable.

\textbf{2. Réactif limitant.}

Quantités de matière initiales :
\[ n(\ce{Zn}) = \frac{m}{M} = \frac{6,54}{65,4} = \SI{0,100}{mol} \]
\[ n(\ce{Cu^2+}) = C \times V = 0,20 \times 0,500 = \SI{0,100}{mol} \]

Tableau d'avancement :
\begin{center}
\begin{tabularx}{\linewidth}{|l|X|X|X|X|}
\hline
\rowcolor{popblueL} État & $\ce{Zn}$ & $\ce{Cu^2+}$ & $\ce{Zn^2+}$ & $\ce{Cu}$ \\
\hline
Initial (mol) & 0,100 & 0,100 & 0 & 0 \\
\hline
En cours (mol) & $0,100 - x$ & $0,100 - x$ & $x$ & $x$ \\
\hline
Final (mol) & $0,100 - x_{max}$ & $0,100 - x_{max}$ & $x_{max}$ & $x_{max}$ \\
\hline
\end{tabularx}
\end{center}

Les deux réactifs donnent la même valeur $x_{max} = \SI{0,100}{mol}$ : le mélange est \textbf{stœchiométrique}, il n'y a pas de réactif en excès.

\textbf{3. Masse de cuivre déposée.}
\[ n(\ce{Cu}) = x_{max} = \SI{0,100}{mol} \]
\[ m(\ce{Cu}) = n \times M = 0,100 \times 63,5 = \SI{6,35}{g} \]

\textbf{Conclusion :} la réaction est totale, les deux réactifs sont entièrement consommés et il se dépose $\SI{6,35}{g}$ de cuivre métallique sur la grenaille de zinc.
\end{exoresolu}

\courssub{3. Action des ions $\ce{H3O+}$ sur les métaux : le couple $\ce{H3O+}/\ce{H2}$}

\begin{definition}[Couple $\ce{H3O+}/\ce{H2}$]
En milieu acide aqueux, le couple de référence est $\ce{H3O+}/\ce{H2}$ (souvent noté $\ce{H+}/\ce{H2}$ par simplification). Sa demi-équation de réduction est :
\[ \ce{2 H3O+ + 2 e- <=> H2 + 2 H2O} \]
L'oxydant est l'ion hydronium $\ce{H3O+}$ (ou proton $\ce{H+}$), le réducteur est le dihydrogène $\ce{H2}$ (gaz).
\end{definition}

\begin{methode}[Étudier l'action des ions $\ce{H3O+}$ sur un métal (couple $\ce{H3O+}/\ce{H2}$)]
Pour savoir si un métal $\ce{M}$ est attaqué par un acide ($\ce{HCl}$, $\ce{H2SO4}$ dilué) :
\begin{enumerate}
\item Repérer les deux couples en présence : $\ce{H3O+}/\ce{H2}$ et $\ce{M^{n+}}/\ce{M}$.
\item Localiser le métal par rapport à l'hydrogène dans la classification (réducteur fort en bas) :
\begin{itemize}
\item métal placé \textbf{au-dessous} de $\ce{H3O+}/\ce{H2}$ ($\ce{Zn}$, $\ce{Fe}$, $\ce{Al}$, $\ce{Mg}$...) : il réduit $\ce{H3O+}$, il y a \textbf{dégagement de dihydrogène} et dissolution du métal ;
\item métal placé \textbf{au-dessus} de $\ce{H3O+}/\ce{H2}$ ($\ce{Cu}$, $\ce{Ag}$, $\ce{Au}$...) : \textbf{aucune réaction} avec les acides non oxydants.
\end{itemize}
\item Écrire les demi-équations : $\ce{2H3O+ + 2e- -> H2 + 2H2O}$ et $\ce{M -> M^{n+} + n e-}$, puis équilibrer les électrons.
\item Caractériser le gaz formé : le dihydrogène produit une \textbf{détonation} (« aboiement ») à l'approche d'une flamme.
\end{enumerate}
\textbf{Réflexe :} seuls les métaux placés \textbf{au-dessous de l'hydrogène} (réducteurs plus forts que $\ce{H2}$) dégagent $\ce{H2}$ avec les acides non oxydants.
\end{methode}

\begin{exoresolu}[Zinc, fer et cuivre face à l'acide chlorhydrique]
Dans trois tubes à essais, on verse $\SI{10}{mL}$ d'acide chlorhydrique à $\SI{1,0}{mol.L^{-1}}$ et on introduit respectivement de la grenaille de zinc, de la limaille de fer et de la tournure de cuivre.
\begin{enumerate}
\item Décrire ce qu'on observe dans chaque tube et identifier le gaz éventuellement dégagé.
\item Écrire l'équation-bilan de chaque réaction qui a lieu.
\item Calculer le volume de gaz dégagé (mesuré dans les conditions où $V_m = \SI{24,0}{L.mol^{-1}}$) dans le tube contenant $\SI{1,31}{g}$ de zinc, l'acide étant en excès.
\end{enumerate}
Donnée : $M(\ce{Zn}) = \SI{65,4}{g.mol^{-1}}$. Classification : $\ce{Cu^2+}/\ce{Cu}$ ; $\ce{H3O+}/\ce{H2}$ ; $\ce{Fe^2+}/\ce{Fe}$ ; $\ce{Zn^2+}/\ce{Zn}$ (réducteur croissant vers le bas).
\tcblower
\textbf{1. Observations.}

\begin{itemize}
\item \textbf{Tube au zinc} : le zinc est au-dessous de l'hydrogène, donc réducteur plus fort que $\ce{H2}$ : effervescence vive, dégagement de $\ce{H2}$ (détonation à la flamme), dissolution progressive du zinc.
\item \textbf{Tube au fer} : le fer est aussi au-dessous de l'hydrogène : effervescence (plus lente qu'avec le zinc), dégagement de $\ce{H2}$, la solution verdit légèrement (apparition des ions $\ce{Fe^2+}$).
\item \textbf{Tube au cuivre} : le cuivre est au-dessus de l'hydrogène, réducteur plus faible que $\ce{H2}$ : \textbf{aucune réaction observable}.
\end{itemize}

\textbf{2. Équations-bilans.}

Tube au zinc : $\ce{Zn -> Zn^2+ + 2e-}$ et $\ce{2H3O+ + 2e- -> H2 + 2H2O}$, d'où :
\[ \ce{Zn(s) + 2H3O+(aq) -> Zn^2+(aq) + H2(g) + 2H2O(l)} \]
Tube au fer :
\[ \ce{Fe(s) + 2H3O+(aq) -> Fe^2+(aq) + H2(g) + 2H2O(l)} \]

\textbf{3. Volume de dihydrogène.}

L'acide est en excès, le zinc est le réactif limitant :
\[ n(\ce{Zn}) = \frac{1,31}{65,4} = \SI{2,00 \times 10^{-2}}{mol} \]
D'après l'équation, $n(\ce{H2}) = n(\ce{Zn}) = \SI{2,00 \times 10^{-2}}{mol}$.
\[ V(\ce{H2}) = n \times V_m = 2,00 \times 10^{-2} \times 24,0 = \SI{0,480}{L} = \SI{480}{mL} \]

\textbf{Conclusion :} le zinc et le fer sont attaqués par l'acide chlorhydrique avec dégagement de dihydrogène, contrairement au cuivre ; le tube au zinc dégage $\SI{480}{mL}$ de $\ce{H2}$.
\end{exoresolu}

\begin{savaistu}[Le dihydrogène, énergie propre et industrie camerounaise]
La réaction des métaux avec les acides pour produire du $\ce{H2}$ est à la base de la production industrielle d'hydrogène (vapeur + méthane $\ce{CH4}$ à haute température chez \textbf{SONARA} ou engrais). Au laboratoire, c'est une méthode classique de préparation du gaz. Le $\ce{H2}$ est un vecteur d'énergie prometteur (pile à combustible, voiture $\ce{H2}$) car sa combustion ne donne que de l'eau : $\ce{2H2 + O2 -> 2H2O}$. Au Cameroun, la corrosion des toitures en tôle (fer) par l'acidité des pluies (gaz $\ce{CO2}$, $\ce{SO2}$) libère lentement de l'hydrogène et forme de la rouille ($\ce{Fe2O3}$), d'où l'importance des peintures protectrices ou de la galvanisation (zinc).
\end{savaistu}

\courssub{4. Construction de la classification électrochimique qualitative}

\begin{methode}[Construire une classification électrochimique qualitative à partir d'expériences]
Pour classer des couples à partir d'observations expérimentales :
\begin{enumerate}
\item Pour chaque expérience « métal A + ion $\ce{B^{n+}}$ », noter si une réaction a lieu (dépôt, coloration, effervescence) ou non.
\item Traduire chaque observation en inégalité de pouvoir réducteur :
\begin{itemize}
\item si A réduit $\ce{B^{n+}}$ (réaction observée) : A est réducteur \textbf{plus fort} que B ;
\item si rien ne se passe : A est réducteur \textbf{moins fort} que B.
\end{itemize}
\item Utiliser les expériences avec les acides pour placer le couple $\ce{H3O+}/\ce{H2}$ : les métaux attaqués par $\ce{H3O+}$ sont au-dessous de l'hydrogène, les autres au-dessus.
\item Croiser toutes les inégalités et ranger les couples par pouvoir réducteur croissant (ou décroissant) en respectant une convention claire et en l'indiquant explicitement.
\item Vérifier la cohérence : toute réaction observée doit correspondre à la règle du $\gamma$ dans la classification obtenue.
\end{enumerate}
\end{methode}

\begin{exoresolu}[Établir une classification à partir de quatre expériences]
On réalise les expériences suivantes :
\begin{itemize}
\item Exp. 1 : une lame de zinc plongée dans une solution de $\ce{CuSO4}$ se recouvre de cuivre.
\item Exp. 2 : une lame de fer plongée dans une solution de $\ce{ZnSO4}$ : aucune réaction.
\item Exp. 3 : une lame de fer plongée dans l'acide chlorhydrique dilué : dégagement gazeux.
\item Exp. 4 : une lame de cuivre plongée dans l'acide chlorhydrique dilué : aucune réaction.
\end{itemize}
\begin{enumerate}
\item Interpréter chaque expérience par une inégalité de pouvoir réducteur.
\item Classer les couples $\ce{Cu^2+}/\ce{Cu}$, $\ce{H3O+}/\ce{H2}$, $\ce{Fe^2+}/\ce{Fe}$, $\ce{Zn^2+}/\ce{Zn}$ par pouvoir réducteur croissant du métal.
\end{enumerate}
\tcblower
\textbf{1. Interprétations.}

\textbf{Exp. 1 :} le zinc réduit les ions $\ce{Cu^2+}$ ($\ce{Zn + Cu^2+ -> Zn^2+ + Cu}$), donc $\ce{Zn}$ est un réducteur plus fort que $\ce{Cu}$ : $\ce{Zn} > \ce{Cu}$.

\textbf{Exp. 2 :} le fer ne réduit pas les ions $\ce{Zn^2+}$, donc $\ce{Fe}$ est un réducteur moins fort que $\ce{Zn}$ : $\ce{Zn} > \ce{Fe}$.

\textbf{Exp. 3 :} le fer réduit les ions $\ce{H3O+}$ en $\ce{H2}$, donc $\ce{Fe}$ est un réducteur plus fort que $\ce{H2}$ : $\ce{Fe} > \ce{H2}$.

\textbf{Exp. 4 :} le cuivre ne réduit pas $\ce{H3O+}$, donc $\ce{Cu}$ est un réducteur moins fort que $\ce{H2}$ : $\ce{H2} > \ce{Cu}$.

\textbf{2. Classification.}

En croisant les inégalités : $\ce{Zn} > \ce{Fe} > \ce{H2} > \ce{Cu}$.

Classification par pouvoir réducteur \textbf{croissant} (du métal le moins réducteur au plus réducteur) :
\begin{center}
\begin{tabularx}{0,9\linewidth}{|X|X|}
\hline
\rowcolor{popblueL} Couple Ox/Red & Pouvoir réducteur du métal \\
\hline
$\ce{Cu^2+}/\ce{Cu}$ & le plus faible \\
\hline
$\ce{H3O+}/\ce{H2}$ & intermédiaire (référence) \\
\hline
$\ce{Fe^2+}/\ce{Fe}$ & fort \\
\hline
$\ce{Zn^2+}/\ce{Zn}$ & le plus fort \\
\hline
\end{tabularx}
\end{center}

\textbf{Conclusion :} le couple $\ce{H3O+}/\ce{H2}$ se place entre les couples $\ce{Cu^2+}/\ce{Cu}$ et $\ce{Fe^2+}/\ce{Fe}$ ; le zinc est le réducteur le plus fort des quatre.
\end{exoresolu}

\begin{aretenir}[Classification électrochimique de référence (extrait)]
\begin{center}
\begin{tabularx}{\linewidth}{|c|X|c|}
\hline
\rowcolor{popblueL} \textbf{Rang} & \textbf{Couple Ox/Red} & \textbf{Pouvoir réducteur du Red} \\
\hline
1 & $\ce{K+}/\ce{K}$ ; $\ce{Ca^2+}/\ce{Ca}$ ; $\ce{Na+}/\ce{Na}$ ; $\ce{Mg^2+}/\ce{Mg}$ & Très fort (métaux alcalins/alcalino-terreux) \\
\hline
2 & $\ce{Al^3+}/\ce{Al}$ ; $\ce{Zn^2+}/\ce{Zn}$ ; $\ce{Fe^2+}/\ce{Fe}$ ; $\ce{Ni^2+}/\ce{Ni}$ ; $\ce{Sn^2+}/\ce{Sn}$ ; $\ce{Pb^2+}/\ce{Pb}$ & Fort (réduisent $\ce{H3O+}$) \\
\hline
3 & $\ce{2H3O+}/\ce{H2}$ & \textbf{Référence} \\
\hline
4 & $\ce{Cu^2+}/\ce{Cu}$ ; $\ce{Ag+}/\ce{Ag}$ ; $\ce{Au^3+}/\ce{Au}$ & Faible (ne réduisent pas $\ce{H3O+}$) \\
\hline
\end{tabularx}
\end{center}
\textbf{Convention :} dans ce tableau, le pouvoir réducteur \textbf{croît vers le bas}. L'oxydant le plus fort est en haut ($\ce{K+}$), le réducteur le plus fort est en bas ($\ce{Au}$... non, $\ce{K}$). \emph{Attention :} certains tableaux inversent le sens. \textbf{Indiquez toujours votre convention.}
\end{aretenir}

\courssub{5. Utilisation de la classification : prévision des réactions}

\begin{methode}[Prévoir une réaction et choisir les bons matériaux grâce à la classification]
Pour exploiter la classification dans un problème concret (protection, corrosion, choix de matériau) :
\begin{enumerate}
\item Identifier toutes les espèces susceptibles de réagir (métaux, ions, acide).
\item Écrire tous les couples présents et les classer selon la convention adoptée.
\item Appliquer la règle du $\gamma$ : seule la réaction entre l'oxydant le plus fort et le réducteur le plus fort est spontanée.
\item Pour \textbf{protéger} un métal $\ce{M}$ de la corrosion : le mettre en contact avec un métal \textbf{plus réducteur} que lui (anode sacrificielle, galvanisation au zinc) qui s'oxydera à sa place.
\item Pour \textbf{stocker} une solution d'ions $\ce{M^{n+}}$ : choisir un récipient en métal \textbf{moins réducteur} que $\ce{M}$, sinon le récipient est attaqué.
\end{enumerate}
\end{methode}

\begin{exoresolu}[Peut-on stocker une solution de sulfate de cuivre dans une cuve en fer ?]
Une coopérative agricole de l'Ouest camerounais utilise une solution de sulfate de cuivre(II) (bouillie bordelaise) comme fongicide pour les cacaoyers. Le responsable propose de la stocker dans une cuve en fer.
\begin{enumerate}
\item Montrer, à l'aide de la classification, que ce choix est mauvais et écrire l'équation-bilan de la réaction qui se produirait.
\item Proposer deux métaux, parmi $\ce{Cu}$, $\ce{Ag}$, $\ce{Zn}$, $\ce{Fe}$, convenant pour fabriquer la cuve. Justifier.
\item La cuve en fer contient $\SI{200}{L}$ de solution à $\SI{0,10}{mol.L^{-1}}$ en $\ce{Cu^2+}$. Quelle masse de fer serait consommée si toute la solution réagissait ?
\end{enumerate}
Donnée : $M(\ce{Fe}) = \SI{55,8}{g.mol^{-1}}$. Classification (réducteur croissant vers le bas) : $\ce{Cu^2+}/\ce{Cu}$ ; $\ce{Ag+}/\ce{Ag}$ au-dessus ; $\ce{Fe^2+}/\ce{Fe}$ ; $\ce{Zn^2+}/\ce{Zn}$.
\tcblower
\textbf{1. Incompatibilité fer / ions $\ce{Cu^2+}$.}

Le fer est un réducteur plus fort que le cuivre : par la règle du $\gamma$, $\ce{Fe}$ réduit spontanément $\ce{Cu^2+}$ :
\[ \ce{Fe(s) + Cu^2+(aq) -> Fe^2+(aq) + Cu(s)} \]
La cuve serait rongée (corrosion, fuites) et la solution épuisée en ions $\ce{Cu^2+}$ : le choix est mauvais.

\textbf{2. Choix du matériau.}

Il faut un métal réducteur \textbf{moins fort} que le cuivre, c'est-à-dire placé au-dessus du couple $\ce{Cu^2+}/\ce{Cu}$ : le \textbf{cuivre} lui-même (même couple, aucune réaction) ou l'\textbf{argent} (réducteur plus faible que $\ce{Cu}$, incapable de réduire $\ce{Cu^2+}$). Le zinc et le fer, plus réducteurs, seraient attaqués.

\textbf{3. Masse de fer consommée.}
\[ n(\ce{Cu^2+}) = C \times V = 0,10 \times 200 = \SI{20}{mol} \]
D'après l'équation, $n(\ce{Fe}) = n(\ce{Cu^2+}) = \SI{20}{mol}$.
\[ m(\ce{Fe}) = n \times M = 20 \times 55,8 = \SI{1116}{g} \approx \SI{1,12}{kg} \]

\textbf{Conclusion :} la cuve en fer est inadaptée : elle perdrait environ $\SI{1,12}{kg}$ de fer et toute la solution serait détruite ; il faut une cuve en cuivre (ou en matière plastique inerte type PEHD).
\end{exoresolu}

\begin{savaistu}[Corrosion et protection au Cameroun : du Mont Cameroun aux pipelines]
La corrosion du fer (rouille) est une réaction d'oxydoréduction où le fer s'oxyde en $\ce{Fe^2+}$/$\ce{Fe^3+}$ par l'oxygène de l'air et l'eau. Au Cameroun, l'humidité et l'acidité des sols (région du Mont Cameroun, zones volcaniques) accélèrent ce phénomène.
\begin{itemize}
\item \textbf{Galvanisation :} les pylônes électriques (ENEO), les toitures et les gaines de forage sont en acier galvanisé (zinc). Le zinc (réducteur plus fort que le fer) s'oxyde \emph{à la place} du fer : protection sacrificielle.
\item \textbf{Anodes sacrificielles :} les pipelines de \textbf{SONARA} (pétrole) et les coques de navires au port de Douala/Kribi sont protégés par des anodes en magnésium ou aluminium (encore plus réducteurs que le zinc).
\item \textbf{Passivation :} l'aluminium (casseroles, structures) forme une couche d'alumine $\ce{Al2O3}$ imperméable qui le protège, malgré sa position forte dans la classification.
\end{itemize}
\end{savaistu}

\courssub{6. Synthèse et applications au quotidien camerounais}

\begin{aretenir}[Synthèse : les 4 piliers de l'oxydoréduction en Première]
\begin{enumerate}
\item \textbf{Couple Ox/Red :} deux espèces conjuguées différant par $n\,e^-$. Écriture $\ce{Ox + n e- <=> Red}$. Notation $\ce{Ox}/\ce{Red}$.
\item \textbf{Classification :} rangement des couples selon le pouvoir oxydant/réducteur. Convention explicite (ex. : réducteur fort en bas). Place du couple $\ce{H3O+}/\ce{H2}$ comme référence.
\item \textbf{Règle du $\gamma$ :} réaction spontanée = diagonale (Ox fort + Red fort). Prévoir le sens d'une réaction, le réactif limitant, les produits.
\item \textbf{Calculs quantitatifs :} équation-bilan $\rightarrow$ quantités initiales $\rightarrow$ tableau d'avancement $\rightarrow$ $x_{max}$ $\rightarrow$ bilan final (masses, concentrations, volumes gazeux).
\end{enumerate}
\end{aretenir}

\begin{methode}[Exploiter quantitativement une réaction d'oxydoréduction (tableau d'avancement)]
Pour les calculs de quantités de matière en oxydoréduction :
\begin{enumerate}
\item Écrire et équilibrer l'équation-bilan ionique de la réaction (égalité des atomes et des charges).
\item Calculer les quantités initiales : $n = \dfrac{m}{M}$ pour un solide, $n = C \times V$ pour un soluté (attention : $V$ en litres).
\item Dresser le tableau d'avancement et déterminer $x_{max}$ en annulant chaque quantité de réactif : la plus petite valeur trouvée est $x_{max}$.
\item En déduire le bilan final : masses ($m = n \times M$), concentrations ($C = \dfrac{n}{V}$), volumes de gaz ($V = n \times V_m$).
\item Vérifier la cohérence (unités, chiffres significatifs, réactif restant positif ou nul).
\end{enumerate}
\end{methode}

\begin{exoresolu}[Dépôt d'argent sur une lame de cuivre]
On plonge une lame de cuivre de masse $m = \SI{3,18}{g}$ dans $V = \SI{100}{mL}$ d'une solution de nitrate d'argent de concentration $C = \SI{0,50}{mol.L^{-1}}$. Il se forme un dépôt gris d'argent et la solution bleuit.
\begin{enumerate}
\item Écrire l'équation-bilan de la réaction.
\item Déterminer l'avancement maximal et le réactif limitant.
\item Calculer la masse d'argent déposée et la concentration finale en ions $\ce{Cu^2+}$ (volume supposé constant).
\end{enumerate}
Données : $M(\ce{Cu}) = \SI{63,5}{g.mol^{-1}}$ ; $M(\ce{Ag}) = \SI{107,9}{g.mol^{-1}}$.
\tcblower
\textbf{1. Équation-bilan.}

Couples : $\ce{Ag+}/\ce{Ag}$ et $\ce{Cu^2+}/\ce{Cu}$. Le cuivre, réducteur plus fort que l'argent, réduit $\ce{Ag+}$ :
\begin{align*}
\ce{Ag+ + e- &-> Ag} \quad (\times 2)\\
\ce{Cu &-> Cu^2+ + 2e-}
\end{align*}
\[ \ce{Cu(s) + 2Ag+(aq) -> Cu^2+(aq) + 2Ag(s)} \]

\textbf{2. Avancement maximal.}

Quantités initiales :
\[ n(\ce{Cu}) = \frac{3,18}{63,5} = \SI{5,01 \times 10^{-2}}{mol} \qquad n(\ce{Ag+}) = 0,50 \times 0,100 = \SI{5,00 \times 10^{-2}}{mol} \]

Tableau d'avancement :
\begin{center}
\begin{tabularx}{\linewidth}{|l|X|X|X|X|}
\hline
\rowcolor{popblueL} État (mol) & $\ce{Cu}$ & $\ce{Ag+}$ & $\ce{Cu^2+}$ & $\ce{Ag}$ \\
\hline
Initial & 0,0501 & 0,0500 & 0 & 0 \\
\hline
Final & $0,0501 - x_{max}$ & $0,0500 - 2x_{max}$ & $x_{max}$ & $2x_{max}$ \\
\hline
\end{tabularx}
\end{center}

Si $\ce{Cu}$ limitant : $x_{max} = 0,0501$ mol. Si $\ce{Ag+}$ limitant : $x_{max} = \dfrac{0,0500}{2} = 0,0250$ mol.

La plus petite valeur est retenue : $x_{max} = \SI{2,50 \times 10^{-2}}{mol}$ ; le réactif limitant est l'ion $\ce{Ag+}$.

\textbf{3. Bilan final.}

Masse d'argent déposée :
\[ m(\ce{Ag}) = 2x_{max} \times M(\ce{Ag}) = 0,0500 \times 107,9 = \SI{5,40}{g} \]
Concentration finale en $\ce{Cu^2+}$ :
\[ [\ce{Cu^2+}]_f = \frac{x_{max}}{V} = \frac{0,0250}{0,100} = \SI{0,250}{mol.L^{-1}} \]

\textbf{Conclusion :} il se dépose $\SI{5,40}{g}$ d'argent, la solution atteint $[\ce{Cu^2+}] = \SI{0,250}{mol.L^{-1}}$ et il reste $0,0501 - 0,0250 = \SI{2,51 \times 10^{-2}}{mol}$ de cuivre non attaqué.
\end{exoresolu}

\begin{attention}[Les 5 erreurs qui coûtent des points au Probatoire]
\begin{enumerate}
\item \textbf{Inverser oxydant et réducteur dans le couple.} Un couple s'écrit toujours Ox/Red : $\ce{Cu^2+}/\ce{Cu}$ et jamais $\ce{Cu}/\ce{Cu^2+}$. Le métal est le réducteur, l'ion est l'oxydant.
\item \textbf{Placer les électrons du mauvais côté.} Dans la demi-équation écrite en réduction, les électrons sont à gauche, avec l'oxydant : $\ce{Cu^2+ + 2e- <=> Cu}$. Des électrons à droite dans l'écriture du couple est une faute.
\item \textbf{Oublier d'égaliser les électrons avant d'additionner.} Avec $\ce{Ag+}/\ce{Ag}$ et $\ce{Cu^2+}/\ce{Cu}$, il faut multiplier la demi-équation de l'argent par 2 : $\ce{Cu + 2Ag+ -> Cu^2+ + 2Ag}$. Écrire $\ce{Cu + Ag+ -> Cu^2+ + Ag}$ est doublement faux (charge et matière non conservées).
\item \textbf{Croire que tout métal réagit avec les acides.} Le cuivre, l'argent, l'or (au-dessus de $\ce{H3O+}/\ce{H2}$ dans la classification « réducteur fort en bas ») ne sont pas attaqués par les acides non oxydants. Toujours comparer la position du métal à celle de l'hydrogène avant de conclure.
\item \textbf{Confondre le sens de la classification.} Avant d'appliquer la règle du $\gamma$, vérifie la convention du tableau (réducteur fort en haut ou en bas ?). Une classification lue à l'envers donne la prévision exactement inverse. Indique toujours le sens choisi dans ta copie.
\end{enumerate}
\end{attention}

\begin{exercice}[Entraînement : Prévoir et calculer]
\begin{enumerate}
\item On plonge une lame d'argent dans une solution de nitrate de fer(II) $\ce{Fe(NO3)2}$. Prédire ce qui se passe. Justifier par la classification.
\item On souhaite produire du dihydrogène en faisant réagir de l'acide chlorhydrique sur un métal. Parmi $\ce{Cu}$, $\ce{Zn}$, $\ce{Ag}$, $\ce{Fe}$, lesquels conviennent ? Écrire l'équation-bilan pour celui qui réagit le plus vite.
\item \textbf{Problème Probatoire type :} On introduit $\SI{2,00}{g}$ de poudre de zinc dans $\SI{100,0}{mL}$ d'une solution de nitrate d'argent $\ce{AgNO3}$ à $\SI{0,100}{mol.L^{-1}}$.
\begin{enumerate}
\item Écrire l'équation-bilan de la réaction.
\item Faire le tableau d'avancement et déterminer le réactif limitant.
\item Calculer la masse d'argent métallique obtenue et la concentration finale en ions $\ce{Zn^2+}$.
\end{enumerate}
Données : $M(\ce{Zn}) = \SI{65,4}{g.mol^{-1}}$ ; $M(\ce{Ag}) = \SI{107,9}{g.mol^{-1}}$.
\end{enumerate}
\end{exercice}

\begin{corrige}[Exercice n]
\textbf{1.} Couples : $\ce{Ag+}/\ce{Ag}$ et $\ce{Fe^2+}/\ce{Fe}$. Dans la classification, $\ce{Ag}$ est au-dessus de $\ce{Fe}$ (réducteur plus faible). L'argent ne réduit pas $\ce{Fe^2+}$. \textbf{Aucune réaction}.

\textbf{2.} Métaux au-dessous de $\ce{H3O+}/\ce{H2}$ : $\ce{Zn}$ et $\ce{Fe}$. Le zinc réagit plus vite (plus réducteur). Équation : $\ce{Zn + 2H3O+ -> Zn^2+ + H2 + 2H2O}$.

\textbf{3.}
\begin{itemize}
\item Couples : $\ce{Ag+}/\ce{Ag}$, $\ce{Zn^2+}/\ce{Zn}$. Zn plus réducteur que Ag. Réaction : $\ce{Zn + 2Ag+ -> Zn^2+ + 2Ag}$.
\item $n(\ce{Zn}) = 2,00 / 65,4 = \SI{3,06 \times 10^{-2}}{mol}$. $n(\ce{Ag+}) = 0,100 \times 0,100 = \SI{1,00 \times 10^{-2}}{mol}$.
\item Tableau : Initial Zn ($3,06\times10^{-2}$), Ag+ ($1,00\times10^{-2}$). Final Zn ($3,06\times10^{-2}-x$), Ag+ ($1,00\times10^{-2}-2x$).
\item $x_{max}(\ce{Zn}) = 3,06\times10^{-2}$ ; $x_{max}(\ce{Ag+}) = 5,00\times10^{-3}$. Limitant = $\ce{Ag+}$. $x_{max} = \SI{5,00 \times 10^{-3}}{mol}$.
\item $m(\ce{Ag}) = 2x_{max} \times M(\ce{Ag}) = 0,0100 \times 107,9 = \SI{1,08}{g}$.
\item $[\ce{Zn^2+}]_f = x_{max} / 0,100 = \SI{0,0500}{mol.L^{-1}}$.
\end{itemize}
\end{corrige}

\newpage

\coursec{Exercices}

\coursec{Couple oxydant-réducteur et classification électrochimique}

\courssub{1. Couple oxydant-réducteur : définition et écriture}

\begin{definition}[Couple oxydant-réducteur]
Un \cle{couple oxydant-réducteur} est un ensemble de deux espèces chimiques (l'une chargée, l'autre neutre ou différemment chargée) qui se transforment l'une en l'autre par gain ou perte d'\cle{électrons}. On le note sous la forme \ce{Ox}/\ce{Red} (ou \ce{Ox}/\ce{Red} avec une barre oblique), où :
\begin{itemize}
    \item \ce{Ox} est la \cle{forme oxydée} (l'\cle{oxydant}) : elle \textbf{accepte} des électrons.
    \item \ce{Red} est la \cle{forme réduite} (le \cle{réducteur}) : elle \textbf{donne} des électrons.
\end{itemize}
\end{definition}

\begin{propriete}[Demi-équation électronique]
La transformation d'un couple s'écrit sous forme d'une \cle{demi-équation électronique} (ou demi-réaction) équilibrée en atomes et en charges :
\[ \ce{Ox + n e- <=> Red} \]
\begin{itemize}
    \item Sens direct (de gauche à droite) : \cle{réduction} (gain de $n$ électrons par l'oxydant).
    \item Sens inverse (de droite à gauche) : \cle{oxydation} (perte de $n$ électrons par le réducteur).
\end{itemize}
\end{propriete}

\begin{exemplebox}[Écriture des couples usuels]
\begin{center}
\begin{tabularx}{\linewidth}{|l|l|l|l|}
\hline
\rowcolor{popblueL}
\textbf{Couple} & \textbf{Oxydant (Ox)} & \textbf{Réducteur (Red)} & \textbf{Demi-équation électronique} \\
\hline
\ce{Cu^2+}/\ce{Cu} & \ce{Cu^2+} & \ce{Cu} & \ce{Cu^2+ + 2e- <=> Cu} \\
\hline
\ce{Fe^2+}/\ce{Fe} & \ce{Fe^2+} & \ce{Fe} & \ce{Fe^2+ + 2e- <=> Fe} \\
\hline
\ce{Zn^2+}/\ce{Zn} & \ce{Zn^2+} & \ce{Zn} & \ce{Zn^2+ + 2e- <=> Zn} \\
\hline
\ce{Ag+}/\ce{Ag} & \ce{Ag+} & \ce{Ag} & \ce{Ag+ + e- <=> Ag} \\
\hline
\ce{Al^3+}/\ce{Al} & \ce{Al^3+} & \ce{Al} & \ce{Al^3+ + 3e- <=> Al} \\
\hline
\ce{H3O+}/\ce{H2} & \ce{H3O+} & \ce{H2} & \ce{2H3O+ + 2e- <=> H2 + 2H2O} \\
\hline
\ce{MnO4-}/\ce{Mn^2+} & \ce{MnO4-} & \ce{Mn^2+} & \ce{MnO4- + 8H+ + 5e- <=> Mn^2+ + 4H2O} \\
\hline
\ce{Cr2O7^2-}/\ce{Cr^3+} & \ce{Cr2O7^2-} & \ce{Cr^3+} & \ce{Cr2O7^2- + 14H+ + 6e- <=> 2Cr^3+ + 7H2O} \\
\hline
\end{tabularx}
\end{center}
\end{exemplebox}

\begin{attention}[Pièges à éviter]
\begin{itemize}
    \item Ne pas confondre \textbf{oxydant} (espèce qui s'\emph{oxyde} ? NON, qui \emph{oxyde} l'autre, donc qui \textbf{se réduit}) et \textbf{réducteur} (espèce qui \emph{réduit} l'autre, donc qui \textbf{s'oxyde}).
    \item La barre oblique \textbf{/} sépare l'oxydant (à gauche) du réducteur (à droite) : \ce{Ox}/\ce{Red}.
    \item Dans la demi-équation, les électrons sont \textbf{à gauche} (réduction). Pour l'oxydation, on retourne la flèche et les électrons passent à droite.
    \item Vérifier systématiquement l'équilibre des \textbf{charges} (somme des charges à gauche = somme des charges à droite).
\end{itemize}
\end{attention}

\courssub{2. Réaction entre un métal et un ion métallique : couples \ce{Mn+}/\ce{M}}

\begin{experience}[Mise en évidence de la réactivité des métaux vis-à-vis des ions métalliques]
\textbf{Matériel :} Béchers, lames de métaux (Zn, Fe, Cu, Al), solutions de sels (\ce{CuSO4} bleu, \ce{FeSO4} vert pâle, \ce{AgNO3} incolore, \ce{ZnSO4} incolore), papier émeri, pince à éprouvette.

\textbf{Protocole :}
\begin{enumerate}
    \item Décaper chaque lame de métal avec du papier émeri (éliminer la couche d'oxyde passivante).
    \item Plonger une lame de zinc dans \ce{CuSO4}, une lame de fer dans \ce{CuSO4}, un fil de cuivre dans \ce{AgNO3}, une lame de cuivre dans \ce{ZnSO4}, etc.
    \item Observer pendant 5 à 10 min : aspect du métal (dépôt, ternissement), couleur de la solution.
\end{enumerate}

\textbf{Observations types :}
\begin{itemize}
    \item \textbf{Zn dans \ce{CuSO4} :} La lame se recouvre d'un dépôt rouge-orangé (cuivre métallique), la solution bleue se décolore. \textbf{Réaction.}
    \item \textbf{Fe dans \ce{CuSO4} :} Le clou se recouvre de cuivre, solution se décolore. \textbf{Réaction.}
    \item \textbf{Cu dans \ce{AgNO3} :} Le fil se recouvre d'un dépôt gris-noir (argent), solution reste incolore. \textbf{Réaction.}
    \item \textbf{Cu dans \ce{ZnSO4} :} Aucune transformation visible. \textbf{Pas de réaction.}
    \item \textbf{Fe dans \ce{ZnSO4} :} Aucune transformation. \textbf{Pas de réaction.}
\end{itemize}

\textbf{Interprétation :}
Une réaction a lieu si le métal mis en présence (réducteur) est \textbf{plus fort réducteur} que le métal de l'ion en solution (réducteur associé à l'oxydant). Le réducteur le plus fort donne ses électrons à l'oxydant le plus fort.
\end{experience}

\begin{methode}[Écrire l'équation-bilan d'une réaction entre un métal et un ion métallique]
\textbf{Exemple :} Lame de fer dans solution de \ce{Cu^2+}.
\begin{enumerate}
    \item Identifier les deux couples en présence : \ce{Fe^2+}/\ce{Fe} et \ce{Cu^2+}/\ce{Cu}.
    \item Écrire les deux demi-équations \textbf{de réduction} :
    \begin{align*}
        \ce{Fe^2+ + 2e- <=> Fe} \\
        \ce{Cu^2+ + 2e- <=> Cu}
    \end{align*}
    \item Déterminer le sens spontané : le réducteur le plus fort (celui qui donne le plus facilement ses électrons) s'oxyde. Ici, le fer réduit les ions \ce{Cu^2+} $\rightarrow$ \ce{Fe} s'oxyde (oxydation), \ce{Cu^2+} se réduit (réduction).
    \item Retourner la demi-équation de l'oxydation (électrons à droite) :
    \ce{Fe <=> Fe^2+ + 2e-}
    \item Additionner les deux demi-équations (électrons s'annulent) :
    \[ \ce{Fe + Cu^2+ -> Fe^2+ + Cu} \]
    \item Vérifier : atomes et charges équilibrés.
\end{enumerate}
\end{methode}

\courssub{3. Action des ions \ce{H3O+} sur les métaux : le couple \ce{H3O+}/\ce{H2}}

\begin{experience}[Réaction des métaux avec une solution acide diluée]
\textbf{Matériel :} Tubes à essai, acide chlorhydrique dilué (\ce{HCl} \SI{0,5}{mol.L^{-1}}), acide sulfurique dilué (\ce{H2SO4} \SI{0,5}{mol.L^{-1}}), grenaille de zinc, limaille de fer, fil de cuivre, lame d'aluminium, allumette, entonnoir à gaz ou tube retourné.

\textbf{Protocole :}
\begin{enumerate}
    \item Introduire \SI{2}{mL} d'acide dans chaque tube.
    \item Ajouter le métal (bien décapé pour l'Al).
    \item Observer : dégagement gazeux (vitesse, abondance), échauffement, aspect du métal.
    \item Recueillir le gaz et tester à la flamme (détonation caractéristique = dihydrogène \ce{H2}).
\end{enumerate}

\textbf{Observations et interprétations :}
\begin{center}
\begin{tabularx}{\linewidth}{|l|X|X|}
\hline
\rowcolor{popblueL}
\textbf{Métal} & \textbf{Observation} & \textbf{Interprétation} \\
\hline
Zinc (\ce{Zn}) & Dégagement gazeux \textbf{rapide et abondant}, solution reste incolore. & \ce{Zn} réduit \ce{H3O+} : \ce{Zn + 2H3O+ -> Zn^2+ + H2 + 2H2O}. Couple \ce{Zn^2+}/\ce{Zn} au-dessus de \ce{H3O+}/\ce{H2}. \\
\hline
Fer (\ce{Fe}) & Dégagement gazeux \textbf{lent}, solution verdit (\ce{Fe^2+}). & \ce{Fe} réduit \ce{H3O+} : \ce{Fe + 2H3O+ -> Fe^2+ + H2 + 2H2O}. Couple \ce{Fe^2+}/\ce{Fe} au-dessus de \ce{H3O+}/\ce{H2}. \\
\hline
Cuivre (\ce{Cu}) & \textbf{Aucune réaction} (même à chaud). & \ce{Cu} ne réduit pas \ce{H3O+}. Couple \ce{Cu^2+}/\ce{Cu} en-dessous de \ce{H3O+}/\ce{H2}. \\
\hline
Argent (\ce{Ag}) & \textbf{Aucune réaction}. & \ce{Ag} ne réduit pas \ce{H3O+}. Couple \ce{Ag+}/\ce{Ag} en-dessous de \ce{H3O+}/\ce{H2}. \\
\hline
Aluminium (\ce{Al}) & Pas de réaction immédiate (couche \ce{Al2O3} passivante), puis réaction violente si décapé. & \ce{Al} est un très fort réducteur (au-dessus de \ce{H3O+}/\ce{H2}), mais passivé par l'oxyde. \\
\hline
\end{tabularx}
\end{center}
\end{experience}

\begin{propriete}[Demi-équation du couple \ce{H3O+}/\ce{H2} en milieu aqueux]
\[ \ce{2H3O+ + 2e- <=> H2 + 2H2O} \]
\begin{itemize}
    \item \ce{H3O+} (ion oxonium) est l'\textbf{oxydant}.
    \item \ce{H2} (dihydrogène) est le \textbf{réducteur}.
    \item C'est le \cle{couple de référence} de la classification électrochimique.
\end{itemize}
\end{propriete}

\begin{savaistu}[Le couple \ce{H3O+}/\ce{H2} au quotidien : vinaigre et bissap]
Au Cameroun, le \emph{vinaigre de palme} ou le \emph{jus de bissap} (hibiscus) sont acides (pH $\approx$ 3-4). Ils contiennent des ions \ce{H3O+}. Si l'on y plonge un métal comme le zinc (galvanisation) ou le fer (casserole non émaillée), une réaction lente se produit : dégagement d'\ce{H2} et formation de sels métalliques (zincates, ferreux). C'est pour cela que les ustensiles en aluminium ou en fer émaillé sont préférés pour la cuisine acide. L'argent et le cuivre, eux, ne réagissent pas : ils sont "nobles" vis-à-vis de l'acide.
\end{savaistu}

\courssub{4. Construction de la classification électrochimique qualitative}

\begin{definition}[Classification électrochimique]
C'est un \cle{tableau} qui range les couples oxydant-réducteur \textbf{par ordre croissant de force d'oxydant} (ou décroissant de force de réducteur).
\begin{itemize}
    \item \textbf{En haut :} Les \textbf{réducteurs les plus forts} (donnent facilement leurs électrons). Les oxydants associés sont les plus faibles.
    \item \textbf{En bas :} Les \textbf{oxydants les plus forts} (acceptent facilement les électrons). Les réducteurs associés sont les plus faibles.
    \item Le couple \ce{H3O+}/\ce{H2} est placé au \textbf{milieu} (référence).
\end{itemize}
\end{definition}

\begin{methode}[Construire une classification à partir d'expériences]
\begin{enumerate}
    \item Lister tous les couples mis en jeu dans les expériences.
    \item Pour chaque expérience où une réaction a lieu : identifier le réducteur qui réagit (s'oxyde) et l'oxydant qui réagit (se réduit).
    \item Appliquer la règle : \textbf{Le réducteur qui réagit est PLUS FORT que le réducteur de l'autre couple.} Il se place \textbf{AU-DESSUS} dans le tableau.
    \item Pour les expériences sans réaction : le réducteur mis en présence est \textbf{PLUS FAIBLE} que le réducteur de l'oxydant en solution. Il se place \textbf{EN-DESSOUS}.
    \item Classer tous les couples du plus fort réducteur (haut) au plus faible (bas).
\end{enumerate}
\end{methode}

\begin{aretenir}[Règle de prévision (Règle du "diagonal")]
Une réaction d'oxydoréduction est spontanée (a lieu) si l'on met en présence :
\begin{center}
\fbox{\textbf{Un réducteur placé AU-DESSUS d'un oxydant dans la classification.}}
\end{center}
Graphiquement : on trace une flèche \textbf{descendante} du réducteur vers l'oxydant.
\begin{itemize}
    \item $\mathrm{Red}_1$ (haut) + $\mathrm{Ox}_2$ (bas) $\rightarrow$ $\mathrm{Ox}_1$ + $\mathrm{Red}_2$ \quad \textbf{Réaction OUI}.
    \item $\mathrm{Red}_2$ (bas) + $\mathrm{Ox}_1$ (haut) $\rightarrow$ Pas de réaction \quad \textbf{Réaction NON}.
\end{itemize}
\end{aretenir}

\courssub{5. Utilisation de la classification : prévision des réactions}

\begin{exemplebox}[Application de la règle de prévision]
Classification de référence (simplifiée) :
\ce{K+}/\ce{K} ; \ce{Ca^2+}/\ce{Ca} ; \ce{Mg^2+}/\ce{Mg} ; \ce{Al^3+}/\ce{Al} ; \ce{Zn^2+}/\ce{Zn} ; \ce{Fe^2+}/\ce{Fe} ; \ce{H3O+}/\ce{H2} ; \ce{Cu^2+}/\ce{Cu} ; \ce{Ag+}/\ce{Ag} ; \ce{Au^3+}/\ce{Au}.

\begin{enumerate}
    \item \textbf{Zn + Cu²⁺ :} \ce{Zn} (haut) au-dessus de \ce{Cu^2+} (bas) $\rightarrow$ \textbf{OUI}. \ce{Zn + Cu^2+ -> Zn^2+ + Cu}.
    \item \textbf{Cu + 2H₃O⁺ :} \ce{Cu} (bas) en-dessous de \ce{H3O+} (haut) $\rightarrow$ \textbf{NON}.
    \item \textbf{Fe + 2Ag⁺ :} \ce{Fe} (haut) au-dessus de \ce{Ag+} (bas) $\rightarrow$ \textbf{OUI}. \ce{Fe + 2Ag+ -> Fe^2+ + 2Ag}.
    \item \textbf{Al + 3H₃O⁺ :} \ce{Al} (très haut) au-dessus de \ce{H3O+} $\rightarrow$ \textbf{OUI}. \ce{2Al + 6H3O+ -> 2Al^3+ + 3H2 + 6H2O}.
\end{enumerate}
\end{exemplebox}

\begin{attention}[Cas particuliers et pièges du Probatoire]
\begin{itemize}
    \item \textbf{Acides non oxydants (\ce{HCl}, \ce{H2SO4} dilué) :} Seul le couple \ce{H3O+}/\ce{H2} joue. Les métaux au-dessus de \ce{H3O+}/\ce{H2} dégagent \ce{H2}. Ceux en dessous ne réagissent pas.
    \item \textbf{Acides oxydants (\ce{HNO3}, \ce{H2SO4} concentré à chaud) :} L'oxydant n'est plus \ce{H3O+} mais \ce{NO3-} ou \ce{SO4^2-} (couples bien plus forts, placés très bas). \textbf{Même le cuivre et l'argent réagissent !} (Produits : \ce{NO2}, \ce{NO}, \ce{SO2}, pas \ce{H2}).
    \item \textbf{Passivation :} L'aluminium, le chrome, le fer (dans \ce{HNO3} concentré) forment une couche d'oxyde protectrice qui empêche la réaction malgré une position favorable dans la classification.
    \item \textbf{Équilibrage :} Toujours vérifier l'équilibre des charges et des atomes (surtout O et H en milieu acide/basique) dans l'équation-bilan finale.
\end{itemize}
\end{attention}

\courssub{6. Synthèse et applications au quotidien camerounais}

\begin{savaistu}[Applications industrielles et domestiques au Cameroun]
\begin{itemize}
    \item \textbf{Galvanisation (Douala, Yaoundé) :} Tôles en fer recouvertes de zinc (\cle{galvanisation à chaud}). Le zinc (au-dessus du fer) se sacrifie pour protéger le fer de la corrosion (rouille) sous le climat tropical humide. Voir exercice "Galvanisation des tôles".
    \item \textbf{Protection cathodique (Pipeline Tchad-Cameroun, SONARA) :} Anodes sacrificielles en magnésium ou zinc (très forts réducteurs) enterrées le long des canalisations en acier. Elles s'oxydent à la place du fer.
    \item \textbf{Piles et batteries :} Pile Leclanché (Zn/\ce{MnO2}), pile alcaline, accumulateurs plomb-acide (voitures, onduleurs). Le principe est toujours : couple réducteur fort (anode -) / couple oxydant fort (cathode +).
    \item \textbf{Métallurgie extractive :} Réduction des minerais (oxydes) par le coke (C/\ce{CO}) ou l'aluminium (thermite) dans les hauts-fourneaux (projet Mbalam/Nabeba pour le fer).
    \item \textbf{Conservation alimentaire :} Boîtes de conserve (fer étamé). L'étain (\ce{Sn^2+}/\ce{Sn} au-dessus de \ce{Fe^2+}/\ce{Fe} ? Non, \ce{Sn} est légèrement au-dessus, mais surtout inerte chimiquement). Si rayure : le fer se sacrifie pour l'étain ? Vérification : \ce{Fe^2+}/\ce{Fe} (-0,44 V) vs \ce{Sn^2+}/\ce{Sn} (-0,14 V). Fe est plus fort réducteur. Si rayure, le fer s'oxyde (protège l'étain). C'est l'inverse de la galvanisation !
    \item \textbf{Pharmacopée traditionnelle :} Préparation de certains remèdes par décoction acide (libération d'ions métalliques \ce{Fe^2+}, \ce{Zn^2+}, \ce{Cu^2+} depuis les casseroles ou mortiers).
\end{itemize}
\end{savaistu}

\begin{aretenir}[Résumé des compétences clés pour le Probatoire]
\begin{enumerate}
    \item Écrire correctement un couple \ce{Ox}/\ce{Red} et sa demi-équation électronique (équilibrée atomes + charges).
    \item Identifier l'oxydant et le réducteur dans un couple.
    \item Construire une classification à partir de résultats expérimentaux (réaction / non-réaction).
    \item Placer le couple \ce{H3O+}/\ce{H2} comme référence.
    \item Prévoir si une réaction a lieu entre un réducteur et un oxydant (règle de la diagonale descendante).
    \item Écrire l'équation-bilan de la réaction prévue (tableau d'avancement si calculs quantitatifs).
    \item Calculer les quantités de matière, réactif limitant, masse de dépôt, volume de gaz, concentration finale.
\end{enumerate}
\end{aretenir}

\courssub{Je vérifie mes connaissances}

\begin{exercice}[QCM : notions fondamentales]
Pour chaque question, choisir la (ou les) bonne(s) réponse(s) et justifier brièvement.
\begin{enumerate}
    \item Un couple oxydant-réducteur est :
    \begin{enumerate}
        \item un couple d'espèces qui échangent des protons \ce{H+} ;
        \item un couple d'espèces qui échangent des électrons ;
        \item un couple formé d'un acide et d'une base.
    \end{enumerate}
    \item Dans le couple \ce{Cu^2+}/\ce{Cu}, l'oxydant est :
    \begin{enumerate}
        \item \ce{Cu^2+} ; \item \ce{Cu} ; \item \ce{H2O}.
    \end{enumerate}
    \item La demi-équation électronique du couple \ce{Zn^2+}/\ce{Zn} s'écrit (sens de la réduction) :
    \begin{enumerate}
        \item \ce{Zn -> Zn^2+ + 2e-} ;
        \item \ce{Zn^2+ + 2e- -> Zn} ;
        \item \ce{Zn^2+ -> Zn + 2e-}.
    \end{enumerate}
    \item Le couple de référence de la classification électrochimique est :
    \begin{enumerate}
        \item \ce{Cu^2+}/\ce{Cu} ; \item \ce{H3O+}/\ce{H2} ; \item \ce{Ag+}/\ce{Ag}.
    \end{enumerate}
\end{enumerate}
\end{exercice}

\begin{exercice}[Vrai ou faux]
Répondre par vrai ou faux en justifiant chaque réponse.
\begin{enumerate}
    \item Dans une réaction d'oxydoréduction, le réducteur gagne des électrons.
    \item Le fer réduit les ions \ce{Cu^2+} : le couple \ce{Fe^2+}/\ce{Fe} est donc placé au-dessus du couple \ce{Cu^2+}/\ce{Cu} dans la classification électrochimique.
    \item Le cuivre est attaqué par l'acide chlorhydrique dilué avec dégagement de dihydrogène.
    \item Un métal placé au-dessus du couple \ce{H3O+}/\ce{H2} est attaqué par les solutions acides diluées.
    \item L'argent réduit les ions \ce{Cu^2+} en cuivre métallique.
\end{enumerate}
\end{exercice}

\begin{exercice}[Texte à trous]
Compléter le texte suivant avec les mots ou expressions : \emph{oxydant}, \emph{réducteur}, \emph{électrons}, \emph{oxydation}, \emph{réduction}, \emph{\ce{H3O+}/\ce{H2}}, \emph{plus fort}.

\medskip
Un couple oxydant-réducteur est formé de deux espèces chimiques qui se transforment l'une en l'autre par transfert d'\ldots. La forme oxydée est l'\ldots, la forme réduite est le \ldots. La perte d'électrons est une \ldots ; le gain d'électrons est une \ldots. Dans la classification électrochimique, les couples sont classés du réducteur le \ldots (en haut) vers l'oxydant le \ldots (en bas). Le couple de référence est le couple \ldots.
\end{exercice}

\begin{exercice}[Définitions et demi-équations]
\begin{enumerate}
    \item Définir : couple oxydant-réducteur ; oxydant ; réducteur.
    \item Écrire les demi-équations électroniques des couples suivants : \ce{Fe^2+}/\ce{Fe} ; \ce{Ag+}/\ce{Ag} ; \ce{Al^3+}/\ce{Al} ; \ce{H3O+}/\ce{H2} (en milieu aqueux).
    \item Pour chaque couple, préciser l'oxydant et le réducteur.
\end{enumerate}
\end{exercice}

\courssub{Je m'entraîne}

\begin{exercice}[Lame de zinc dans le sulfate de cuivre]
Au laboratoire du lycée, on plonge une lame de zinc bien décapée dans un bécher contenant une solution bleue de sulfate de cuivre \ce{CuSO4}.
\begin{enumerate}
    \item Décrire les observations attendues au bout de quelques minutes (aspect de la lame, couleur de la solution).
    \item Écrire les demi-équations électroniques mises en jeu.
    \item En déduire l'équation-bilan de la réaction.
    \item Conclure sur la position relative des couples \ce{Zn^2+}/\ce{Zn} et \ce{Cu^2+}/\ce{Cu} dans la classification électrochimique.
\end{enumerate}
\end{exercice}

\begin{exercice}[Construction d'une classification]
On réalise les quatre expériences suivantes :
\begin{itemize}
    \item Expérience 1 : une lame de fer plongée dans une solution de \ce{CuSO4} se recouvre de cuivre ;
    \item Expérience 2 : un fil de cuivre plongé dans une solution de nitrate d'argent \ce{AgNO3} se recouvre d'un dépôt gris d'argent ;
    \item Expérience 3 : de la grenaille de zinc versée dans une solution d'acide chlorhydrique dilué produit un dégagement gazeux qui détone à la flamme ;
    \item Expérience 4 : un fil de cuivre plongé dans l'acide chlorhydrique dilué ne subit aucune transformation.
\end{itemize}
\begin{enumerate}
    \item Pour chaque expérience, écrire l'équation-bilan de la réaction lorsqu'elle a lieu.
    \item Établir la classification des cinq couples \ce{Fe^2+}/\ce{Fe}, \ce{Cu^2+}/\ce{Cu}, \ce{Ag+}/\ce{Ag}, \ce{Zn^2+}/\ce{Zn} et \ce{H3O+}/\ce{H2}, du réducteur le plus fort au réducteur le plus faible.
    \item Placer précisément le couple \ce{H3O+}/\ce{H2} et justifier sa position.
\end{enumerate}
\end{exercice}

\begin{exercice}[Prévision de réactions]
On donne la classification électrochimique simplifiée (du réducteur le plus fort au plus faible) :
\ce{Zn^2+}/\ce{Zn} ; \ce{Fe^2+}/\ce{Fe} ; \ce{H3O+}/\ce{H2} ; \ce{Cu^2+}/\ce{Cu} ; \ce{Ag+}/\ce{Ag}.

Prévoir si une réaction a lieu dans chacun des cas suivants. Justifier la réponse et écrire l'équation-bilan des réactions possibles :
\begin{enumerate}
    \item \ce{Ag} + \ce{Cu^2+} ;
    \item \ce{Zn} + \ce{Fe^2+} ;
    \item \ce{Cu} + \ce{H3O+} ;
    \item \ce{Fe} + \ce{Ag+}.
\end{enumerate}
\end{exercice}

\begin{exercice}[Le choix du bidon]
Un artisan de Nkongsamba dispose d'un bidon en fer. Peut-il l'utiliser sans risque pour stocker :
\begin{enumerate}
    \item une solution de sulfate de fer(II) \ce{FeSO4} ?
    \item une solution de sulfate de cuivre \ce{CuSO4} ?
    \item de l'acide chlorhydrique dilué \ce{HCl} ?
\end{enumerate}
Justifier chaque réponse à l'aide de la classification électrochimique et écrire les équations-bilans des réactions éventuelles.
\end{exercice}

\courssub{Je me prépare au Probatoire}

\begin{exercice}[Le fer et le sulfate de cuivre]
Un élève de Première C du Lycée de Bafoussam réalise l'expérience suivante : il plonge un clou en fer de masse $m = \SI{5,6}{g}$ dans un volume $V = \SI{250}{mL}$ d'une solution de sulfate de cuivre(II) de concentration $C = \SI{0,5}{mol.L^{-1}}$.

\medskip
\textbf{Données :} $M(\ce{Fe}) = \SI{56}{g.mol^{-1}}$ ; $M(\ce{Cu}) = \SI{63,5}{g.mol^{-1}}$.

\medskip
Classification (réducteur le plus fort en haut) : \ce{Fe^2+}/\ce{Fe} puis \ce{Cu^2+}/\ce{Cu}.

\begin{enumerate}
    \item Expliquer pourquoi une réaction se produit. Écrire les demi-équations électroniques et l'équation-bilan de la réaction.
    \item Calculer les quantités de matière initiales de fer et d'ions \ce{Cu^2+}. Construire le tableau d'avancement et déterminer le réactif limitant.
    \item En déduire la masse de cuivre déposée sur le clou à la fin de la réaction.
    \item Calculer la concentration finale en ions \ce{Fe^2+} dans la solution (on supposera le volume constant).
    \item Décrire deux observations expérimentales qui permettent de suivre l'avancement de cette réaction.
\end{enumerate}
\end{exercice}

\begin{exercice}[La galvanisation des tôles à Douala]
Dans les quartiers de Douala, les toitures en tôle sont exposées à un climat chaud et humide qui favorise la corrosion du fer. Pour protéger les tôles, on les recouvre d'une couche de zinc : c'est la \emph{galvanisation}.

\medskip
On donne la classification électrochimique simplifiée (réducteur le plus fort en haut) :
\ce{Zn^2+}/\ce{Zn} ; \ce{Fe^2+}/\ce{Fe} ; \ce{H3O+}/\ce{H2} ; \ce{Cu^2+}/\ce{Cu}.

\medskip
\textbf{Données :} $M(\ce{Zn}) = \SI{65,4}{g.mol^{-1}}$ ; $M(\ce{Fe}) = \SI{56}{g.mol^{-1}}$ ; $M(\ce{H}) = \SI{1}{g.mol^{-1}}$.

\begin{enumerate}
    \item Une tôle galvanisée est rayée : le fer est mis à nu au contact de l'eau de pluie légèrement acide. En utilisant la classification, expliquer pourquoi c'est le zinc qui est attaqué et non le fer. Écrire l'équation-bilan de la réaction entre le zinc et les ions \ce{H3O+}.
    \item Justifier l'expression : \og le zinc joue le rôle de réducteur sacrificiel \fg.
    \item On fait réagir une masse $m = \SI{13,08}{g}$ de zinc avec un excès d'acide chlorhydrique dilué. Calculer le volume de dihydrogène dégagé, mesuré dans les conditions normales de température et de pression ($V_m = \SI{22,4}{L.mol^{-1}}$).
    \item Un fabricant peu scrupuleux propose des tôles recouvertes de cuivre. Montrer, à l'aide de la classification, que ce choix ne protège pas le fer : écrire l'équation-bilan de la réaction qui se produirait entre le fer et les ions \ce{Cu^2+}.
\end{enumerate}
\end{exercice}

\begin{exercice}[Étude comparative de trois métaux]
Au laboratoire du Lycée Classique de Dschang, on dispose de trois métaux : fer, zinc et cuivre, ainsi que de solutions de sulfate de zinc \ce{ZnSO4}, de sulfate de fer(II) \ce{FeSO4} et d'acide chlorhydrique dilué. On réalise les expériences suivantes :

\begin{center}
\begin{tabularx}{\linewidth}{|c|X|X|}
\hline
\rowcolor{popblueL}
\textbf{N°} & \textbf{Expérience} & \textbf{Observation} \\
\hline
1 & Lame de zinc dans \ce{FeSO4} & Dépôt grisâtre de fer sur le zinc \\
\hline
2 & Lame de fer dans \ce{ZnSO4} & Aucune transformation \\
\hline
3 & Grenaille de zinc dans \ce{HCl} dilué & Dégagement gazeux abondant \\
\hline
4 & Fil de cuivre dans \ce{HCl} dilué & Aucune transformation \\
\hline
5 & Clou en fer dans \ce{HCl} dilué & Dégagement gazeux lent \\
\hline
\end{tabularx}
\end{center}

\textbf{Données :} $M(\ce{Fe}) = \SI{56}{g.mol^{-1}}$ ; $M(\ce{Zn}) = \SI{65,4}{g.mol^{-1}}$ ; $M(\ce{Cu}) = \SI{63,5}{g.mol^{-1}}$ ; $V_m = \SI{22,4}{L.mol^{-1}}$ (CNTP).

\begin{enumerate}
    \item Interpréter les expériences 1 et 2 : écrire l'équation-bilan de la réaction de l'expérience 1 et expliquer l'absence de réaction dans l'expérience 2.
    \item Identifier le gaz dégagé dans les expériences 3 et 5 et proposer un test expérimental pour le caractériser.
    \item Établir la classification des quatre couples \ce{Zn^2+}/\ce{Zn}, \ce{Fe^2+}/\ce{Fe}, \ce{H3O+}/\ce{H2} et \ce{Cu^2+}/\ce{Cu}, du réducteur le plus fort au réducteur le plus faible. Justifier la place du couple \ce{H3O+}/\ce{H2}.
    \item On plonge une lame de fer de masse $m = \SI{2,8}{g}$ dans $V = \SI{200}{mL}$ d'acide chlorhydrique de concentration $C = \SI{0,5}{mol.L^{-1}}$. Déterminer le réactif limitant à l'aide d'un tableau d'avancement, puis calculer le volume de gaz dégagé (CNTP) et la concentration finale en ions \ce{Fe^2+}.
    \item Citer une application industrielle ou domestique camerounaise qui exploite la classification électrochimique.
\end{enumerate}
\end{exercice}



\newpage

\coursec{Corrigés des exercices}

\begin{corrige}[Exercice 1 — QCM : notions fondamentales]
\begin{enumerate}
    \item \textbf{Bonne réponse : (b).} Un couple oxydant-réducteur met en jeu un transfert d'\cle{électrons} entre un oxydant et un réducteur. L'échange de protons \ce{H+} (réponses a et c) caractérise un couple \textbf{acide-base}, ce qui est hors sujet ici.
    \item \textbf{Bonne réponse : (a) \ce{Cu^2+}.} Dans la notation conventionnelle \ce{Ox}/\ce{Red}, l'\textbf{oxydant} (forme oxydée) est écrit à \textbf{gauche} de la barre oblique. \ce{Cu^2+} peut capter 2 électrons : \ce{Cu^2+ + 2e- <=> Cu}.
    \item \textbf{Bonne réponse : (b) \ce{Zn^2+ + 2e- -> Zn}.} Une demi-équation électronique s'écrit par convention dans le sens de la \textbf{réduction} : l'oxydant et les électrons sont à \textbf{gauche} de la flèche. La réponse (a) est une oxydation ; la réponse (c) n'est pas équilibrée en charges ($+2$ à gauche, $-2$ à droite).
    \item \textbf{Bonne réponse : (b) \ce{H3O+}/\ce{H2}.} C'est le \cle{couple de référence} de la classification électrochimique : il sépare les métaux attaqués par les acides dilués non oxydants (au-dessus) des métaux nobles (en dessous).
\end{enumerate}
\textbf{Point de vigilance :} au Probatoire, une justification même brève est souvent exigée au QCM ; une réponse juste non justifiée peut ne rapporter que la moitié des points.
\end{corrige}

\begin{corrige}[Exercice 2 — Vrai ou faux]
\begin{enumerate}
    \item \textbf{Faux.} Le réducteur \textbf{cède} (perd) des électrons : il subit une \textbf{oxydation}. C'est l'oxydant qui gagne des électrons (réduction). Moyen mnémotechnique : \emph{le réducteur réduit l'autre, donc il s'oxyde lui-même}.
    \item \textbf{Vrai.} Si le fer réduit les ions \ce{Cu^2+}, c'est que \ce{Fe} est un réducteur \textbf{plus fort} que \ce{Cu}. Or la classification range les réducteurs du plus fort (en haut) au plus faible (en bas) : le couple \ce{Fe^2+}/\ce{Fe} est donc placé \textbf{au-dessus} du couple \ce{Cu^2+}/\ce{Cu}.
    \item \textbf{Faux.} Le cuivre est placé \textbf{en dessous} du couple \ce{H3O+}/\ce{H2} : c'est un métal \emph{noble}. Il ne peut pas réduire \ce{H3O+} en \ce{H2}. L'acide chlorhydrique dilué (non oxydant) est \textbf{sans action} sur le cuivre. (Attention : l'acide nitrique, lui, attaque le cuivre, mais c'est \ce{NO3-} qui joue le rôle d'oxydant, pas \ce{H3O+}.)
    \item \textbf{Vrai.} C'est une conséquence directe de la règle de prévision : tout métal dont le couple est au-dessus de \ce{H3O+}/\ce{H2} est un réducteur plus fort que \ce{H2} ; il réduit donc \ce{H3O+} avec dégagement de dihydrogène : \ce{M + n H3O+ -> M^{n+} + (n/2) H2 + n H2O}.
    \item \textbf{Faux.} Le couple \ce{Ag+}/\ce{Ag} est placé \textbf{en dessous} de \ce{Cu^2+}/\ce{Cu} : l'argent est un réducteur \textbf{plus faible} que le cuivre. C'est la réaction inverse qui est spontanée : \ce{Cu + 2Ag+ -> Cu^2+ + 2Ag}.
\end{enumerate}
\textbf{Point de vigilance :} dans un vrai/faux, la justification vaut autant que la réponse. Toujours appuyer la réponse sur la position relative des couples dans la classification.
\end{corrige}

\begin{corrige}[Exercice 3 — Texte à trous]
Un couple oxydant-réducteur est formé de deux espèces chimiques qui se transforment l'une en l'autre par transfert d'\textbf{électrons}. La forme oxydée est l'\textbf{oxydant}, la forme réduite est le \textbf{réducteur}. La perte d'électrons est une \textbf{oxydation} ; le gain d'électrons est une \textbf{réduction}. Dans la classification électrochimique, les couples sont classés du réducteur le \textbf{plus fort} (en haut) vers l'oxydant le \textbf{plus fort} (en bas). Le couple de référence est le couple \textbf{\ce{H3O+}/\ce{H2}}.

\medskip
\textbf{Point de vigilance :} ne pas inverser \emph{oxydation} et \emph{réduction} : l'\textbf{ox}ydation est la perte d'électrons (le nombre d'oxydation \emph{augmente}), la \textbf{réduction} est le gain d'électrons (le nombre d'oxydation \emph{diminue}).
\end{corrige}

\begin{corrige}[Exercice 4 — Définitions et demi-équations]
\begin{enumerate}
    \item \textbf{Définitions :}
    \begin{itemize}
        \item \textbf{Couple oxydant-réducteur :} ensemble de deux espèces chimiques, noté \ce{Ox}/\ce{Red}, qui se transforment l'une en l'autre par transfert d'électrons selon \ce{Ox + n e- <=> Red}.
        \item \textbf{Oxydant :} espèce chimique capable de \textbf{capter} un ou plusieurs électrons (forme oxydée du couple).
        \item \textbf{Réducteur :} espèce chimique capable de \textbf{céder} un ou plusieurs électrons (forme réduite du couple).
    \end{itemize}
    \item \textbf{Demi-équations électroniques (sens de la réduction) :}
    \begin{align*}
        \ce{Fe^2+}/\ce{Fe} &: \quad \ce{Fe^2+ + 2e- <=> Fe} \\
        \ce{Ag+}/\ce{Ag} &: \quad \ce{Ag+ + e- <=> Ag} \\
        \ce{Al^3+}/\ce{Al} &: \quad \ce{Al^3+ + 3e- <=> Al} \\
        \ce{H3O+}/\ce{H2} &: \quad \ce{2H3O+ + 2e- <=> H2 + 2H2O}
    \end{align*}
    \textbf{Vérification des charges} (exemple du couple \ce{H3O+}/\ce{H2}) : à gauche $2\times(+1) + 2\times(-1) = 0$ ; à droite $0 + 2\times 0 = 0$. L'équilibre est respecté.
    \item \textbf{Identification oxydant / réducteur :}
    \begin{center}
\begin{tabularx}{\linewidth}{|l|X|X|}
\hline
\rowcolor{popblueL}
\textbf{Couple} & \textbf{Oxydant (capte \ce{e-})} & \textbf{Réducteur (cède \ce{e-})} \\
\hline
\ce{Fe^2+}/\ce{Fe} & \ce{Fe^2+} & \ce{Fe} \\
\hline
\ce{Ag+}/\ce{Ag} & \ce{Ag+} & \ce{Ag} \\
\hline
\ce{Al^3+}/\ce{Al} & \ce{Al^3+} & \ce{Al} \\
\hline
\ce{H3O+}/\ce{H2} & \ce{H3O+} & \ce{H2} \\
\hline
\end{tabularx}
    \end{center}
\end{enumerate}
\textbf{Point de vigilance :} pour le couple \ce{H3O+}/\ce{H2}, l'erreur classique est d'écrire \ce{H3O+ + e- <=> H2} : l'atome d'hydrogène n'est pas équilibré. Il faut \textbf{2 \ce{H3O+}} pour former \textbf{1 \ce{H2}}, et on ajoute \ce{H2O} pour équilibrer l'oxygène.
\end{corrige}

\begin{corrige}[Exercice 5 — Lame de zinc dans le sulfate de cuivre]
\begin{enumerate}
    \item \textbf{Observations attendues :}
    \begin{itemize}
        \item Sur la lame de zinc : apparition d'un \textbf{dépôt rouge-orangé} (cuivre métallique) ; la lame s'amincit progressivement (le zinc se dissout).
        \item Dans la solution : la teinte \textbf{bleue} caractéristique des ions \ce{Cu^2+} \textbf{pâlit puis disparaît} (les ions \ce{Zn^2+} formés sont incolores).
    \end{itemize}
    \item \textbf{Demi-équations électroniques :}
    \begin{itemize}
        \item Le zinc métallique cède des électrons (oxydation) : \ce{Zn -> Zn^2+ + 2e-}
        \item Les ions cuivre(II) captent des électrons (réduction) : \ce{Cu^2+ + 2e- -> Cu}
    \end{itemize}
    \item \textbf{Équation-bilan :} on additionne les deux demi-équations ; les 2 électrons échangés s'éliminent :
    \[ \ce{Zn(s) + Cu^2+(aq) -> Zn^2+(aq) + Cu(s)} \]
    Vérification : atomes équilibrés (1 Zn, 1 Cu de chaque côté) ; charges : $+2$ de chaque côté.
    \item \textbf{Conclusion :} puisque le zinc \textbf{réduit} les ions \ce{Cu^2+}, le zinc est un réducteur \textbf{plus fort} que le cuivre. Le couple \ce{Zn^2+}/\ce{Zn} se place donc \textbf{au-dessus} du couple \ce{Cu^2+}/\ce{Cu} dans la classification électrochimique.
\end{enumerate}
\textbf{Point de vigilance :} ne pas écrire \ce{SO4^2-} dans l'équation-bilan : c'est un ion \textbf{spectateur}, il ne participe pas au transfert d'électrons.
\end{corrige}

\begin{corrige}[Exercice 6 — Construction d'une classification]
\begin{enumerate}
    \item \textbf{Équations-bilans :}
    \begin{itemize}
        \item \textbf{Expérience 1} (réaction) : \ce{Fe(s) + Cu^2+(aq) -> Fe^2+(aq) + Cu(s)}
        \item \textbf{Expérience 2} (réaction) : \ce{Cu(s) + 2Ag+(aq) -> Cu^2+(aq) + 2Ag(s)} \quad (multiplier la demi-équation \ce{Ag+ + e- -> Ag} par 2 pour compenser les 2 électrons cédés par \ce{Cu}).
        \item \textbf{Expérience 3} (réaction, gaz qui détone = \ce{H2}) : \ce{Zn(s) + 2H3O+(aq) -> Zn^2+(aq) + H2(g) + 2H2O(l)}
        \item \textbf{Expérience 4 :} aucune réaction.
    \end{itemize}
    \item \textbf{Raisonnement expérience par expérience :}
    \begin{itemize}
        \item Exp. 1 : \ce{Fe} réduit \ce{Cu^2+} $\Rightarrow$ \ce{Fe} réducteur plus fort que \ce{Cu} $\Rightarrow$ \ce{Fe^2+}/\ce{Fe} \textbf{au-dessus} de \ce{Cu^2+}/\ce{Cu}.
        \item Exp. 2 : \ce{Cu} réduit \ce{Ag+} $\Rightarrow$ \ce{Cu^2+}/\ce{Cu} \textbf{au-dessus} de \ce{Ag+}/\ce{Ag}.
        \item Exp. 3 : \ce{Zn} réduit \ce{H3O+} $\Rightarrow$ \ce{Zn^2+}/\ce{Zn} \textbf{au-dessus} de \ce{H3O+}/\ce{H2}.
        \item Exp. 4 : \ce{Cu} ne réduit pas \ce{H3O+} $\Rightarrow$ \ce{Cu^2+}/\ce{Cu} \textbf{en dessous} de \ce{H3O+}/\ce{H2}.
    \end{itemize}
    \textbf{Classification obtenue (réducteur le plus fort en haut) :}
    \begin{center}
\begin{tabular}{c}
    \ce{Zn^2+}/\ce{Zn} \\
    \ce{Fe^2+}/\ce{Fe} \\
    \ce{H3O+}/\ce{H2} \\
    \ce{Cu^2+}/\ce{Cu} \\
    \ce{Ag+}/\ce{Ag} \\
\end{tabular}
    \end{center}
    \item \textbf{Place du couple \ce{H3O+}/\ce{H2} :} il se situe \textbf{entre \ce{Fe^2+}/\ce{Fe} et \ce{Cu^2+}/\ce{Cu}}. Justification : le zinc (et le fer, de façon analogue) dégage \ce{H2} au contact de l'acide dilué, donc leurs couples sont au-dessus ; le cuivre et l'argent n'y réagissent pas, donc leurs couples sont en dessous. Le couple \ce{H3O+}/\ce{H2} joue ainsi son rôle de \textbf{frontière de référence} entre métaux « attaquables » et métaux « nobles ».
\end{enumerate}
\textbf{Point de vigilance :} une absence de réaction est une \textbf{information} au même titre qu'une réaction : elle permet de placer le couple du métal introduit \emph{en dessous} du couple de l'ion en solution.
\end{corrige}

\begin{corrige}[Exercice 7 — Prévision de réactions]
On applique la règle de prévision : \textbf{une réaction spontanée a lieu entre le réducteur d'un couple placé au-dessus et l'oxydant d'un couple placé en dessous} (diagonale descendante).

Classification rappelée (haut $\rightarrow$ bas) : \ce{Zn^2+}/\ce{Zn} ; \ce{Fe^2+}/\ce{Fe} ; \ce{H3O+}/\ce{H2} ; \ce{Cu^2+}/\ce{Cu} ; \ce{Ag+}/\ce{Ag}.

\begin{enumerate}
    \item \ce{Ag} + \ce{Cu^2+} : le réducteur \ce{Ag} (couple \ce{Ag+}/\ce{Ag}, tout en bas) est \textbf{en dessous} de l'oxydant \ce{Cu^2+}. $\Rightarrow$ \textbf{Aucune réaction.}
    \item \ce{Zn} + \ce{Fe^2+} : \ce{Zn} est \textbf{au-dessus} de \ce{Fe^2+}. $\Rightarrow$ \textbf{Réaction :}
    \[ \ce{Zn(s) + Fe^2+(aq) -> Zn^2+(aq) + Fe(s)} \]
    \item \ce{Cu} + \ce{H3O+} : \ce{Cu} est \textbf{en dessous} de \ce{H3O+}. $\Rightarrow$ \textbf{Aucune réaction} (le cuivre est noble vis-à-vis des acides dilués non oxydants).
    \item \ce{Fe} + \ce{Ag+} : \ce{Fe} est \textbf{au-dessus} de \ce{Ag+}. $\Rightarrow$ \textbf{Réaction :}
    \[ \ce{Fe(s) + 2Ag+(aq) -> Fe^2+(aq) + 2Ag(s)} \]
    Détail de l'équilibrage : \ce{Fe -> Fe^2+ + 2e-} et $2\times(\ce{Ag+ + e- -> Ag})$ ; les 2 électrons s'éliminent.
\end{enumerate}
\textbf{Point de vigilance :} toujours équilibrer les \textbf{électrons échangés} avant d'additionner les demi-équations ; pour \ce{Ag+}/\ce{Ag} (1 électron) face à un métal qui cède 2 électrons, le coefficient 2 sur \ce{Ag+} et \ce{Ag} est obligatoire.
\end{corrige}

\begin{corrige}[Exercice 8 — Le choix du bidon]
Le bidon est en \textbf{fer}, réducteur du couple \ce{Fe^2+}/\ce{Fe}, placé au-dessus de \ce{H3O+}/\ce{H2} et de \ce{Cu^2+}/\ce{Cu}.

\begin{enumerate}
    \item \textbf{Solution de \ce{FeSO4} : OUI, stockage possible.} Le fer métallique est mis en présence des ions \ce{Fe^2+}, c'est-à-dire l'oxydant de \textbf{son propre couple}. Il n'y a pas de transfert spontané d'électrons au sein d'un même couple : \textbf{aucune réaction}. Le bidon n'est pas attaqué.
    \item \textbf{Solution de \ce{CuSO4} : NON, stockage dangereux.} \ce{Fe} est au-dessus de \ce{Cu^2+} : le fer réduit les ions cuivre(II). Le bidon se perce lentement (le fer se dissout) et du cuivre se dépose sur les parois :
    \[ \ce{Fe(s) + Cu^2+(aq) -> Fe^2+(aq) + Cu(s)} \]
    \item \textbf{Acide chlorhydrique dilué : NON, stockage dangereux.} \ce{Fe} est au-dessus de \ce{H3O+} : le fer réduit les ions oxonium. Le bidon est rongé et il se dégage du \textbf{dihydrogène}, gaz inflammable (risque d'explosion en milieu confiné) :
    \[ \ce{Fe(s) + 2H3O+(aq) -> Fe^2+(aq) + H2(g) + 2H2O(l)} \]
\end{enumerate}
\textbf{Conclusion pratique :} l'artisan ne peut stocker sans risque que la solution de \ce{FeSO4}. Pour l'acide ou le sulfate de cuivre, il devra utiliser un récipient en matière plastique ou en verre.
\end{corrige}

\begin{corrige}[Exercice 9 — Le fer et le sulfate de cuivre (type Probatoire)]
\begin{enumerate}
    \item \textbf{Justification et équations :} la classification place \ce{Fe^2+}/\ce{Fe} \textbf{au-dessus} de \ce{Cu^2+}/\ce{Cu} : le fer, réducteur le plus fort, réduit spontanément les ions \ce{Cu^2+}.
    \begin{align*}
        \text{Oxydation : } & \ce{Fe -> Fe^2+ + 2e-} \\
        \text{Réduction : } & \ce{Cu^2+ + 2e- -> Cu}
    \end{align*}
    \textbf{Équation-bilan :} \ce{Fe(s) + Cu^2+(aq) -> Fe^2+(aq) + Cu(s)}.
    \item \textbf{Quantités initiales :}
    \begin{align*}
        n_i(\ce{Fe}) &= \frac{m}{M(\ce{Fe})} = \frac{5,6}{56} = \SI{0,10}{mol} \\
        n_i(\ce{Cu^2+}) &= C \times V = 0,5 \times 0,250 = \SI{0,125}{mol}
    \end{align*}
    \textbf{Tableau d'avancement :}
    \begin{center}
\begin{tabular}{|l|c|c|c|c|}
\hline
\rowcolor{popblueL}
\textbf{Équation} & \ce{Fe} + & \ce{Cu^2+} -> & \ce{Fe^2+} + & \ce{Cu} \\
\hline
État initial ($x=0$) & $0{,}10$ & $0{,}125$ & $0$ & $0$ \\
\hline
En cours ($x$) & $0{,}10 - x$ & $0{,}125 - x$ & $x$ & $x$ \\
\hline
État final ($x_{\max}$) & $0{,}10 - x_{\max}$ & $0{,}125 - x_{\max}$ & $x_{\max}$ & $x_{\max}$ \\
\hline
\end{tabular}
    \end{center}
    \textbf{Réactif limitant :} hypothèse 1 : $0{,}10 - x_{\max} = 0 \Rightarrow x_{\max} = \SI{0,10}{mol}$ ; hypothèse 2 : $0{,}125 - x_{\max} = 0 \Rightarrow x_{\max} = \SI{0,125}{mol}$. On retient la \textbf{plus petite} valeur : $x_{\max} = \SI{0,10}{mol}$. \textbf{Le fer est le réactif limitant} ; il reste $0{,}125 - 0{,}10 = \SI{0,025}{mol}$ d'ions \ce{Cu^2+} en solution.
    \item \textbf{Masse de cuivre déposée :} $n(\ce{Cu}) = x_{\max} = \SI{0,10}{mol}$, d'où :
    \[ m(\ce{Cu}) = n \times M(\ce{Cu}) = 0,10 \times 63,5 = \boxed{\SI{6,35}{g}} \]
    \item \textbf{Concentration finale en \ce{Fe^2+} :} $n_f(\ce{Fe^2+}) = x_{\max} = \SI{0,10}{mol}$ dans $V = \SI{0,250}{L}$ :
    \[ [\ce{Fe^2+}]_f = \frac{0,10}{0,250} = \boxed{\SI{0,40}{mol.L^{-1}}} \]
    \item \textbf{Observations permettant de suivre la réaction :}
    \begin{itemize}
        \item la \textbf{décoloration progressive} de la solution bleue (consommation des ions \ce{Cu^2+}) ;
        \item la \textbf{croissance du dépôt rouge-orangé} de cuivre sur le clou (on peut aussi peser le clou séché : sa masse évolue).
    \end{itemize}
\end{enumerate}
\textbf{Barème indicatif (sur 5 pts) :} justification + équation-bilan (1 pt) ; quantités initiales (1 pt) ; tableau d'avancement + réactif limitant (1 pt) ; masse de cuivre (1 pt) ; concentration finale + observations (1 pt).

\textbf{Points de vigilance :} convertir $V$ en litres ; vérifier que la solution reste bleue pâle à la fin (il reste des \ce{Cu^2+} : le fer est limitant) ; ne pas confondre la masse de cuivre déposée avec la variation de masse du clou.
\end{corrige}

\begin{corrige}[Exercice 10 — La galvanisation des tôles à Douala (type Probatoire)]
\begin{enumerate}
    \item \textbf{Explication :} au niveau de la rayure, zinc et fer sont en contact et baignent dans l'eau de pluie acide (présence d'ions \ce{H3O+}). Or la classification place \ce{Zn^2+}/\ce{Zn} \textbf{au-dessus} de \ce{Fe^2+}/\ce{Fe} : le zinc est un réducteur \textbf{plus fort} que le fer. C'est donc le zinc qui s'oxyde préférentiellement ; le fer, mis à nu, reste intact tant qu'il reste du zinc.
    \textbf{Équation-bilan :}
    \[ \ce{Zn(s) + 2H3O+(aq) -> Zn^2+(aq) + H2(g) + 2H2O(l)} \]
    \item \textbf{« Réducteur sacrificiel » :} le zinc se \emph{sacrifie} : il s'oxyde (se consomme) \textbf{à la place} du fer. Tant qu'il reste du zinc en contact avec le fer, le fer ne peut pas être oxydé : il est protégé. On parle aussi d'\cle{anode sacrificielle} (protection cathodique).
    \item \textbf{Volume de dihydrogène :}
    \begin{align*}
        n(\ce{Zn}) &= \frac{m}{M(\ce{Zn})} = \frac{13,08}{65,4} = \SI{0,200}{mol} \\
        \text{D'après l'équation : } & n(\ce{H2}) = n(\ce{Zn}) = \SI{0,200}{mol} \quad (\text{acide en excès}) \\
        V(\ce{H2}) &= n \times V_m = 0,200 \times 22,4 = \boxed{\SI{4,48}{L}}
    \end{align*}
    \item \textbf{Tôle recouverte de cuivre :} le couple \ce{Cu^2+}/\ce{Cu} est \textbf{en dessous} de \ce{Fe^2+}/\ce{Fe} : le cuivre est un réducteur \emph{plus faible} que le fer. En cas de rayure, c'est le \textbf{fer} qui s'oxyde au contact du cuivre (corrosion galvanique accélérée) ; au contact d'ions \ce{Cu^2+} :
    \[ \ce{Fe(s) + Cu^2+(aq) -> Fe^2+(aq) + Cu(s)} \]
    Le cuivre \textbf{ne protège pas} le fer, il aggrave sa corrosion : ce choix est à rejeter.
\end{enumerate}
\textbf{Barème indicatif (sur 5 pts) :} explication + équation (1,5 pt) ; rôle sacrificiel (0,5 pt) ; calcul de $n(\ce{Zn})$ (1 pt) ; volume de \ce{H2} (1 pt) ; comparaison Cu/Fe + équation (1 pt).

\textbf{Points de vigilance :} bien vérifier le calcul $13,08 / 65,4 = 0,200$ exactement ; préciser « acide en excès » pour justifier $n(\ce{H2}) = n(\ce{Zn})$ ; ne pas écrire que le cuivre « s'oxyde à la place du fer » : c'est l'inverse.
\end{corrige}

\begin{corrige}[Exercice 11 — Étude comparative de trois métaux (type Probatoire)]
\begin{enumerate}
    \item \textbf{Expérience 1 :} un dépôt grisâtre de fer apparaît sur le zinc : le zinc a réduit les ions \ce{Fe^2+} :
    \[ \ce{Zn(s) + Fe^2+(aq) -> Zn^2+(aq) + Fe(s)} \]
    \textbf{Expérience 2 :} aucune transformation : le fer est un réducteur \textbf{plus faible} que le zinc (conséquence de l'expérience 1) ; il ne peut donc pas réduire les ions \ce{Zn^2+}. Pas de réaction.
    \item \textbf{Gaz dégagé :} le \textbf{dihydrogène \ce{H2}}. \textbf{Test :} on approche une flamme (allumette enflammée) de l'ouverture du tube : le gaz produit une \textbf{détonation caractéristique} (petite explosion avec bruit sec). \emph{Sécurité : effectuer le test sur un petit volume de gaz.}
    \item \textbf{Classification :}
    \begin{itemize}
        \item Exp. 1 : \ce{Zn} réduit \ce{Fe^2+} $\Rightarrow$ \ce{Zn^2+}/\ce{Zn} au-dessus de \ce{Fe^2+}/\ce{Fe} ;
        \item Exp. 3 et 5 : \ce{Zn} et \ce{Fe} réduisent \ce{H3O+} $\Rightarrow$ leurs couples sont au-dessus de \ce{H3O+}/\ce{H2} ;
        \item Exp. 4 : \ce{Cu} ne réduit pas \ce{H3O+} $\Rightarrow$ \ce{Cu^2+}/\ce{Cu} en dessous de \ce{H3O+}/\ce{H2}.
    \end{itemize}
    D'où, du réducteur le plus fort au plus faible :
    \[ \ce{Zn^2+}/\ce{Zn} \;>\; \ce{Fe^2+}/\ce{Fe} \;>\; \ce{H3O+}/\ce{H2} \;>\; \ce{Cu^2+}/\ce{Cu} \]
    \textbf{Place de \ce{H3O+}/\ce{H2} :} entre \ce{Fe^2+}/\ce{Fe} et \ce{Cu^2+}/\ce{Cu}, car le fer et le zinc sont attaqués par l'acide dilué (dégagement de \ce{H2}) alors que le cuivre ne l'est pas.
    \item \textbf{Étude quantitative :} équation-bilan : \ce{Fe(s) + 2H3O+(aq) -> Fe^2+(aq) + H2(g) + 2H2O(l)}.
    \begin{align*}
        n_i(\ce{Fe}) &= \frac{2,8}{56} = \SI{0,050}{mol} \\
        n_i(\ce{H3O+}) &= C \times V = 0,5 \times 0,200 = \SI{0,100}{mol}
    \end{align*}
    \textbf{Tableau d'avancement :}
    \begin{center}
\begin{tabular}{|l|c|c|c|c|}
\hline
\rowcolor{popblueL}
\textbf{Équation} & \ce{Fe} + & \ce{2H3O+} -> & \ce{Fe^2+} + & \ce{H2} \\
\hline
État initial & $0{,}050$ & $0{,}100$ & $0$ & $0$ \\
\hline
En cours & $0{,}050 - x$ & $0{,}100 - 2x$ & $x$ & $x$ \\
\hline
État final & $0{,}050 - x_{\max}$ & $0{,}100 - 2x_{\max}$ & $x_{\max}$ & $x_{\max}$ \\
\hline
\end{tabular}
    \end{center}
    \textbf{Réactif limitant :} hypothèse 1 (fer) : $x_{\max} = \SI{0,050}{mol}$ ; hypothèse 2 (\ce{H3O+}) : $x_{\max} = \frac{0,100}{2} = \SI{0,050}{mol}$. Les deux valeurs sont \textbf{égales} : les réactifs sont en \textbf{proportions stœchiométriques} et sont entièrement consommés. $x_{\max} = \SI{0,050}{mol}$.
    \textbf{Volume de gaz (CNTP) :} $n(\ce{H2}) = x_{\max} = \SI{0,050}{mol}$ :
    \[ V(\ce{H2}) = 0,050 \times 22,4 = \boxed{\SI{1,12}{L}} \]
    \textbf{Concentration finale en \ce{Fe^2+} :} $n_f(\ce{Fe^2+}) = \SI{0,050}{mol}$ dans $V = \SI{0,200}{L}$ :
    \[ [\ce{Fe^2+}]_f = \frac{0,050}{0,200} = \boxed{\SI{0,25}{mol.L^{-1}}} \]
    \item \textbf{Application camerounaise :} la \textbf{galvanisation des tôles de toiture} (Douala, Yaoundé) : le zinc, réducteur plus fort que le fer, protège la tôle de la corrosion sous le climat tropical humide. Autre exemple : les \textbf{anodes sacrificielles} en magnésium ou zinc protégeant les canalisations en acier (pipeline Tchad--Cameroun, installations de la SONARA).
\end{enumerate}
\textbf{Barème indicatif (sur 10 pts) :} interprétation exp. 1--2 + équation (1,5 pt) ; gaz + test (1 pt) ; classification justifiée (2 pts) ; quantités initiales (1 pt) ; tableau d'avancement + réactif limitant (1,5 pt) ; volume de \ce{H2} (1 pt) ; concentration finale (1 pt) ; application (1 pt).

\textbf{Points de vigilance :} attention au coefficient \textbf{2} devant \ce{H3O+} dans le tableau d'avancement ($0{,}100 - 2x$) ; reconnaître le cas particulier des proportions stœchiométriques (les deux réactifs disparaissent ensemble) ; ne pas oublier les unités (\si{L}, \si{mol.L^{-1}}) et la conversion $V = \SI{200}{mL} = \SI{0,200}{L}$.
\end{corrige}

\newpage

\coursec{Fiche bilan}

\coursec{Leçon 06 : Couple oxydant-réducteur et classification électrochimique}

\begin{popfigure}
\begin{tikzpicture}[
  mindmap, grow cyclic,
  every node/.style={concept, font=\small},
  concept color=popblue,
  root concept/.append style={concept color=popdark, font=\bfseries\small, text=white, minimum size=2.8cm},
  level 1/.append style={level distance=3.6cm, sibling angle=51.4, minimum size=2.2cm, font=\scriptsize},
  level 1 concept/.append style={text=white},
]
\node[root concept]{Couples oxydant/réducteur et classification électrochimique}
  child[concept color=popblue]{ node {Couple Ox/Red : définition, écriture $\mathrm{Ox} + n\,e^- \rightleftharpoons \mathrm{Red}$} }
  child[concept color=poporange]{ node {Couples $\ce{M^{n+}}/\ce{M}$ : métal + ion métallique} }
  child[concept color=popgreen]{ node {Couple $\ce{H3O+}/\ce{H2}$ : action des acides sur les métaux} }
  child[concept color=poppurple]{ node {Classification électrochimique qualitative} }
  child[concept color=poppink]{ node {Règle du gamma : prévoir les réactions} }
  child[concept color=popgold]{ node {Expériences clés : lame de Cu dans $\ce{Ag+}$, Zn dans HCl} }
  child[concept color=popmuted]{ node {Applications : corrosion, protection, métallurgie} };
\end{tikzpicture}
\legende{Carte mentale de la leçon : les sept idées forces à retenir.}
\end{popfigure}

\begin{savaistu}[Situation déclenchante : la corrosion et la récupération des métaux]
Au Cameroun, les infrastructures métalliques (ponts sur la Sanaga, toitures en tôle, armatures du barrage de Lom Pangar) subissent la \cle{corrosion} : le fer se transforme en rouille ($\ce{Fe2O3 * n H2O}$) au contact de l'air humide. Comment l'éviter ? On utilise le \cle{zinc} (galvanisation) ou la \cle{protection cathodique}. À l'inverse, en métallurgie (ex. traitement des minerais à Figuil ou récupération de l'argent en bijouterie), on cherche à \emph{faire réagir} un métal avec des ions métalliques. Tous ces phénomènes relèvent d'une même logique : les \cle{réactions d'oxydoréduction} entre couples \cle{oxydant/réducteur}. Cette leçon te donne les outils pour les comprendre et les prévoir.
\end{savaistu}

\courssub{1. Couple oxydant-réducteur : définition et écriture}

\begin{definition}[Couple oxydant/réducteur (couple redox)]
Un \cle{couple oxydant/réducteur} (noté \cle{Ox/Red}) est un ensemble de deux espèces chimiques conjuguées : un \cle{oxydant} (\cle{Ox}) et un \cle{réducteur} (\cle{Red}), qui peuvent se transformer l'un en l'autre par échange de $n$ électrons ($n \in \mathbb{N}^*$).
La transformation s'écrit sous forme de \cle{demi-équation électronique} (ou demi-réaction) :
\[ \mathrm{Ox} + n\,e^- \rightleftharpoons \mathrm{Red} \]
L'oxydant est l'espèce qui \emph{capte} les électrons (il s'\emph{réduit}) ; le réducteur est l'espèce qui \emph{cède} les électrons (il s'\emph{oxyde}).
\end{definition}

\begin{propriete}[Conventions d'écriture]
\begin{itemize}
  \item Le couple se note toujours \textbf{Ox/Red} (l'oxydant en premier, le réducteur en second), séparés par une barre oblique : $\ce{Cu^2+}/\ce{Cu}$, $\ce{Fe^3+}/\ce{Fe^2+}$, $\ce{H3O+}/\ce{H2}$.
  \item Pour un métal $\ce{M}$ et son cation $\ce{M^{n+}}$, le couple s'écrit $\ce{M^{n+}}/\ce{M}$.
  \item La demi-équation s'équilibre en atomes \emph{et} en charges.
\end{itemize}
\end{propriete}

\begin{methode}[Écrire et équilibrer une demi-équation $\mathrm{Ox} + n\,e^- \rightleftharpoons \mathrm{Red}$]
\begin{enumerate}
  \item Écrire les formules de l'oxydant et du réducteur (données par le couple).
  \item Équilibrer l'élément principal (autre que O et H).
  \item Équilibrer l'oxygène en ajoutant des molécules d'eau $\ce{H2O}$.
  \item Équilibrer l'hydrogène en ajoutant des ions $\ce{H3O+}$ (milieu aqueux acide).
  \item Équilibrer les charges en ajoutant des électrons $e^-$ \textbf{toujours du côté de l'oxydant} (celui qui les capte).
  \item Vérifier : même nombre d'atomes de chaque élément et même charge totale des deux côtés.
\end{enumerate}
\end{methode}

\begin{exemplebox}[Demi-équations classiques à connaître par cœur]
\begin{itemize}
  \item Couple $\ce{Cu^2+}/\ce{Cu}$ : $\ce{Cu^2+ + 2e- <=> Cu}$
  \item Couple $\ce{Ag+}/\ce{Ag}$ : $\ce{Ag+ + e- <=> Ag}$
  \item Couple $\ce{Zn^2+}/\ce{Zn}$ : $\ce{Zn^2+ + 2e- <=> Zn}$
  \item Couple $\ce{Fe^2+}/\ce{Fe}$ : $\ce{Fe^2+ + 2e- <=> Fe}$
  \item Couple $\ce{Fe^3+}/\ce{Fe^2+}$ : $\ce{Fe^3+ + e- <=> Fe^2+}$
  \item Couple $\ce{MnO4-}/\ce{Mn^2+}$ (milieu acide) : $\ce{MnO4- + 8H3O+ + 5e- <=> Mn^2+ + 12H2O}$
  \item Couple $\ce{Cr2O7^2-}/\ce{Cr^3+}$ (milieu acide) : $\ce{Cr2O7^2- + 14H3O+ + 6e- <=> 2Cr^3+ + 21H2O}$
  \item Couple $\ce{H3O+}/\ce{H2}$ : $\ce{2H3O+ + 2e- <=> H2 + 2H2O}$
\end{itemize}
\end{exemplebox}

\courssub{2. Réaction entre un métal et un ion métallique : couples $\ce{M^{n+}}/\ce{M}$}

\begin{experience}[Mise en évidence de la réaction $\ce{Cu} + \ce{Ag+}$ : « l'arbre d'argent »]
\begin{itemize}
  \item \textbf{Matériel :} Bécher, solution de nitrate d'argent $\ce{AgNO3}$ (\SI{0,1}{mol.L^{-1}}), lame de cuivre brillante, pince.
  \item \textbf{Protocole :} Plonger la lame de Cu dans la solution incolore d'$\ce{AgNO3}$. Observer pendant quelques minutes.
  \item \textbf{Observations :} La lame de cuivre se couvre d'un dépôt grisâtre brillant (cristaux d'argent métallique). La solution devient bleue (apparition des ions $\ce{Cu^2+}$).
  \item \textbf{Interprétation :} Le cuivre métal ($\ce{Cu}$, réducteur) cède des électrons aux ions argent ($\ce{Ag+}$, oxydant). Il y a réaction spontanée.
  \item \textbf{Couples en jeu :} $\ce{Ag+}/\ce{Ag}$ et $\ce{Cu^2+}/\ce{Cu}$.
  \item \textbf{Demi-équations :}
  \begin{align*}
    \text{Oxydation (Cu)} &: \ce{Cu -> Cu^2+ + 2e-} \\
    \text{Réduction (Ag+)} &: \ce{2Ag+ + 2e- -> 2Ag}
  \end{align*}
  \item \textbf{Équation bilan :} $\ce{Cu + 2Ag+ -> Cu^2+ + 2Ag}$
\end{itemize}
\end{experience}

\begin{experience}[Réaction $\ce{Zn} + \ce{Cu^2+}$]
\begin{itemize}
  \item \textbf{Observation :} Un morceau de zinc plongé dans une solution de sulfate de cuivre ($\ce{CuSO4}$, bleue) se couvre de cuivre rouge ; la solution se décolore.
  \item \textbf{Équation bilan :} $\ce{Zn + Cu^2+ -> Zn^2+ + Cu}$
  \item \textbf{Conclusion :} Le zinc est un réducteur plus fort que le cuivre (il « chasse » le cuivre de sa solution).
\end{itemize}
\end{experience}

\begin{propriete}[Règle générale pour les couples $\ce{M^{n+}}/\ce{M}$]
Lorsqu'on met en présence un métal $\ce{M_1}$ (réducteur) et un ion métallique $\ce{M_2^{n+}}$ (oxydant), une réaction spontanée a lieu \textbf{si le métal $\ce{M_1}$ est un meilleur réducteur que le métal $\ce{M_2}$} (c'est-à-dire si le couple $\ce{M_2^{n+}}/\ce{M_2}$ est un meilleur oxydant que $\ce{M_1^{n+}}/\ce{M_1}$). On verra comment classer cela précisément à la section 4.
\end{propriete}

\courssub{3. Action des ions $\ce{H3O+}$ sur les métaux : le couple $\ce{H3O+}/\ce{H2}$}

\begin{experience}[Action d'un acide sur différents métaux]
\begin{itemize}
  \item \textbf{Matériel :} Tubes à essai, acide chlorhydrique dilué ($\ce{HCl}$), \SI{1}{mol.L^{-1}}, morceaux de zinc ($\ce{Zn}$), fer ($\ce{Fe}$), cuivre ($\ce{Cu}$), argent ($\ce{Ag}$), allumette.
  \item \textbf{Protocole :} Introduire chaque métal dans un tube contenant \SI{5}{mL} d'acide. Observer. Approcher une flamme de l'embouchure si dégagement gazeux.
  \item \textbf{Observations :}
  \begin{center}
  \begin{tabularx}{\linewidth}{|l|X|X|}\hline
  \rowcolor{popblueL} \textbf{Métal} & \textbf{Observation} & \textbf{Test à la flamme} \\\hline
  $\ce{Zn}$, $\ce{Fe}$, $\ce{Al}$, $\ce{Mg}$ & Dégagement gazeux rapide (bulles), solution qui s'échauffe & \textbf{Détonation caractéristique (« aboiement »)} $\to$ \cle{Dihydrogène $\ce{H2}$} \\\hline
  $\ce{Cu}$, $\ce{Ag}$, $\ce{Au}$, $\ce{Pt}$ & Aucune réaction visible, pas de bulles, pas de chauffage & --- \\\hline
  \end{tabularx}
  \end{center}
  \item \textbf{Interprétation :} Les ions $\ce{H3O+}$ (oxydant) sont réduits en $\ce{H2}$ (réducteur) par les métaux $\ce{Zn}$, $\ce{Fe}$, etc. (réducteurs). Le couple en jeu est $\ce{H3O+}/\ce{H2}$.
  \item \textbf{Demi-équation :} $\ce{2H3O+ + 2e- <=> H2 + 2H2O}$
  \item \textbf{Équations bilan (exemples) :}
  \begin{align*}
    \ce{Zn + 2H3O+ &-> Zn^2+ + H2 ^ + 2H2O} \\
    \ce{Fe + 2H3O+ &-> Fe^2+ + H2 ^ + 2H2O} \\
    \ce{2Al + 6H3O+ &-> 2Al^3+ + 3H2 ^ + 6H2O}
  \end{align*}
\end{itemize}
\end{experience}

\begin{aretenir}[Identification du dihydrogène $\ce{H2}$]
Le gaz $\ce{H2}$ est incolore, inodore, peu soluble dans l'eau. Il se reconnaît par le \textbf{test de l'aboiement} : une flamme approchée de l'embouchure du tube provoque une \textbf{petite explosion sèche} (détonation) avec formation d'eau $\ce{2H2 + O2 -> 2H2O}$. \textbf{Attention :} ne jamais faire ce test sur un gros volume de gaz (risque d'explosion violente).
\end{aretenir}

\begin{attention}[Piège fréquent : « L'acide attaque le métal »]
On dit souvent « l'acide attaque le métal ». Rigoureusement, ce sont les \textbf{ions $\ce{H30+}$} (l'oxydant du couple $\ce{H3O+}/\ce{H2}$) qui oxydent le métal. L'anion de l'acide ($\ce{Cl-}$, $\ce{SO4^2-}$, $\ce{NO3-}$) est \emph{spectateur} (sauf cas particuliers comme $\ce{HNO3}$ concentré ou $\ce{H2SO4}$ chaud où l'anion devient oxydant). En Première, on travaille avec des acides dilués : seul $\ce{H3O+}$ oxyde.
\end{attention}

\courssub{4. Construction de la classification électrochimique qualitative}

\begin{definition}[Classification électrochimique qualitative]
C'est un \textbf{tableau classant les couples oxydant/réducteur par pouvoir oxydant décroissant de l'oxydant} (donc par pouvoir réducteur croissant du réducteur), établi à partir de résultats expérimentaux (réactions métal/ion métallique, métal/acide).
\end{definition}

\begin{methode}[Construire la classification à partir d'expériences]
\begin{enumerate}
  \item Lister les couples étudiés : $\ce{Ag+}/\ce{Ag}$, $\ce{Cu^2+}/\ce{Cu}$, $\ce{H3O+}/\ce{H2}$, $\ce{Zn^2+}/\ce{Zn}$, $\ce{Fe^2+}/\ce{Fe}$, etc.
  \item Pour chaque paire de couples, noter si réaction il y a (expériences sections 2 et 3).
  \item \textbf{Règle :} Si $\ce{Ox_1}$ réagit avec $\ce{Red_2}$, alors $\ce{Ox_1}$ est un \emph{meilleur oxydant} que $\ce{Ox_2}$ (donc $\ce{Ox_1}/\ce{Red_1}$ se place \textbf{au-dessus} de $\ce{Ox_2}/\ce{Red_2}$).
  \item Ranger les couples du plus fort oxydant (haut) au plus faible oxydant (bas).
\end{enumerate}
\end{methode}

\begin{exemplebox}[Construction pas à pas (résumé des résultats expérimentaux)]
\begin{enumerate}
  \item $\ce{Cu}$ réduit $\ce{Ag+}$ $\to$ $\ce{Ag+}/\ce{Ag}$ au-dessus de $\ce{Cu^2+}/\ce{Cu}$.
  \item $\ce{Zn}$ réduit $\ce{Cu^2+}$ $\to$ $\ce{Cu^2+}/\ce{Cu}$ au-dessus de $\ce{Zn^2+}/\ce{Zn}$.
  \item $\ce{Zn}$, $\ce{Fe}$ réduisent $\ce{H3O+}$ (dégagent $\ce{H2}$) $\to$ $\ce{H3O+}/\ce{H2}$ au-dessus de $\ce{Zn^2+}/\ce{Zn}$ et $\ce{Fe^2+}/\ce{Fe}$.
  \item $\ce{Cu}$, $\ce{Ag}$ ne réduisent \emph{pas} $\ce{H3O+}$ $\to$ $\ce{Cu^2+}/\ce{Cu}$ et $\ce{Ag+}/\ce{Ag}$ au-dessus de $\ce{H3O+}/\ce{H2}$.
\end{enumerate}
\emph{Ordre final (haut $\to$ bas) :} $\ce{Ag+}/\ce{Ag}$ > $\ce{Cu^2+}/\ce{Cu}$ > $\ce{H3O+}/\ce{H2}$ > $\ce{Fe^2+}/\ce{Fe}$ > $\ce{Zn^2+}/\ce{Zn}$ > $\ce{Al^3+}/\ce{Al}$ > $\ce{Mg^2+}/\ce{Mg}$.
\end{exemplebox}

\begin{popfigure}
\begin{tikzpicture}[
  scale=0.9,
  every node/.style={font=\small},
  couple/.style={rectangle, draw=popdark, thick, fill=popcream, rounded corners, minimum width=3.5cm, minimum height=0.8cm, align=center},
  arrow/.style={-Stealth, thick, popdark},
  label/.style={font=\scriptsize, text=popmuted}
]
% Couples (Top to Bottom)
\node[couple, fill=poppink] (Ag) at (0, 6) {$\ce{Ag+}/\ce{Ag}$ \\ \textbf{Oxydant le plus fort}};
\node[couple, fill=poporangeL] (Cu) at (0, 4.5) {$\ce{Cu^2+}/\ce{Cu}$};
\node[couple, fill=popgoldL] (H) at (0, 3) {$\ce{H3O+}/\ce{H2}$ \\ \textbf{Repère central}};
\node[couple, fill=popgreenL] (Fe) at (0, 1.5) {$\ce{Fe^2+}/\ce{Fe}$};
\node[couple, fill=popblueL] (Zn) at (0, 0) {$\ce{Zn^2+}/\ce{Zn}$};
\node[couple, fill=poppurpleL] (Al) at (0, -1.5) {$\ce{Al^3+}/\ce{Al}$};
\node[couple, fill=poppinkL] (Mg) at (0, -3) {$\ce{Mg^2+}/\ce{Mg}$ \\ \textbf{Réducteur le plus fort}};

% Arrows and labels
\draw[arrow] (Ag.south) -- (Cu.north) node[midway, right, label] {$\ce{Ag+}$ oxyde $\ce{Cu}$};
\draw[arrow] (Cu.south) -- (H.north) node[midway, right, label] {$\ce{Cu^2+}$ oxyde $\ce{H2}$};
\draw[arrow] (H.south) -- (Fe.north) node[midway, right, label] {$\ce{H3O+}$ oxyde $\ce{Fe}$, $\ce{Zn}$...};
\draw[arrow] (Fe.south) -- (Zn.north);
\draw[arrow] (Zn.south) -- (Al.north);
\draw[arrow] (Al.south) -- (Mg.north);

% Side annotations
\node[left=0.5cm of Ag, label, rotate=90, anchor=south] {Pouvoir oxydant $\uparrow$};
\node[left=0.5cm of Mg, label, rotate=90, anchor=north] {Pouvoir réducteur $\uparrow$};
\node[right=0.5cm of Ag, label, rotate=90, anchor=south] {Métaux \textbf{nobles}};
\node[right=0.5cm of Cu, label, rotate=90, anchor=center] {(Cu, Ag, Au, Pt)};
\node[right=0.5cm of H, label, rotate=90, anchor=center] {Non attaqués par $\ce{H+}$};
\node[right=0.5cm of Zn, label, rotate=90, anchor=north] {Métaux \textbf{non nobles}};
\node[right=0.5cm of Mg, label, rotate=90, anchor=north] {(Zn, Fe, Al, Mg...)};
\node[right=0.5cm of Mg, label, rotate=90, anchor=north] {Attaqués par $\ce{H+}$ $\to$ $\ce{H2}$};
\end{tikzpicture}
\legende{Classification électrochimique qualitative simplifiée (extrait). Les couples réels en contiennent beaucoup plus (ex. $\ce{K+}/\ce{K}$, $\ce{Ca^2+}/\ce{Ca}$, $\ce{Cl2}/\ce{Cl-}$, $\ce{MnO4-}/\ce{Mn^2+}$).}
\end{popfigure}

\begin{propriete}[Position du couple $\ce{H3O+}/\ce{H2}$ et distinction nobles / non nobles]
Le couple $\ce{H3O+}/\ce{H2}$ occupe une \textbf{position charnière} (repère central).
\begin{itemize}
  \item \textbf{Au-dessus} (oxydants plus forts que $\ce{H3O+}$) : $\ce{Ag+}$, $\ce{Cu^2+}$, $\ce{Au^3+}$, $\ce{Pt^2+}$... Les métaux correspondants (\textbf{métaux nobles}) ne sont \emph{pas} attaqués par les acides dilués (pas de $\ce{H2}$).
  \item \textbf{En-dessous} (oxydants plus faibles que $\ce{H3O+}$) : $\ce{Fe^2+}$, $\ce{Zn^2+}$, $\ce{Al^3+}$, $\ce{Mg^2+}$, $\ce{Ca^2+}$... Les métaux correspondants (\textbf{métaux non nobles}) \emph{sont} attaqués par les acides dilués avec dégagement de $\ce{H2}$.
\end{itemize}
\end{propriete}

\courssub{5. Utilisation de la classification : prévision des réactions (Règle du $\gamma$)}

\begin{definition}[Règle du gamma ($\gamma$) — Prévision d'une réaction d'oxydoréduction]
Soit deux couples $\ce{Ox_1}/\ce{Red_1}$ et $\ce{Ox_2}/\ce{Red_2}$ placés dans la classification avec $\ce{Ox_1}/\ce{Red_1}$ \textbf{au-dessus} de $\ce{Ox_2}/\ce{Red_2}$.
\begin{itemize}
  \item L'oxydant du couple du haut ($\ce{Ox_1}$) peut oxyder le réducteur du couple du bas ($\ce{Red_2}$).
  \item La réaction spontanée est : $\ce{Ox_1 + Red_2 -> Red_1 + Ox_2}$.
  \item Graphiquement, on trace un \textbf{$\gamma$ (gamma)} partant de $\ce{Ox_1}$ (haut gauche) vers $\ce{Red_2}$ (bas droite). Les produits ($\ce{Red_1}$, $\ce{Ox_2}$) se trouvent aux autres extrémités du $\gamma$ et \textbf{ne réagissent pas entre eux}.
\end{itemize}
\end{definition}

\begin{attention}[Erreur fatale : sens du gamma]
\textbf{Le gamma part toujours de l'oxydant du couple DU HAUT vers le réducteur du couple DU BAS.}
\begin{itemize}
  \item \textbf{FAUX :} « Le réducteur du couple du haut réagit avec l'oxydant du couple du bas. » (Ce serait l'inverse, réaction non spontanée).
  \item \textbf{FAUX :} « On trace le gamma du bas vers le haut. »
  \item \textbf{Astuce :} L'oxydant « tombe » sur le réducteur (haut $\to$ bas).
\end{itemize}
\end{attention}

\begin{methode}[Prévoir et écrire l'équation bilan d'une réaction redox]
\begin{enumerate}
  \item Identifier les deux couples mis en présence (ex: $\ce{Fe^2+}/\ce{Fe}$ et $\ce{Cu^2+}/\ce{Cu}$).
  \item Les placer dans la classification électrochimique.
  \item Appliquer la règle du $\gamma$ : l'oxydant du couple le plus haut réagit avec le réducteur du couple le plus bas.
  \item Écrire les deux demi-équations (réduction pour l'oxydant du haut, oxydation pour le réducteur du bas).
  \item Multiplier les demi-équations par les coefficients nécessaires pour avoir \textbf{le même nombre d'électrons} des deux côtés.
  \item Additionner les deux demi-équations et simplifier les électrons (et les espèces communes).
  \item Vérifier la conservation des éléments et des charges.
\end{enumerate}
\end{methode}

\begin{exoresolu}[Prévoir la réaction entre $\ce{Fe}$ et $\ce{Cu^2+}$]
\textbf{Énoncé :} On plonge un clou de fer dans une solution de sulfate de cuivre ($\ce{CuSO4}$). Prédire ce qui se passe et écrire l'équation bilan.
\tcblower
\textbf{Solution :}
\begin{enumerate}
  \item Couples : $\ce{Cu^2+}/\ce{Cu}$ et $\ce{Fe^2+}/\ce{Fe}$.
  \item Classification : $\ce{Cu^2+}/\ce{Cu}$ est \textbf{au-dessus} de $\ce{Fe^2+}/\ce{Fe}$.
  \item Gamma : $\ce{Cu^2+}$ (oxydant du haut) + $\ce{Fe}$ (réducteur du bas) $\to$ $\ce{Cu}$ + $\ce{Fe^2+}$.
  \item Demi-équations :
  \begin{align*}
    \text{Réduction :} &\quad \ce{Cu^2+ + 2e- -> Cu} \\
    \text{Oxydation :} &\quad \ce{Fe -> Fe^2+ + 2e-}
  \end{align*}
  \item Électrons déjà équilibrés (2 et 2).
  \item Bilan : $\ce{Cu^2+ + Fe -> Cu + Fe^2+}$.
  \item Vérif : Charges : $2+ + 0 = 0 + 2+$ (OK). Atomes : 1 Cu, 1 Fe (OK).
  \item \textbf{Observation attendue :} Le clou se couvre de cuivre rouge, la solution bleue se décolore (devient verte $\ce{Fe^2+}$).
\end{enumerate}
\end{exoresolu}

\begin{exoresolu}[Réaction entre $\ce{Ag}$ et $\ce{H3O+}$]
\textbf{Énoncé :} Une pièce en argent est plongée dans de l'acide chlorhydrique dilué. Y a-t-il réaction ?
\tcblower
\textbf{Solution :}
\begin{enumerate}
  \item Couples : $\ce{Ag+}/\ce{Ag}$ et $\ce{H3O+}/\ce{H2}$.
  \item Classification : $\ce{Ag+}/\ce{Ag}$ est \textbf{au-dessus} de $\ce{H3O+}/\ce{H2}$.
  \item Gamma : L'oxydant du haut est $\ce{Ag+}$ ; le réducteur du bas est $\ce{H2}$.
  \item Les réactifs mis en présence sont $\ce{Ag}$ (réducteur du haut) et $\ce{H3O+}$ (oxydant du bas). Ce n'est \textbf{PAS} le sens du gamma.
  \item \textbf{Conclusion : Pas de réaction spontanée.} L'argent ne « dissout » pas dans l'acide chlorhydrique dilué. (C'est pour cela qu'on l'utilise en bijouterie et qu'il ne rouille pas).
\end{enumerate}
\end{exoresolu}

\courssub{6. Synthèse et applications au quotidien camerounais}

\begin{savaistu}[Applications concrètes au Cameroun]
\begin{itemize}
  \item \textbf{Corrosion du fer (rouille) :} $\ce{Fe}$ (réducteur) + $\ce{O2}$ (oxydant, couple $\ce{O2}/\ce{H2O}$ très haut dans la classification) + $\ce{H2O}$ $\to$ $\ce{Fe^2+}$ puis $\ce{Fe2O3} \cdot n\,\ce{H2O}$. C'est une réaction redox spontanée (gamma immense).
  \item \textbf{Galvanisation (toitures, pylônes ENEO, barrages) :} On recouvre le fer d'une couche de \textbf{zinc} ($\ce{Zn}$). $\ce{Zn}$ est \emph{plus bas} que $\ce{Fe}$ dans la classification ($\ce{Zn}$ meilleur réducteur). Si la couche craque, $\ce{Zn}$ s'oxyde \emph{à la place} du fer (anode sacrificielle) : $\ce{Zn -> Zn^2+ + 2e-}$ protège $\ce{Fe}$.
  \item \textbf{Récupération de l'argent (bijouterie, déchets électroniques) :} On utilise du cuivre (fil de Cu) pour précipiter l'argent depuis une solution d'$\ce{Ag+}$ : $\ce{Cu + 2Ag+ -> Cu^2+ + 2Ag}$ (arbre d'argent).
  \item \textbf{Piles et batteries :} Toute pile (pile saline, alcaline, lithium-ion, batterie de voiture au plomb) fonctionne sur le principe de deux couples redox séparés (anode et cathode) reliés par un fil (électrons) et un pont salin (ions). La tension ($E$) dépend de l'écart dans la classification.
  \item \textbf{Métallurgie (Figuil, fer) :} Réduction du minerai de fer ($\ce{Fe2O3}$) par le coke ($\ce{C}$ / $\ce{CO}$) à haute température : $\ce{Fe2O3 + 3CO -> 2Fe + 3CO2}$. Le carbone/$\ce{CO}$ est un réducteur très fort (très bas dans la classification à haute T).
  \item \textbf{Vinaigre et bissap :} L'acidité (ions $\ce{H3O+}$) du vinaigre ou du jus de bissap (acide citrique, acide malique) peut attaquer les casseroles en aluminium ($\ce{Al}$ est non noble, au-dessus de $\ce{H2}$) $\to$ $\ce{2Al + 6H3O+ -> 2Al^3+ + 3H2 + 6H2O}$. Préférer l'inox (fer + chrome, passivé) ou l'émail.
\end{itemize}
\end{savaistu}

\begin{aretenir}[L'essentiel de la leçon — Synthèse structurée]
\textbf{1. Couple Ox/Red}
\begin{itemize}
  \item Deux espèces conjuguées $\ce{Ox}$ (oxydant) et $\ce{Red}$ (réducteur) liées par $\ce{Ox + n e- <=> Red}$.
  \item Écriture : $\ce{Ox/Red}$ (ex: $\ce{Cu^2+}/\ce{Cu}$, $\ce{H3O+}/\ce{H2}$).
  \item Demi-équation : équilibrer atomes (O par $\ce{H2O}$, H par $\ce{H3O+}$) puis charges par $e^-$ (côté Ox).
\end{itemize}

\textbf{2. Couples métalliques $\ce{M^{n+}}/\ce{M}$}
\begin{itemize}
  \item Expérience clé : $\ce{Cu + 2Ag+ -> Cu^2+ + 2Ag}$ (arbre d'argent), $\ce{Zn + Cu^2+ -> Zn^2+ + Cu}$.
  \item Le métal solide est le réducteur ; l'ion est l'oxydant.
\end{itemize}

\textbf{3. Couple $\ce{H3O+}/\ce{H2}$}
\begin{itemize}
  \item Demi-équation : $\ce{2H3O+ + 2e- <=> H2 + 2H2O}$.
  \item Métaux au-dessus de $\ce{H2}$ (Zn, Fe, Al, Mg...) : réagissent avec acides dilués $\to$ $\ce{H2}$ (aboiement).
  \item Métaux en-dessous (Cu, Ag, Au, Pt) : pas de réaction avec acides dilués.
\end{itemize}

\textbf{4. Classification électrochimique}
\begin{itemize}
  \item Classement par pouvoir oxydant \textbf{décroissant} (haut $\to$ bas).
  \item $\ce{H3O+}/\ce{H2}$ = repère central (sépare nobles / non nobles).
  \item Construction par expériences croisées (qui oxyde qui ?).
\end{itemize}

\textbf{5. Règle du $\gamma$ (Prévision)}
\begin{itemize}
  \item $\gamma$ part de \textbf{$\ce{Ox}$ du couple DU HAUT} vers \textbf{$\ce{Red}$ du couple DU BAS}.
  \item Réaction : $\ce{Ox_{haut} + Red_{bas} -> Red_{haut} + Ox_{bas}}$.
  \item Produits stables (ne réagissent pas entre eux).
\end{itemize}

\textbf{6. Méthode « Équation bilan »}
Couples $\to$ Classification $\to$ Gamma $\to$ Demi-équations $\to$ Égaliser $e^-$ $\to$ Additionner $\to$ Vérifier.
\end{aretenir}

\begin{aretenir}[Checklist avant le Probatoire]
\begin{itemize}
  \item[$\square$] Je sais définir un couple oxydant/réducteur et l'écrire sous la forme $\ce{Ox/Red}$.
  \item[$\square$] Je sais écrire et équilibrer la demi-équation $\ce{Ox + n e- <=> Red}$ d'un couple donné (milieu acide : $\ce{H3O+}$, $\ce{H2O}$).
  \item[$\square$] Je connais par cœur les couples $\ce{Cu^2+}/\ce{Cu}$, $\ce{Ag+}/\ce{Ag}$, $\ce{Zn^2+}/\ce{Zn}$, $\ce{Fe^2+}/\ce{Fe}$, $\ce{Fe^3+}/\ce{Fe^2+}$, $\ce{H3O+}/\ce{H2}$.
  \item[$\square$] Je sais décrire l'expérience de la lame de cuivre dans une solution de $\ce{Ag+}$ (dépôt d'Ag, bleuissement) et écrire l'équation bilan.
  \item[$\square$] Je sais décrire l'action de l'acide chlorhydrique dilué sur Zn, Fe, Cu, Ag et identifier $\ce{H2}$ par l'aboiement.
  \item[$\square$] Je sais que seuls les métaux placés \textbf{au-dessus} de $\ce{H2}$ dans la classification sont attaqués par $\ce{H3O+}$.
  \item[$\square$] Je sais construire une classification électrochimique qualitative à partir de résultats expérimentaux (tableau de réactivités).
  \item[$\square$] Je sais placer le couple $\ce{H3O+}/\ce{H2}$ entre les métaux nobles (en dessous) et non nobles (au-dessus).
  \item[$\square$] Je sais appliquer la \textbf{règle du $\gamma$} (haut $\to$ bas : $\ce{Ox_{haut}}$ + $\ce{Red_{bas}}$) pour prévoir une réaction.
  \item[$\square$] Je sais combiner deux demi-équations (égalité des électrons) pour obtenir l'équation bilan.
  \item[$\square$] Je sais citer des applications : corrosion/protection (galvanisation, anode sacrificielle), récupération des métaux, piles, métallurgie.
  \item[$\square$] Je vérifie \textbf{systématiquement} la conservation des éléments et des charges dans mes équations finales.
\end{itemize}
\end{aretenir}

\begin{exercice}[Entraînement — Prévoir et écrire les bilans]
\begin{enumerate}
  \item On plonge une lame de zinc dans une solution de nitrate de fer(II) $\ce{Fe(NO3)2}$. Prédire ce qui se passe. Écrire les demi-équations et l'équation bilan.
  \item On ajoute de la poudre de fer dans une solution de nitrate d'argent $\ce{AgNO3}$. Prédire. Écrire le bilan.
  \item On met du cuivre en contact avec une solution de $\ce{Zn^2+}$. Réaction ? Justifier par la classification.
  \item Écrire la demi-équation du couple $\ce{MnO4-}/\ce{Mn^2+}$ en milieu acide. Équilibrer.
  \item \textbf{Défi :} On a trois couples : $\ce{X^2+}/\ce{X}$, $\ce{Y^2+}/\ce{Y}$, $\ce{Z^2+}/\ce{Z}$. Expériences : $\ce{X}$ réduit $\ce{Y^2+}$ ; $\ce{Z}$ réduit $\ce{X^2+}$ ; $\ce{Y}$ ne réduit pas $\ce{Z^2+}$. Classer les trois couples du plus fort oxydant au plus faible. Où placer $\ce{H3O+}/\ce{H2}$ si $\ce{X}$ réagit avec l'acide mais $\ce{Y}$ ne réagit pas ?
\end{enumerate}
\end{exercice}

\begin{corrige}[Exercice 1]
\begin{enumerate}
  \item \textbf{Cuivre :} Couples $\ce{Fe^2+}/\ce{Fe}$ et $\ce{Zn^2+}/\ce{Zn}$. Classification : $\ce{Fe^2+}/\ce{Fe}$ au-dessus de $\ce{Zn^2+}/\ce{Zn}$. Gamma : $\ce{Fe^2+}$ (Ox haut) + $\ce{Zn}$ (Red bas) $\to$ $\ce{Fe}$ + $\ce{Zn^2+}$.
  \begin{align*}
    \text{Red :} &\quad \ce{Fe^2+ + 2e- -> Fe} \\
    \text{Ox :} &\quad \ce{Zn -> Zn^2+ + 2e-} \\
    \text{Bilan :} &\quad \ce{Zn + Fe^2+ -> Zn^2+ + Fe}
  \end{align*}
  Observation : Dépôt de fer gris sur le zinc, solution qui se décolore (si $\ce{Fe^2+}$ coloré) ou incolore.
  \item \textbf{Argent :} Couples $\ce{Ag+}/\ce{Ag}$ et $\ce{Fe^2+}/\ce{Fe}$. $\ce{Ag+}/\ce{Ag}$ bien au-dessus. Gamma : $\ce{Ag+}$ + $\ce{Fe}$ $\to$ $\ce{Ag}$ + $\ce{Fe^2+}$.
  \begin{align*}
    \text{Red :} &\quad \ce{Ag+ + e- -> Ag} \quad (\times 2) \\
    \text{Ox :} &\quad \ce{Fe -> Fe^2+ + 2e-} \\
    \text{Bilan :} &\quad \ce{Fe + 2Ag+ -> Fe^2+ + 2Ag}
  \end{align*}
  Observation : Dépôt d'argent (gris/noir) sur le fer, solution qui verdit ($\ce{Fe^2+}$).
  \item \textbf{Cuivre / Zinc :} Couples $\ce{Cu^2+}/\ce{Cu}$ et $\ce{Zn^2+}/\ce{Zn}$. $\ce{Cu^2+}/\ce{Cu}$ au-dessus. Réactifs : $\ce{Cu}$ (Red haut) et $\ce{Zn^2+}$ (Ox bas). \textbf{Pas le sens du gamma.} $\to$ \textbf{Pas de réaction.}
  \item \textbf{$\ce{MnO4-}/\ce{Mn^2+}$ :} $\ce{MnO4- + 8H3O+ + 5e- <=> Mn^2+ + 12H2O}$.
  \item \textbf{Défi :}
  \begin{itemize}
    \item $\ce{X}$ réduit $\ce{Y^2+}$ $\to$ $\ce{Y^2+}/\ce{Y}$ au-dessus de $\ce{X^2+}/\ce{X}$.
    \item $\ce{Z}$ réduit $\ce{X^2+}$ $\to$ $\ce{X^2+}/\ce{X}$ au-dessus de $\ce{Z^2+}/\ce{Z}$.
    \item $\ce{Y}$ ne réduit pas $\ce{Z^2+}$ $\to$ cohérent (si $\ce{Y}$ au-dessus de $\ce{Z}$, $\ce{Y}$ est moins bon réducteur que $\ce{Z}$).
    \item Ordre : $\ce{Y^2+}/\ce{Y}$ > $\ce{X^2+}/\ce{X}$ > $\ce{Z^2+}/\ce{Z}$.
    \item $\ce{X}$ réagit avec acide $\to$ $\ce{X}$ au-dessus de $\ce{H2}$ $\to$ $\ce{X^2+}/\ce{X}$ au-dessus de $\ce{H3O+}/\ce{H2}$.
    \item $\ce{Y}$ ne réagit pas $\to$ $\ce{Y}$ en-dessous de $\ce{H2}$ $\to$ $\ce{Y^2+}/\ce{Y}$ en-dessous de $\ce{H3O+}/\ce{H2}$.
    \item \textbf{Contradiction ?} Non. $\ce{X}$ au-dessus de $\ce{H2}$, $\ce{Y}$ en-dessous. Mais on a dit $\ce{Y^2+}/\ce{Y}$ au-dessus de $\ce{X^2+}/\ce{X}$. Donc $\ce{Y^2+}/\ce{Y}$ > $\ce{X^2+}/\ce{X}$ > $\ce{H3O+}/\ce{H2}$ > $\ce{Z^2+}/\ce{Z}$.
    \item Attente : $\ce{Y}$ est un métal noble (ex: Cu), $\ce{X}$ est un métal non noble (ex: Fe), $\ce{Z}$ est un métal très non noble (ex: Zn).
  \end{itemize}
\end{enumerate}
\end{corrige}

\findecours

\end{document}
